\documentclass[11pt]{article}
\usepackage[T1]{fontenc}
\usepackage[utf8]{inputenc}
\usepackage[letterpaper,margin=1in,headheight=15pt]{geometry}
\usepackage{amsmath,amsthm,mathtools}
\usepackage{newtxtext,newtxmath}
\usepackage{booktabs,array,longtable}
\usepackage{microtype}
\usepackage[dvipsnames]{xcolor}
\usepackage{xurl,hyperref}
\usepackage{aliascnt}
\usepackage[nameinlink,noabbrev]{cleveref}
\usepackage{enumitem}
\usepackage{listings}
\usepackage{fancyhdr}
\hypersetup{colorlinks=true,linkcolor=MidnightBlue,citecolor=MidnightBlue,
 urlcolor=MidnightBlue,pdftitle={Arithmetic Local--Global Dichotomies for ProveIt's Three-Variable Keller Map},
 pdfauthor={Research draft prepared with ChatGPT},pdfsubject={Integral local-global failure and exact p-adic fiber distributions}}
\pagestyle{fancy}\fancyhf{}
\fancyhead[L]{\small Arithmetic local--global dichotomies}
\fancyhead[R]{\small ProveIt research}
\fancyfoot[C]{\thepage}
\setlength{\parskip}{4pt}
\setlength{\parindent}{1.15em}
\setlength{\emergencystretch}{2em}
\numberwithin{equation}{section}
\newtheorem{theorem}{Theorem}[section]
\theoremstyle{plain}
\newaliascnt{proposition}{theorem}
\newtheorem{proposition}[proposition]{Proposition}
\aliascntresetthe{proposition}
\crefname{proposition}{proposition}{propositions}
\theoremstyle{plain}
\newaliascnt{lemma}{theorem}
\newtheorem{lemma}[lemma]{Lemma}
\aliascntresetthe{lemma}
\crefname{lemma}{lemma}{lemmas}
\theoremstyle{plain}
\newaliascnt{corollary}{theorem}
\newtheorem{corollary}[corollary]{Corollary}
\aliascntresetthe{corollary}
\crefname{corollary}{corollary}{corollaries}
\theoremstyle{definition}
\newaliascnt{definition}{theorem}
\newtheorem{definition}[definition]{Definition}
\aliascntresetthe{definition}
\crefname{definition}{definition}{definitions}
\theoremstyle{definition}
\newaliascnt{question}{theorem}
\newtheorem{question}[question]{Research question}
\aliascntresetthe{question}
\crefname{question}{research question}{research questions}
\theoremstyle{definition}
\newaliascnt{example}{theorem}
\newtheorem{example}[example]{Example}
\aliascntresetthe{example}
\crefname{example}{example}{examples}
\theoremstyle{remark}
\newaliascnt{remark}{theorem}
\newtheorem{remark}[remark]{Remark}
\aliascntresetthe{remark}
\crefname{remark}{remark}{remarks}
\newcommand{\Z}{\mathbb Z}\newcommand{\Q}{\mathbb Q}
\newcommand{\R}{\mathbb R}\newcommand{\C}{\mathbb C}
\newcommand{\F}{\mathbb F}\newcommand{\A}{\mathbb A}\newcommand{\PP}{\mathbb P}
\newcommand{\Zp}{\Z_p}\newcommand{\Ztwo}{\Z_2}
\newcommand{\Qp}{\Q_p}\newcommand{\Oh}{\mathcal O}
\newcommand{\Spec}{\operatorname{Spec}}\newcommand{\Disc}{\operatorname{Disc}}
\newcommand{\Gal}{\operatorname{Gal}}\newcommand{\im}{\operatorname{im}}
\newcommand{\diag}{\operatorname{diag}}\newcommand{\lcm}{\operatorname{lcm}}
\newcommand{\vol}{\operatorname{vol}}
\newcommand{\target}{\boldsymbol b}
\newcommand{\psha}{\nolinkurl{e21766d04c2b8a9b2cdba0cd43e563024b3bd1b9}}
\newcommand{\rsha}{\nolinkurl{caf5685db1099cb90e608556d4b8d825734f29cb}}
\newcommand{\wtag}{\textup{[write]}}% text added when the report was written into ProveIt
\newcommand{\rpath}[1]{\nolinkurl{#1}}
% Part II (batch 58, manuscript 01)
\usepackage{needspace}
\newcommand{\dhsha}{\nolinkurl{1085b506d65e207a05b7e9c861bb1fe88432fe38}}
\newcommand{\Hsplit}{\mathcal H^{\mathrm{sp}}}
\newcommand{\Isplit}{\mathcal I^{\mathrm{sp}}}
\newcommand{\calB}{\mathcal B}
% Part III (batch 64, manuscript 02)
\newcommand{\nssha}{\nolinkurl{6d04e1e385f2fe7d1cfbdbd8fd8d4e45bd6b6c72}}
\newcommand{\NS}{\mathrm{ns}}\newcommand{\spc}{\mathrm{sp}}
\lstset{basicstyle=\ttfamily\footnotesize,columns=fullflexible,
 keepspaces=true,showstringspaces=false,breaklines=true,
 frame=single,rulecolor=\color{black!20},xleftmargin=0.4em,xrightmargin=0.4em,
 aboveskip=0.8em,belowskip=0.8em}
\setlist[enumerate]{itemsep=4pt,topsep=5pt}
\begin{document}
\hypersetup{pageanchor=false}
\begin{titlepage}
\setlength{\parindent}{0pt}
\vspace*{0.1in}
\noindent{\large\scshape ProveIt research article}\par
\vspace{0.25in}
\noindent{\huge\bfseries Arithmetic Local--Global\par Dichotomies for ProveIt's\par Three-Variable Keller Map}\par
\vspace{0.25in}
\noindent{\Large Split integral Hasse failures, a dense parameter family,\par and exact $p$-adic fiber laws}\par
\vspace{0.3in}
\noindent Research draft prepared for Vladimir Reshetnikov\par
\noindent September 24, 2026\par
\vspace{0.3in}
\begin{abstract}
We study the arithmetic image of the explicit polynomial map with constant
Jacobian determinant $-2$ developed in ProveIt's Jacobian-conjecture project.
The complex geometry and finite-field fiber laws are prior results
(\wtag{} of Gao and of the \texttt{royvanrijn/jacobian-research} notes
respectively; neither is part of ProveIt); our focus
is the distinction between rational, integral, and locally integral inverse
images. We construct an explicit family of targets having exactly three
rational preimages, none integral, but at least one integral preimage over
$\Z_p$ for every prime $p$. More generally, a two-parameter split-root family
admits an exact greatest-common-divisor criterion for local solubility.
Its locally soluble parameters have density $27/(4\pi^2)$ and its targets
are Zariski dense in the plane with third coordinate $2$. The failure
persists over every prescribed ring of $S$-integers and under integral
stable equivalence.

At the exceptional prime $2$, we prove that integral fiber cardinality is
determined modulo $8$, and has the Haar probability law
$(21,7,3,1)/32$ for cardinalities $0,1,2,3$. Consequently the $2$-adic image
has measure $11/32$. A general inverse-ball lemma turns a small finite
certificate into a theorem at every precision. In contrast, the rational
image satisfies a Hasse principle over every number field, by a cubic
Galois argument. We also give a quantitative density-zero bound for rationally
soluble integral targets and a precise formalization roadmap. All new
claims have explicit proofs; computational identities and finite tables
are supplied as replayable certificates, not represented as Lean proofs.
\end{abstract}
\vfill
\noindent\small\textbf{Status.} The split integral family, its parameter density,
and the dyadic refinement are candidate new contributions relative to the
sources reviewed here. Global priority has not been established. This is
not a new disproof of the Jacobian conjecture, nor a solution of the plane
Jacobian problem.

\smallskip
\noindent\wtag{} Batch~36, manuscript~03, placed in \texttt{1a1396d4d}: a
report of ProveIt's research-report collection that continues the formal
project \rpath{Algebra/JacobianConjecture}. AI-assisted, unrefereed, not
formalized; labels carry the prefix \texttt{alg:}. See
Section~\ref{alg:sub:attribution} and Appendix~\ref{alg:app:report}.
[Added 30 September 2026, batch~58:] Part~II
(Section~\ref{alg:dh:sec:relation}, labels \texttt{alg:dh:}, not
formalized) answers Research question~\ref{alg:sub:q-density} and, for
completely split fibers in square boxes, Research
question~\ref{alg:sub:q-height}.
[Added 30 September 2026, batch~64:] Part~III
(Section~\ref{alg:ns:sec:relation}, labels \texttt{alg:ns:}, not
formalized) answers Research questions~\ref{alg:sub:q-split} and, for
square boxes in all sectors, \ref{alg:sub:q-height}.
\end{titlepage}
\hypersetup{pageanchor=true}
\tableofcontents
\newpage

\section{The research question and its exact scope}\label{alg:sec:scope}

\subsection{From a formal polynomial witness to an arithmetic problem}
ProveIt contains independent Lean and Rocq/Coq developments for a
three-variable polynomial map announced by Levent Alp\"oge, including
its constant Jacobian determinant, explicit collisions, symmetry, and
stable changes of variables \cite{proveit,proveitresearch}. We use the
following compact presentation, over $\Z$:
\begin{equation}\label{alg:eq:F}
\begin{gathered}
 u=1+xy,\qquad h=u^2z+y^2(1+3u),\\
 F(x,y,z)=(P,Q,R)=(uh,\;y+3xh,\;x(5-3u-x^2z)).
\end{gathered}
\end{equation}
Thus
\begin{align*}
P&=(1+xy)^3z+y^2(1+xy)(4+3xy),\\
Q&=y+3x(1+xy)^2z+3xy^2(4+3xy),\\
R&=2x-3x^2y-x^3z.
\end{align*}
The formal identity $\det JF=-2$ makes $F$ an \emph{\'etale} morphism over
$\Z[1/2]$: there the Jacobian determinant is a unit. Over $\Z$ it is not
\'etale at $2$. This distinction will matter in the local counting theorem,
but it will \emph{not} explain away the global integral obstruction.

A collision establishes noninjectivity. It does not answer which integral
targets occur as values, how much congruences detect, or whether rational
and integral local--global principles agree. We therefore pose the concrete
question
\begin{equation}\label{alg:eq:localglobal-question}
\target\in\Z^3,\qquad
\target\in F(\Z_p^3)\ \text{for every prime }p
\quad\stackrel{?}{\Longrightarrow}\quad
\target\in F(\Z^3).
\end{equation}
We answer it negatively by explicit complete fibers. We then classify a
large family of such failures and determine the local image at the unique
bad prime.

\subsection{What is inherited, and what is contributed here}
The choice of this problem followed a review of the actual repository
statements rather than an assumption that every unexplained phenomenon
was open. In particular, Gao's Theorems 3.3--3.4 already give the generic
geometric degree and complete complex fiber stratification of this map
\cite{gao}. The archived note \emph{Exact finite-field value distribution}
in \texttt{royvanrijn/jacobian-research} already gives the cubic inverse
model and exact finite-field histograms, including characteristic $3$
\cite{finitefield}. We reproduce short proofs because the arithmetic
arguments depend on their precise boundary conditions, not because these
are being claimed as new theorems.

There is also prior work on a \emph{different} integral Jacobian-one map
with a locally integral fiber of geometric degree five having \emph{no
rational point} \cite{degreefive}. Our examples instead have three rational
points, all of which fail global integrality. Consequently they neither
supersede that stronger rational obstruction nor conflict with its
minimum-degree discussion. The two constructions isolate different
failures of local--global reasoning.

\begin{center}\small
\begin{tabular}{p{0.34\textwidth}p{0.28\textwidth}p{0.28\textwidth}}
\toprule
Statement & Role in this article & Evidence boundary\\
\midrule
Map $F$, determinant, collisions & Repository input; rederived & Prior formal developments; exact identity here\\
Complex fiber sizes $3,1,0$ & Prior geometry & Gao; self-contained root proof here\\
Odd finite-field histograms & Prior arithmetic input & Archived note; proof and replay here\\
Split integral Hasse family and density & Candidate new refinement & Explicit formulas and elementary proofs\\
$S$-integer and stable-equivalence persistence & Consequences of the family & Exact algebraic arguments\\
Dyadic law $(21,7,3,1)/32$ & Candidate new refinement & Inverse-ball theorem plus finite certificate\\
Rational Hasse principle; thin image bound & Applications, not priority claims & Cubic algebra, Chebotarev, elementary sieve\\
\bottomrule
\end{tabular}
\end{center}

Our source review is dated September 24, 2026, with ProveIt pinned to
\psha\ and the second research repository pinned to \rsha.
The review did not locate the same split-plane or dyadic statements in
the inspected material. That is a bounded literature finding, not a proof
of worldwide novelty. Nor do we claim that the question in
\eqref{alg:eq:localglobal-question} was a previously published named open
problem. It is a natural unresolved arithmetic question arising from the
repo, and the proofs below resolve it for this map.

\subsection{Attribution and the formal project}\label{alg:sub:attribution}
\wtag{} This subsection was added when the manuscript was written into
ProveIt. It makes the attribution of the inherited material precise.
\begin{itemize}[itemsep=2pt,topsep=3pt]\raggedright
\item \emph{The map} was announced by Alp\"oge. ProveIt's Jacobian project
formalizes it, in Lean 4 as \texttt{counterexample}
(\rpath{Algebra/JacobianConjecture/Lean/JacobianConjecture/Counterexample.lean},
lines 95--101, namespace \texttt{LeanProofs.JacobianCounterexample}) and
in Rocq/Coq as \texttt{counterexample}
(\rpath{Algebra/JacobianConjecture/Coq/Counterexample.v}, line 110),
with exactly the coordinates $P,Q,R$ displayed below.
\item \emph{The determinant and the collisions} are kernel-checked there:
\texttt{jacobianDet\_counterexample} ($\det JF=-2$ as a polynomial identity
over every commutative ring), the integral collision
\texttt{collision} with \texttt{collision\textsubscript{0}}${}=(-1,1,5)$,
\texttt{collision\textsubscript{1}}${}=(0,-2,-16)$ and
\texttt{collisionValue}${}=(0,-2,0)$, the rational triple collision
\texttt{counterexample\_triple\_collision} over $(-\tfrac14,0,0)$, the
mirror symmetry \texttt{counterexample\_equivariant}, the torus action
\texttt{counterexample\_scaling}
$F(x/s,sy,s^2z)=(s^2P,sQ,R/s)$, the one-parameter collision family
\texttt{collisionFamily}, and the refutation
\texttt{jacobianConjectureInDimensionThree\_false}; the Coq development
proves the corresponding statements (\texttt{jacobian\_det\_is\_minus\_two},
\texttt{counterexample\_scaling},
\texttt{jacobian\_conjecture\_dimension\_three\_is\_false}).
\item \emph{The complex fiber stratification} (fibers of size $3$, $1$,
$0$; \cref{alg:cor:geometric}) is Gao's \cite[Theorems 3.3--3.4]{gao}. It
is \emph{not} in ProveIt. At the write it was checked on arXiv that Gao's
Theorem~3.3 concerns this very map (component degrees $7,6,4$,
determinant $-2$, three points in the generic fiber) and that his
polynomial $c_3=27v_1^2v_3^2-18v_1v_2v_3+16v_1+v_2^3v_3-v_2^2$ is
$-\Delta/4$ in the notation of \eqref{alg:eq:discriminant}.
\item \emph{The cubic inverse model and the finite-field histograms}
(\cref{alg:thm:finitefields} and the count \eqref{alg:eq:irredcount}) are
from the archived note \cite{finitefield} of
\texttt{royvanrijn/jacobian-research}, which already uses the variable
$t=y+1/x$ and the cubic \eqref{alg:eq:inversecubic}. They are \emph{not}
in ProveIt. At the write it was checked that the pinned note states the
formulas \eqref{alg:eq:fflaw} and \eqref{alg:eq:fflaw3}.
\item \emph{Nothing arithmetic about the fibers} is in ProveIt: the
project has no material on fiber stratification, finite fields, $p$-adic
or dyadic fibers, or Hasse principles, at the pin and at the placement
commit. The theorems of \cref{alg:sec:split,alg:sec:plane,alg:sec:stability,alg:sec:rational,alg:sec:ball,alg:sec:dyadic,alg:sec:sparse}
are this manuscript's (with the priority caveats stated above), and
\emph{none of them is formalized}. The stable representatives used in
\cref{alg:prop:stable} are the project's: its degree-six one is
kernel-checked (\texttt{simplerCounterexample}), the lower-degree ones
have SymPy certificates only.
\end{itemize}

\subsection{Notation and the headline results}
We write $\target=(A,B,C)$ for a target. The letter $p$ denotes a rational
prime, while $q$ is a finite field's cardinality. Haar measure on $\Z_p^d$
is normalized to give the whole space measure one. For a finite set $S$
of primes, $\Z[S^{-1}]$ denotes the rationals whose denominators use only
primes in $S$.

The main explicit example is especially small:
\begin{equation}\label{alg:eq:headline-example}
F^{-1}(0,-4,2)(\overline\Q)=
\left\{
\left(-\tfrac12,2,36\right),
\left(\tfrac16,-4,-288\right),
\left(\tfrac13,-4,0\right)
\right\}.
\end{equation}
Every prime sees an integral point in this fiber, but no point is integral
over $\Z$. The complete proof, including exhaustiveness, is in
\cref{alg:sec:split}. The larger split-plane family and density theorem follow
in \cref{alg:sec:plane}. The dyadic law is proved in \cref{alg:sec:dyadic}:
\[
\Pr_{\target\in\Z_2^3}\bigl(\#F^{-1}(\target)(\Z_2)=j\bigr)
=\frac1{32}(21,7,3,1)_j,\qquad j=0,1,2,3.
\]
Both results depend on distinguishing actual source points from roots of
an inverse polynomial which have escaped to infinity.

\section{A complete simple-root model}\label{alg:sec:root}

\subsection{An elementary coordinate change}
The following calculation supplies all the inverse formulas needed later.
It is valid over any field of characteristic different from $2$.
On the source chart $x\ne0$, introduce
\begin{equation}\label{alg:eq:wt}
w=\frac1x,\qquad t=y+\frac1x,\qquad C=R(x,y,z).
\end{equation}
Then $u=xt$, and solving the definition of $R$ gives
\begin{equation}\label{alg:eq:sourcewt}
x=\frac1w,\qquad y=t-w,\qquad
z=5w^2-3tw-Cw^3.
\end{equation}
Direct substitution into \eqref{alg:eq:F} now yields the much smaller formulas
\begin{equation}\label{alg:eq:targetwt}
 A=wt+t^2-Ct^3,\qquad B=2w+4t-3Ct^2,\qquad R=C.
\end{equation}
For example, substitution first gives $u=t/w$ and
$h=w^2+tw-Ct^2w$, from which both first coordinates follow immediately.
Eliminating $w$ gives
\begin{equation}\label{alg:eq:inversecubic}
g_{A,B,C}(T)=CT^3-2T^2+BT-2A,
\qquad g_{A,B,C}(t)=0,
\end{equation}
and the derivative is
\begin{equation}\label{alg:eq:derivative}
d(t):=g'_{A,B,C}(t)=3Ct^2-4t+B=2w=\frac2x.
\end{equation}
In particular, every finite source point on this chart determines a
\emph{simple} root. A repeated root cannot be reconstructed by dividing
by its derivative.

\begin{proposition}[A short determinant proof]\label{alg:prop:jacobian}
As an identity in $\Z[x,y,z]$, one has $\det JF=-2$.
\end{proposition}
\begin{proof}
In the localization where $x$ is invertible, the change
$(x,y,z)\mapsto(w,t,C)$ has determinant
$(-x^{-2})\cdot1\cdot(-x^3)=x=1/w$.
From \eqref{alg:eq:targetwt}, the determinant of
$(w,t,C)\mapsto(A,B,C)$ is
\[
\det\begin{pmatrix}
t&w+2t-3Ct^2&-t^3\\
2&4-6Ct&-3t^2\\
0&0&1
\end{pmatrix}=-2w.
\]
The chain rule gives $-2$ in the localization. Since $\Z[x,y,z]$
injects into that localization, it gives the same polynomial identity
globally. This is a derivation of the identity, independent of the
accompanying symbolic expansion check.
\end{proof}

\begin{theorem}[Complete field-valued inverse]\label{alg:thm:root}
Let $k$ be a field with $\operatorname{char}k\ne2$, and
$(A,B,C)\in k^3$. Every simple root $t\in k$ of
$g_{A,B,C}$ determines exactly one preimage with $x\ne0$, namely
\begin{equation}\label{alg:eq:reconstruction}
 w=\frac{g'_{A,B,C}(t)}2,\qquad
 (x,y,z)=\left(\frac1w,\ t-w,\ 5w^2-3tw-Cw^3\right).
\end{equation}
These are all the preimages with $x\ne0$. There is a preimage with
$x=0$ exactly when $C=0$, and then it is uniquely
\begin{equation}\label{alg:eq:boundarypoint}
 (0,B,A-4B^2).
\end{equation}
Equivalently, the complete fiber is in natural bijection with the simple
$k$-rational projective roots of
\begin{equation}\label{alg:eq:binarycubic}
G_{A,B,C}(T,S)=CT^3-2T^2S+BTS^2-2AS^3.
\end{equation}
\end{theorem}
\begin{proof}
The forward construction and simplicity follow from
\eqref{alg:eq:wt}--\eqref{alg:eq:derivative}. Conversely, set $w=g'(t)/2\ne0$.
The second expression in \eqref{alg:eq:targetwt} is $B$. The first expression
minus $A$ equals $g(t)/2$ after this substitution, so it vanishes.
This proves that \eqref{alg:eq:reconstruction} is a two-sided inverse
on the chart. If $x=0$, \eqref{alg:eq:F} reduces to
$F(0,y,z)=(z+4y^2,y,0)$, proving \eqref{alg:eq:boundarypoint}.
Finally, the projective point $[1:0]$ is a root of $G$ exactly when
$C=0$. On the chart $T=1$, its derivative with respect to $S$ is $-2$,
so it is always simple. All other projective roots are the affine roots
already classified.
\end{proof}

\begin{remark}[Why a root count can be wrong]
A cubic may have three roots counted with multiplicity while the actual
fiber has only one point, or none. Also, setting $C=0$ lowers the degree
of the affine inverse polynomial but \emph{adds} the source point
\eqref{alg:eq:boundarypoint}. Both issues must be handled before making any
arithmetic inference. In particular, a resultant equation alone is not
an exhaustive inverse description.
\end{remark}

\subsection{The known geometric stratification}
The discriminant of the binary cubic, or of $g$ with its cubic coefficient
retained, is
\begin{equation}\label{alg:eq:discriminant}
\Delta=4B^2-4CB^3-64A+72CAB-108C^2A^2.
\end{equation}
Over an algebraically closed field of characteristic different from
$2$ and $3$, \cref{alg:thm:root} immediately gives the following result.

\begin{corollary}[Previously known complex fiber law]\label{alg:cor:geometric}
There are three preimages if $\Delta\ne0$, one if $\Delta=0$ but the
binary cubic is not a cube, and none on the triple-root locus
\begin{equation}\label{alg:eq:triple}
 C\ne0,\qquad 3BC=4,\qquad 27AC^2=4.
\end{equation}
There is no fiber of cardinality two over an algebraically closed field.
\end{corollary}
\begin{proof}
A separable degree-three divisor on $\PP^1$ has three simple points.
A degree-three divisor with a repeated point, but not supported at a
single point, has one double point and one simple point. A cube has no simple point. Comparing
$C(T-t_0)^3$ with $g$ gives $t_0=2/(3C)$ and then
\eqref{alg:eq:triple}. A triple root at infinity is impossible because the
coefficient of $T^2S$ is $-2$.
\end{proof}
This recovers Gao's stratification \cite[Theorem 3.4]{gao}; indeed his
polynomial $c_3$ is $-\Delta/4$. We use a different inverse coordinate
because the derivative identity \eqref{alg:eq:derivative} is especially
convenient for denominators. The geometric law does not imply that every
\emph{integral} fiber has size $0,1$, or $3$; integrality can discard just
one of three rational points.

\subsection{The projective chart at infinity}\label{alg:sec:schemechart}
For completeness, the projective-root picture can be made into an
explicit scheme description over $\Z[1/2]$. Let $Y$ be the incidence
subscheme $G=0$ in $\A^3\times\PP^1$, and let $Y^{\mathrm{simp}}$ be the
open set where the corresponding root is simple.
The map from the source is
\[
 (x,y,z)\longmapsto\bigl(F(x,y,z),[1+xy:x]\bigr).
\]
The projective pair is globally defined because $(1+xy)-yx=1$.
On $S\ne0$, \cref{alg:thm:root} gives its inverse.
On $T\ne0$, put $s=S/T$; the incidence equation becomes
\[
 C=2s-Bs^2+2As^3.
\]
Set
\[
 k=1-Bs+3As^2,\qquad Y_0=B-3As.
\]
Simplicity is precisely $k\ne0$, because the derivative of
$C-2s+Bs^2-2As^3$ is $-2k$. An inverse chart is
\begin{equation}\label{alg:eq:infinitychart}
 x=\frac{s}{k},\qquad y=Y_0,\qquad
 z=Ak^3-k(k+3)Y_0^2.
\end{equation}
Indeed $1+xy=1/k$, $h=Ak$, and the three coordinates of $F$ become
$A,B,2s-Bs^2+2As^3$. These identities prove that the two chart inverses
glue and that the source is isomorphic to $Y^{\mathrm{simp}}$.
This is a coordinate interpretation, not a separate novelty claim.
No scheme-theoretic machinery from this paragraph is needed in the
integral arguments below.

\section{An explicit infinite integral Hasse failure}\label{alg:sec:split}

\subsection{The family and all its points}
\begin{theorem}[A split three-point obstruction]\label{alg:thm:family}
For every integer $n\ge2$, let
\begin{equation}\label{alg:eq:bn}
 \target_n=(0,-2n(n-1),2).
\end{equation}
Then
\[
 \target_n\in F(\Z_p^3)\quad\text{for every prime }p,
 \qquad \target_n\in F(\Q^3)\subset F(\R^3),
 \qquad \target_n\notin F(\Z^3).
\]
The complete geometric fiber consists of three rational points. More
precisely, with
\begin{equation}\label{alg:eq:familyD}
\begin{array}{c|c|c}
i&t_i&D_i\\ \hline
0&0&-n(n-1)\\
1&n&n(2n-1)\\
2&1-n&(n-1)(2n-1)
\end{array}
\end{equation}
those points are
\begin{equation}\label{alg:eq:familypoints}
 v_i(n)=\left(D_i^{-1},\;t_i-D_i,\;
                  5D_i^2-3t_iD_i-2D_i^3\right).
\end{equation}
\end{theorem}
\begin{proof}
The inverse cubic factors exactly as
\[
 g_{\target_n}(T)=2T^3-2T^2-2n(n-1)T
              =2T(T-n)(T+n-1).
\]
For $n\ge2$, its roots $0,n,1-n$ are distinct. Half its derivative at
these roots is respectively the three values in \eqref{alg:eq:familyD}.
Since the third target coordinate is nonzero, there is no $x=0$ point.
The complete inverse theorem therefore gives exactly
\eqref{alg:eq:familypoints}, even over $\overline\Q$.

All second and third coordinates of these points are integers. Their
first coordinates are $1/D_i$, with $|D_i|>1$, so none is an integer
point. This proves global insolubility without a bounded search.

The three integers $n,n-1,2n-1$ are pairwise coprime: subtracting adjacent
integers proves the first assertion, and
$2n-(2n-1)=1$, $(2n-1)-2(n-1)=1$ prove the other two.
A prime divides at most one of them. If it divides $n$, select $v_2$;
if it divides $n-1$, select $v_1$; if it divides $2n-1$, select $v_0$.
The selected denominator is a $p$-adic unit. If it divides none, all
three points are $p$-adically integral. Thus every prime has an integral
preimage. The displayed points themselves are rational and real.
\end{proof}

\begin{example}[The smallest displayed parameter]
At $n=2$, the points are exactly those in
\eqref{alg:eq:headline-example}. At the prime $2$, use $(1/3,-4,0)$.
At the prime $3$, use $(-1/2,2,36)$. At every other prime all three
points are integral. A single rational point integral at every prime
would be a $\Z$-point, but here the locally selected point changes with
the prime. The quantifiers
\[
 \forall p\;\exists v_p
 \qquad\text{and}\qquad
 \exists v\;\forall p
\]
are genuinely different, despite the fiber having only three points.
\end{example}

\subsection{A denominator criterion for every split finite fiber}
The preceding mechanism is elementary and reusable.

\begin{lemma}[Denominator covering]\label{alg:lem:denominator}
Let $f\in\Z[X_1,\ldots,X_d]^d$ and $\target\in\Z^d$. Suppose its complete
geometric fiber is a finite nonempty set of rational points
$v_1,\ldots,v_r$. Let $d_i$ be the least positive common denominator of
the coordinates of $v_i$. Then
\begin{align*}
 \target\in f(\Z^d)&\iff d_i=1\text{ for some }i,\\
 \target\in f(\Z_p^d)&\iff p\nmid d_i\text{ for some }i,\\
 \target\in f(\Z_p^d)\text{ for every }p
   &\iff\gcd(d_1,\ldots,d_r)=1.
\end{align*}
\end{lemma}
\begin{proof}
A rational vector is integral over $\Z_p$ exactly when its reduced
coordinate denominators are all prime to $p$. Because every geometric
point is rational, extension to $\Q_p$ creates no additional points in
the fiber. Hence local membership is exactly the stated disjunction over
the existing finite set. A prime fails every disjunct precisely when it
divides their greatest common divisor. The global assertion is the
corresponding denominator-one condition.
\end{proof}
The split hypothesis matters: a nonsplit factor can acquire new points
over some completions. We will not silently apply the lemma to a cubic
with an irreducible quadratic factor.

\subsection{An actual solution of every congruence}
\begin{corollary}\label{alg:cor:congruence}
For every $n\ge2$ and every positive integer $m$, there is an integer
vector $a$ with $F(a)\equiv\target_n\pmod m$, although
$F(a)=\target_n$ has no integer solution.
\end{corollary}
\begin{proof}
For each prime power $p^e\mid m$, select a point $v_i(n)$ with $p\nmid D_i$.
Reduce its integer coordinates normally and its first coordinate by
$D_i^{-1}\bmod p^e$. This gives a source modulo $p^e$. Combine these source
coordinates using the Chinese remainder theorem. Since $F$ has integer
coefficients, its value has the required reduction at every prime power.
\end{proof}
The accompanying function \texttt{congruence\_preimage(n,m)} performs exactly
this construction. For example,
\[
 F(20587,20252,13536)\equiv(0,-4,2)\pmod{21600}.
\]
This is an exact certificate, not a numerical approximation to a nonexistent
integer solution.

\section{A dense two-parameter family and its exact density}\label{alg:sec:plane}

\subsection{Complete classification on the integral split-root locus}
The one-parameter family is the specialization $(a,b,c)=(0,n,1-n)$
of a much larger construction.

\begin{theorem}[Split-plane classification]\label{alg:thm:plane}
Let $a,b,c\in\Z$ be pairwise distinct with $a+b+c=1$, and define
\begin{equation}\label{alg:eq:planemap}
 \Phi(a,b)=\bigl(abc,\;2(ab+ac+bc),\;2\bigr),\qquad c=1-a-b.
\end{equation}
For $r\in\{a,b,c\}$ put
\[
 D_r=\prod_{s\in\{a,b,c\},\ s\ne r}(r-s).
\]
Then the complete geometric fiber of $F$ over $\Phi(a,b)$ consists of
\begin{equation}\label{alg:eq:planepoints}
 \left(D_r^{-1},\ r-D_r,\ 5D_r^2-3rD_r-2D_r^3\right),
 \qquad r=a,b,c.
\end{equation}
It has no integral point. Writing
\begin{equation}\label{alg:eq:primitivecondition}
 g=\gcd(a-b,a-c)=\gcd(a-b,3a-1),
\end{equation}
we have
\begin{equation}\label{alg:eq:alllocalcriterion}
 \Phi(a,b)\in F(\Z_p^3)\text{ for every prime }p
 \quad\Longleftrightarrow\quad g=1.
\end{equation}
More precisely, a prime has no integral preimage exactly when it divides $g$.
\end{theorem}
\begin{proof}
The coefficient identities for three roots of sum $1$ give
\[
 g_{\Phi(a,b)}(T)=2(T-a)(T-b)(T-c).
\]
Thus \cref{alg:thm:root} gives \eqref{alg:eq:planepoints}, with no boundary point.
If $|D_r|=1$, both other roots must be integers at distance one from $r$.
Because they are distinct, they must be $r-1$ and $r+1$. Their sum with
$r$ would be $3r=1$, an impossibility for integral $r$. Hence all three
first coordinates are nonintegral.

For the exact gcd, put $u=a-b$ and $v=a-c$. Up to signs the three
$D$'s are
\[
 uv,\qquad u(v-u),\qquad v(v-u).
\]
After writing $u=gu_0$, $v=gv_0$, the integers $u_0,v_0,v_0-u_0$ are
pairwise coprime. No prime can divide all three of their pairwise products.
Consequently
\begin{equation}\label{alg:eq:gcdD}
 \gcd(|D_a|,|D_b|,|D_c|)=g^2.
\end{equation}
All other coordinates in \eqref{alg:eq:planepoints} are integral. Apply
\cref{alg:lem:denominator}. Finally,
$(a-b)+(a-c)=3a-1$, proving the second expression for $g$.
\end{proof}

\begin{remark}[Exact local multiplicities in this family]
If $p\nmid g$ and exactly two roots agree modulo $p$, only the point
corresponding to the third root is integral over $\Z_p$. If all three
root residues are distinct, all three points are integral. If $p\mid g$,
no point is integral. In particular, for every locally soluble target of
this family, the integral fiber has exactly one point at $p=2$ and at
$p=3$. At $2$ three distinct residues are impossible. At $3$ three distinct
residues would sum to $0$, whereas the root sum is $1$.
\end{remark}

\subsection{A positive density theorem in the parameter plane}
Density here refers to ordered integer parameters $(a,b)$, not to all
integer targets in $\A^3$. The distinction is essential.

\begin{theorem}[Parameter density]\label{alg:thm:densityparams}
Let $M(H)$ count pairs $(a,b)\in[-H,H]^2\cap\Z^2$ such that the three
integers $a,b,1-a-b$ are distinct and
$\gcd(a-b,3a-1)=1$. As $H\to\infty$ through positive integers,
\begin{equation}\label{alg:eq:paramdensity}
 M(H)=\frac{27}{4\pi^2}(2H+1)^2+O(H\log H).
\end{equation}
Every counted pair gives a locally integral but globally nonintegral
fiber of the fixed map $F$.
\end{theorem}
\begin{proof}
First ignore coincidences among the roots, and let $L=2H+1$.
Since $3a-1\ne0$, the gcd in question is positive and at most $3H+1$.
The M\"obius identity for coprimality gives the count
\[
 \sum_{d\le3H+1}\mu(d)
 \#\{(a,b)\in[-H,H]^2:d\mid a-b,\ d\mid3a-1\}.
\]
The inner set is empty when $3\mid d$. Otherwise $3$ is invertible
modulo $d$, and there is exactly one residue pair:
\[
 a\equiv b\equiv3^{-1}\pmod d.
\]
Its count is $L^2/d^2+O(L/d+1)$, uniformly in $d$. Summing the errors gives
$O(H\log H)$, while extending the absolutely convergent main sum to
infinity costs $O(H)$. Its value is
\[
 \sum_{3\nmid d}\frac{\mu(d)}{d^2}
 =\prod_{p\ne3}\left(1-\frac1{p^2}\right)
 =\frac{1}{\zeta(2)(1-3^{-2})}
 =\frac{27}{4\pi^2}.
\]
The three root-coincidence equations
$a=b$, $2a+b=1$, and $a+2b=1$ account for only $O(H)$ pairs.
Removing them does not affect the displayed asymptotic. The last
assertion is \cref{alg:thm:plane}.
\end{proof}
The approximate value of the density is $0.683918$. Exact enumeration,
used only as a check on the proof, gives the following data.
\begin{center}
\begin{tabular}{r r r}
\toprule
$H$ & $M(H)$ & $M(H)/(2H+1)^2$\\
\midrule
10&310&0.702948\\
25&1764&0.678201\\
50&7034&0.689540\\
100&27706&0.685775\\
250&171878&0.684770\\
\bottomrule
\end{tabular}
\end{center}

\begin{corollary}[Zariski density in a fixed target plane]\label{alg:cor:zariski}
The set
\[
 \{\Phi(a,b):a,b,1-a-b\text{ distinct},\ \gcd(a-b,3a-1)=1\}
\]
is Zariski dense in the plane $C=2$. Every target in it is an integral
Hasse failure for $F$ in the sense of \eqref{alg:eq:localglobal-question}.
\end{corollary}
\begin{proof}
A nonzero polynomial in two variables has only $O(H)$ zeros in
$[-H,H]^2\cap\Z^2$: view it as a polynomial in one variable and separate
the finitely many values where all coefficient polynomials vanish.
By \cref{alg:thm:densityparams}, the allowed parameter pairs have positive
two-dimensional density, so are Zariski dense in $\A^2$.

The map $\Phi$ is dominant onto the target plane. Indeed, for any
$(A,B)\in\overline\Q^2$, the monic cubic
\[
 T^3-T^2+\frac B2T-A
\]
has three roots over $\overline\Q$, counted with multiplicity, whose sum
is $1$; those roots give a parameter pair mapping to $(A,B,2)$.
A polynomial relation vanishing on the allowed target set would therefore
pull back to a relation vanishing on a dense parameter set, hence would
vanish identically on the target plane.
\end{proof}

\begin{remark}[What the density does not say]
The map from ordered distinct root triples to a fixed target has six
permutations. Moreover, a parameter box has coefficient heights
$|A|=O(H^3)$ and $|B|=O(H^2)$, not a uniform target box. Thus
\eqref{alg:eq:paramdensity} is not a density formula in target space.
It does prove that the integral obstruction is much more than an isolated
collision or a single exceptional arithmetic progression.
\end{remark}

\section{\texorpdfstring{$S$}{S}-integers, stable equivalence, and congruence closure}\label{alg:sec:stability}

\subsection{Inverting finitely many primes does not remove the obstruction}
\begin{theorem}[Persistence over every prescribed $S$]\label{alg:thm:S}
For every finite set $S$ of rational primes, infinitely many $n\ge2$
satisfy
\[
 \target_n\in F(\Z_p^3)\ \text{for every }p,
 \qquad
 \target_n\notin F\bigl(\Z[S^{-1}]^3\bigr).
\]
In particular, this failure persists for the \'etale map over
$\Z[1/2]$.
\end{theorem}
\begin{proof}
Choose distinct primes $\ell_1,\ell_2$ outside $S$. By the Chinese
remainder theorem, there are infinitely many $n\ge2$ satisfying
$n\equiv0\pmod{\ell_1}$ and $n\equiv1\pmod{\ell_2}$.
In \eqref{alg:eq:familyD}, the denominators $D_0,D_1$ are divisible by
$\ell_1$, and $D_2$ is divisible by $\ell_2$. Their numerators are one,
so no cancellation is possible. None of the three complete rational
preimages lies in $\Z[S^{-1}]^3$. Local integrality at every prime was
proved for every $n$ in \cref{alg:thm:family}.
\end{proof}
For example, $n=6$ gives the target $(0,-60,2)$ and denominators
$-30,66,55$. None is a power of two up to sign, so this example has no
$\Z[1/2]$-point. The obstruction is therefore not a defect caused solely
by bad reduction at $2$.

\begin{corollary}[The same dense obstruction over every fixed $S$]\label{alg:cor:Sdensity}
For a fixed finite set $S$, let $M_S(H)$ count the parameter pairs in
\cref{alg:thm:densityparams} whose target has no preimage over
$\Z[S^{-1}]$. Then
\[
 M_S(H)=\frac{27}{4\pi^2}(2H+1)^2+O_S(H\log H).
\]
Their targets are Zariski dense in $C=2$ and have integral preimages over
$\Z_p$ for every prime, including the primes of $S$.
\end{corollary}
\begin{proof}
All second and third coordinates in \eqref{alg:eq:planepoints} are integers.
A preimage belongs to $\Z[S^{-1}]^3$ exactly when some $D_r$ has no prime
factor outside $S$. Then both nonzero integer differences $r-s$ and
$r-t$ have all prime factors in $S$ as well.

In the parameter box, all differences have absolute value at most $3H+1$.
The number of signed nonzero integers of this size supported on $S$ is
at most
\[
 2\prod_{p\in S}\left(1+\left\lfloor\frac{\log(3H+1)}{\log p}\right\rfloor\right)
 =O_S\bigl((\log H)^{|S|}\bigr).
\]
Given two such differences $u=r-s$, $v=r-t$, the root sum forces
$3r=1+u+v$ and hence determines the root triple, up to its finitely many
orderings. Thus only $O_S((\log H)^{2|S|})$ parameter pairs in the whole
box can give an $S$-integral point. Removing this set from the locally
soluble pairs of \cref{alg:thm:densityparams} preserves its stated asymptotic;
for fixed $S$ the polylogarithmic error is absorbed into $O_S(H\log H)$.
Positive parameter density implies target Zariski density by the proof
of \cref{alg:cor:zariski}. The all-prime local assertion was not weakened
when these pairs were removed.
\end{proof}

\subsection{Transport to the repository's stabilized maps}
\begin{proposition}[Integral stable-equivalence invariance]\label{alg:prop:stable}
Let $R$ be either $\Z$ or $\Z[S^{-1}]$, let $f\in R[X]^d$, and let
\[
 h=\beta\circ(f\times\operatorname{id}_{\A^m})\circ\alpha,
\]
where $\alpha,\beta$ are polynomial automorphisms over $R$ with polynomial
inverses over $R$. For a target $\target\in R^d$, put
$\target'=\beta(\target,0)$. Over every $R$-algebra $T$ there is a
bijection
\[
 h^{-1}(\target')(T)\simeq f^{-1}(\target)(T).
\]
Consequently integral local--global failures transport to $h$ over $R$.
\end{proposition}
\begin{proof}
Apply $\beta^{-1}$ to the target equation and then $\alpha$ to a source
point. The equation becomes
$(f\times\operatorname{id})(v,u)=(\target,0)$, which forces $u=0$ and
$f(v)=\target$. Each step is invertible over every $T$ because all
polynomial coefficients and inverse coefficients lie in $R$.
\end{proof}
The source and target shears in ProveIt's stable degree reductions are
of this form \cite{proveitresearch}. Thus the arithmetic obstruction is
not tied to the expanded degree-seven presentation. It transfers to each
such integral stable representative. The proposition is unconditional;
when applied to a particular repository reduction, the equality defining
that reduction retains the proof status documented by the repository.
We do not upgrade its symbolic certificates to kernel-checked proofs.

\subsection{The exact profinite closure}
\begin{theorem}[Closure of the integral image]\label{alg:thm:closure}
For any polynomial map $f\in\Z[X_1,\ldots,X_d]^r$,
\begin{equation}\label{alg:eq:closure}
 \overline{f(\Z^d)}^{\,\widehat\Z^r}
 =f(\widehat\Z^d)
 =\prod_p f(\Z_p^d).
\end{equation}
For the map $F$, the difference between this closure and $F(\Z^3)$,
intersected with $\Z^3$, contains every $\target_n$ and is Zariski dense
in the target plane $C=2$.
\end{theorem}
\begin{proof}
The profinite integers $\widehat\Z=\prod_p\Z_p$ are compact, and
polynomial evaluation is continuous. Hence $f(\widehat\Z^d)$ is closed.
The diagonal image of $\Z^d$ is dense by the Chinese remainder theorem,
so its image is dense in the compact image: every neighborhood of
$f(x)$ pulls back to a neighborhood of $x$ meeting $\Z^d$.
This proves the first equality. Polynomial evaluation on the product
is coordinatewise, and choices of preimages can be made separately at
each prime, proving the second equality. The final assertions follow
from \cref{alg:thm:family,alg:cor:zariski}.
\end{proof}
Thus even \emph{all} congruences, not merely all congruences below a fixed
bound, fail to characterize the integer image. On a split fiber,
\cref{alg:lem:denominator} identifies exactly what they miss: local points
can live on different rational components.

\section{The contrasting rational Hasse principle}\label{alg:sec:rational}

\subsection{A cubic obstruction at infinitely many primes}
We use one standard external theorem in this section: Chebotarev's
density theorem. In the form needed here, every conjugacy class in the
Galois group of a finite Galois extension of number fields occurs as
Frobenius at infinitely many unramified finite places. The factorization
of a separable polynomial at a good place is given by the cycle lengths
of that Frobenius action. See Milne \cite[Theorem 8.31 and Example 8.36]{milne}.
No effective version is required.

\begin{lemma}\label{alg:lem:cubicHasse}
If an irreducible cubic over a number field $K$ has a root in every finite
completion $K_v$, a contradiction follows. In fact, it has no root at
infinitely many finite places.
\end{lemma}
\begin{proof}
Its Galois group acts transitively on three roots and is either the cyclic
group of order three or the full symmetric group $S_3$. In either case
it contains a three-cycle. Chebotarev supplies infinitely many good
places with that Frobenius cycle type. The reduction is an irreducible
cubic there. After excluding the finitely many places where its monic
normalization is not integral or its discriminant is not a unit, a local
root would be integral and would reduce to a root of this irreducible
cubic. This is impossible.
\end{proof}

\begin{theorem}[Rational local--global equality]\label{alg:thm:rationalHasse}
For every number field $K$ and $\target\in K^3$,
\begin{equation}\label{alg:eq:rationalHasse}
 \target\in F(K^3)
 \quad\Longleftrightarrow\quad
 \target\in F(K_v^3)\ \text{for every finite place }v.
\end{equation}
Archimedean places can be added without changing the equality.
\end{theorem}
\begin{proof}
The forward implication is immediate. Write $\target=(A,B,C)$ and
assume local solubility at every finite place. If $C=0$, the point
$(0,B,A-4B^2)$ lies in $K^3$. Suppose $C\ne0$.
By \cref{alg:thm:root}, local solubility means that $g_{A,B,C}$ has a simple
root in every $K_v$.

An irreducible $g$ is impossible by \cref{alg:lem:cubicHasse}. A reducible
cubic has a root in $K$. If some $K$-root is simple, reconstruction
finishes the proof. If a $K$-root has multiplicity two, the remaining
linear factor is also defined over $K$ and is simple. The only remaining
possibility is a triple root, in which case no simple root exists even
geometrically and every local fiber is empty. This contradicts the
assumption. Thus a simple $K$-root always exists.
\end{proof}

\begin{remark}[Reconciliation with the integral counterexamples]
The targets $\target_n$ do have rational preimages, so they satisfy
\cref{alg:thm:rationalHasse}. What fails is simultaneous integrality of one
of those points. The earlier degree-five construction with no rational
preimage \cite{degreefive} is of a different kind: a product of an
irreducible cubic and its quadratic discriminant-field polynomial can
have a root at every completion without a rational root. A single cubic
cannot realize that covering mechanism. We claim neither a new general
Hasse theorem for arbitrary polynomial maps nor a contradiction to the
known degree-five examples.
\end{remark}

\subsection{A practical exact rational inverse algorithm}
The proof gives a short, complete rational algorithm. On input
$(A,B,C)\in\Q^3$, factor $g_{A,B,C}$ over $\Q$. Retain each rational
root of multiplicity one, reconstruct it by \eqref{alg:eq:reconstruction},
and add \eqref{alg:eq:boundarypoint} when $C=0$. To decide integer or
$S$-integer membership, inspect the reduced denominators of this finite
list. An irreducible quadratic factor can be ignored for the rational
list, but not for a local integral-image problem.

The file \texttt{code/rational\_inverse.py} implements this procedure
using exact rational arithmetic and SymPy's rational factorization.
It verifies every reconstructed point by substitution. The mathematical
completeness of the procedure is \cref{alg:thm:root}, rather than any
heuristic search bound.

\section{Known finite-field laws and their arithmetic use}\label{alg:sec:finitefields}

\subsection{Odd-characteristic fiber counts}
For a finite field $k=\F_q$ of odd characteristic, let
\[
 N_j(q)=\#\{\target\in k^3:\#F^{-1}(\target)(k)=j\}.
\]
The following formulas are prior results from \cite{finitefield}.
We include their elementary proof to fix exactly which roots are counted.

\begin{theorem}[Finite-field law, credited to the archived note]\label{alg:thm:finitefields}
If $\operatorname{char}k>3$, the only fiber sizes are $0,1,3$, and
\begin{equation}\label{alg:eq:fflaw}
 N_0=\frac{(q-1)(q^2+2)}3,\qquad
 N_1=\frac{q^3+q^2-2q+2}2,\qquad
 N_3=\frac{(q-1)(q^2+2)}6.
\end{equation}
If $\operatorname{char}k=3$, the only fiber sizes are again $0,1,3$, but
\begin{equation}\label{alg:eq:fflaw3}
 N_0=\frac{q^2(q-1)}3,\qquad
 N_1=\frac{q^2(q+1)}2,\qquad
 N_3=\frac{q^2(q-1)}6.
\end{equation}
\end{theorem}
\begin{proof}
When $C\ne0$, a cubic has $0,1$, or $3$ simple roots in its ground field:
two distinct simple rational roots force the third root to be rational
and simple. A repeated-root cubic has at most one simple root.
When $C=0$, the boundary point is always present; its affine quadratic
has either zero or two simple rational roots. This gives the same list
of possible fiber sizes.

For $C\ne0$, a three-point fiber is determined by an unordered triple
of distinct roots $\{r,s,t\}$ with nonzero sum. The sum determines
$C=2/(r+s+t)$, and the other elementary symmetric functions determine
$B,A$. For $C=0$, a three-point fiber is determined by an unordered
pair of distinct roots of $-2T^2+BT-2A$. Therefore
\begin{equation}\label{alg:eq:triplecount}
 N_3=\binom q2+\binom q3-Z_q,
\end{equation}
where $Z_q$ counts unordered distinct zero-sum triples.

There are $q^2$ ordered zero-sum triples in total. In characteristic
not three, each equality event $r=s$, $s=t$, or $t=r$ has $q$ solutions;
all their pairwise and triple intersections consist only of $(0,0,0)$.
The union has $3q-2$ elements. Hence
$6Z_q=q^2-3q+2=(q-1)(q-2)$.
In characteristic three, any equality in a zero-sum triple forces all
three entries equal, and there are $q$ such triples. Thus
$6Z_q=q^2-q$ instead. Substitution in \eqref{alg:eq:triplecount} gives
both values of $N_3$.

Finally, counting targets and source points respectively gives
$q^3=N_0+N_1+N_3$ and $q^3=N_1+3N_3$. Thus $N_0=2N_3$ and
$N_1=q^3-3N_3$, proving every formula.
\end{proof}

\subsection{A fixed-leading-coefficient refinement}
We will need one more known count from \cite{finitefield}, proved here
independently by trace counting. Fix $C\ne0$ and assume
$\operatorname{char}k>3$. As $(A,B)$ varies, $g/C$ ranges over all monic
cubics with the fixed coefficient $-2/C$ of $T^2$. The irreducible ones
number
\begin{equation}\label{alg:eq:irredcount}
 \frac{q^2-1}{3}.
\end{equation}
Indeed, the trace map $\F_{q^3}\to\F_q$ has $q^2$ elements of any prescribed
trace. Exactly one of those lies in $\F_q$, since the trace of a ground-field
element is three times that element. Every other element has degree three
and belongs to a Frobenius orbit of size three. This proves
\eqref{alg:eq:irredcount}.
Across all $C\ne0$, the number of target residue classes whose inverse
binary cubic has \emph{no} projective root is consequently
\begin{equation}\label{alg:eq:Iq}
 I(q)=\frac{(q-1)(q^2-1)}3
     =\frac{(q-1)^2(q+1)}3.
\end{equation}
These classes form a useful congruence obstruction to rational solubility.
They should not be confused with \emph{all} omitted targets of $F$ over
$k$, because a repeated projective root is still a root but contributes
no finite source point.

\section{An inverse-ball theorem at every precision}\label{alg:sec:ball}

\subsection{The exact image of a sufficiently small ball}
We now establish the analytic bridge which makes a finite dyadic
calculation conclusive. The lemma is a quantitative version of the
$p$-adic inverse-function argument, with its image lattice retained
explicitly. We give a proof rather than relying on a convergence claim.

\begin{lemma}[Exact inverse ball]\label{alg:lem:ball}
Let $f:\Z_p^d\to\Z_p^d$ be given by polynomials with coefficients in
$\Z_p$, let $a\in\Z_p^d$, and set $J=Df(a)$. Suppose
$v_p(\det J)=e<\infty$. For every integer $k>e$, the restriction of $f$
to $a+p^k\Z_p^d$ is a bijection onto
\begin{equation}\label{alg:eq:ballimage}
 f(a)+p^kJ\Z_p^d.
\end{equation}
This image is a union of target residue classes modulo $p^{k+e}$.
For every $m\ge k+e$, a target residue class modulo $p^m$ has either
zero or exactly $p^e$ preimages modulo $p^m$ within the source class
$a\pmod{p^k}$. It has $p^e$ precisely when its lifts lie in
\eqref{alg:eq:ballimage}.
\end{lemma}
\begin{proof}
Taylor expansion, with integral polynomial coefficients, gives
\[
 f(a+p^kt)=f(a)+p^kJt+p^{2k}R(t),
\]
where $R:\Z_p^d\to\Z_p^d$ is an integral polynomial map and is
$1$-Lipschitz for the maximum norm. The adjugate formula shows that
$p^eJ^{-1}$ is integral. Define
\[
 \Psi(t)=t+p^kJ^{-1}R(t).
\]
Its correction term is Lipschitz with constant at most $p^{-(k-e)}<1$.
For every $v\in\Z_p^d$, the fixed-point iteration
\[
 t\longmapsto v-p^kJ^{-1}R(t)
\]
is a contraction of the complete space $\Z_p^d$ into itself. It has a
unique fixed point, proving that $\Psi$ is bijective. The same strict
inequality gives
$\|\Psi(t)-\Psi(t')\|_p=\|t-t'\|_p$, so $\Psi$ induces a bijection
modulo every $p^r$. Since
\[
 f(a+p^kt)=f(a)+p^kJ\Psi(t),
\]
we have both injectivity and the exact image \eqref{alg:eq:ballimage}.

Also $p^e\Z_p^d\subset J\Z_p^d$, so the image contains a full target
ball of radius $p^{-(k+e)}$ around each of its points. It is therefore
a union of residue classes modulo $p^{k+e}$.

For the finite count, the Smith normal form of $J$ has diagonal entries
with valuations $e_1,\ldots,e_d\ge0$ and $\sum e_i=e$. With
$r=m-k\ge e$, the congruence $Jt\equiv v\pmod{p^r}$ has
$p^{e_1+\cdots+e_d}=p^e$ solutions whenever it is soluble.
The bijection induced by $\Psi$ leaves that count unchanged. Finally,
$p^m\Z_p^d\subset p^kJ\Z_p^d$ for $m\ge k+e$, so solubility modulo
$p^m$ is equivalent to membership in the actual image lattice. This
proves the last assertion without extrapolating from finite precision.
\end{proof}

\begin{corollary}[Uniform determinant valuation]\label{alg:cor:precision}
Suppose $v_p(\det Df(a))=e$ for every $a\in\Z_p^d$, and fix $k>e$.
Then the actual integral fiber cardinality
$J_f(\target)=\#f^{-1}(\target)(\Z_p)$ is finite and determined by
$\target\bmod p^{k+e}$. For every $m\ge k+e$,
\begin{equation}\label{alg:eq:finiteactual}
 \#\{a\bmod p^m:f(a)\equiv\target\pmod{p^m}\}
 =p^eJ_f(\target).
\end{equation}
In particular, both the image and its fiber-cardinality strata are clopen.
\end{corollary}
\begin{proof}
Partition the source into its $p^{kd}$ residue balls modulo $p^k$.
Each ball contributes either one actual preimage or none, and its image
membership depends only on the target class modulo $p^{k+e}$ by
\cref{alg:lem:ball}. Sum the point counts and the finite congruence counts
over these disjoint source balls.
\end{proof}

\subsection{The good-prime consequences}
For $F$, every odd prime has $e=0$, so we can take $k=1$.
Each source point modulo $p$ lifts uniquely above every target lift
with that residue. Thus, exactly and not just asymptotically,
\begin{equation}\label{alg:eq:oddpadiclaw}
 \Pr_{\target\in\Z_p^3}
 \bigl(\#F^{-1}(\target)(\Z_p)=j\bigr)=\frac{N_j(p)}{p^3},
 \qquad p\text{ odd},\ j=0,1,3.
\end{equation}
This is the direct $p$-adic consequence of the previously known
finite-field formulas. In particular,
\begin{equation}\label{alg:eq:oddimage}
 \rho_p:=\vol(F(\Z_p^3))=
 \begin{cases}
 7/9,&p=3,\\[2pt]
 \displaystyle\frac23+\frac1{3p}-\frac2{3p^2}+\frac2{3p^3},&p>3.
 \end{cases}
\end{equation}
The same lifting proof works over the ring of integers of an unramified
extension of $\Q_p$, with $p$ replaced in the finite-field counts by its
residue cardinality. We do not use that extension below.

\section{The exceptional prime: the exact dyadic law}\label{alg:sec:dyadic}

\subsection{From 64 source balls to 512 target classes}
At $p=2$, the determinant has constant valuation one. Taking $e=1$,
$k=2$ in \cref{alg:lem:ball} partitions the source into 64 balls
$a+4\Z_2^3$. Their images are
\begin{equation}\label{alg:eq:dyadicballs}
 F(a)+4J_F(a)\Z_2^3.
\end{equation}
Because $v_2(\det J_F(a))=1$, the matrix $J_F(a)\bmod2$ has rank two.
Thus each image in \eqref{alg:eq:dyadicballs} consists of exactly four
classes modulo $8$, specifically
\begin{equation}\label{alg:eq:finiteballimage}
 \{F(a)+4J_F(a)h\pmod8:h\in\{0,1\}^3\}.
\end{equation}
Every target in such an image has exactly one preimage in that source
ball. Therefore the desired integral fiber cardinality is obtained by
counting how many of the 64 four-element sets contain the target class.

\begin{theorem}[Exact dyadic integral law]\label{alg:thm:dyadic}
For the fixed map $F$ in \eqref{alg:eq:F}, the number
$j(\target)=\#F^{-1}(\target)(\Z_2)$ is determined by
$\target\bmod8$. Its Haar distribution is
\begin{equation}\label{alg:eq:dyadiclaw}
\begin{array}{c|rrrr}
 j&0&1&2&3\\ \hline
 \Pr(j(\target)=j)&21/32&7/32&3/32&1/32.
\end{array}
\end{equation}
In particular,
\begin{equation}\label{alg:eq:rho2}
 \vol(F(\Z_2^3))=\frac{11}{32}.
\end{equation}
For every $m\ge3$, the possible fiber sizes of
$F:(\Z/2^m\Z)^3\to(\Z/2^m\Z)^3$ are $0,2,4,6$, with counts
\begin{equation}\label{alg:eq:dyadicfinite}
 2^{3m}\left(\frac{21}{32},\frac7{32},\frac3{32},\frac1{32}\right)
\end{equation}
respectively. Modulo $8$ is necessary for image membership: modulo $4$
does not suffice.
\end{theorem}
\begin{proof}
The reduction to \eqref{alg:eq:finiteballimage} proves that a finite table
at modulus $8$ settles the infinite-precision assertion. The exact table
is below. Each row fixes the parity of $(A,B,C)$ and thus contains 64
target classes modulo $8$.
\begin{center}
\begin{tabular}{c r r r r}
\toprule
$(A,B,C)\bmod2$ & $j=0$ & $j=1$ & $j=2$ & $j=3$\\
\midrule
$(0,0,0)$&0&48&0&16\\
$(0,0,1)$&32&32&0&0\\
$(0,1,0)$&48&0&16&0\\
$(0,1,1)$&64&0&0&0\\
$(1,0,0)$&32&32&0&0\\
$(1,0,1)$&64&0&0&0\\
$(1,1,0)$&48&0&16&0\\
$(1,1,1)$&48&0&16&0\\
\midrule
Total&336&112&48&16\\
\bottomrule
\end{tabular}
\end{center}
Here is a complete finite enumeration defining and checking the total
row using only integer arithmetic. The derivative is not needed in this
version: \cref{alg:cor:precision} says that each actual preimage contributes
exactly two source residues modulo $8$.
\begin{lstlisting}[language=Python]
from collections import Counter
from itertools import product

def F(x, y, z):
    u = 1 + x*y
    h = u*u*z + y*y*(1 + 3*u)
    return u*h, y + 3*x*h, x*(5 - 3*u - x*x*z)

count = Counter(tuple(v % 8 for v in F(*a))
                for a in product(range(8), repeat=3))
assert all(c % 2 == 0 for c in count.values())
actual = {b: count[b] // 2
          for b in product(range(8), repeat=3)}
assert Counter(actual.values()) == {0:336, 1:112, 2:48, 3:16}
\end{lstlisting}
The accompanying certificate independently constructs all 64 sets in
\eqref{alg:eq:finiteballimage} using the product-rule Jacobian and verifies
equality \emph{for every target residue}, not merely agreement of the
total histogram. It also reproduces each parity row. This finite integer
enumeration is the computational part of the proof; the all-precision
part is \cref{alg:lem:ball}.

Divide the total row by 512 to obtain \eqref{alg:eq:dyadiclaw} and sum its
nonzero columns to obtain \eqref{alg:eq:rho2}. Equation
\eqref{alg:eq:finiteactual} multiplies actual cardinalities by two at every
modulus $2^m$, $m\ge3$, proving \eqref{alg:eq:dyadicfinite}.
Finally, the target classes $(0,1,0)$ and $(0,1,4)$ agree modulo $4$;
the former has two integral $2$-adic preimages and the latter none,
as certified by the same pointwise table. Hence image membership
cannot be determined modulo $4$.
\end{proof}

\subsection{A two-point integral fiber and the importance of the domain}
The target $(0,1,0)$ has three rational preimages:
\[
 (0,1,-4),\qquad (2,-1/2,5/4),\qquad (-2,1,2).
\]
Only the first and third are in $\Z_2^3$, so its dyadic integral fiber
has two points. This is compatible with the three-point rational fiber
and the known geometric law. In fact, it also has exactly two ordinary
integer preimages.

\begin{remark}[\wtag{} An intake observation: the project's collision fiber]\label{alg:rem:collisionfiber}
This remark is \emph{not} a claim of the manuscript. It was observed when
the manuscript was taken into ProveIt and checked there by exact
computation (SymPy~1.14.0: evaluation, and a lex Gr\"obner basis whose
$x$-polynomial is $x(x-1)(x+1)$). Apply \cref{alg:thm:root} to the
project's collision value $(0,-2,0)$ (\texttt{collisionValue},
\cref{alg:sub:attribution}). Since $C=0$ there is the boundary point
$(0,-2,-16)$, and $g(T)=-2T(T+1)$ has the simple roots $t=0$ and $t=-1$,
which reconstruct $(-1,1,5)$ and $(1,-2,8)$. So
\[
 F^{-1}(0,-2,0)(\overline\Q)=\{(-1,1,5),\ (0,-2,-16),\ (1,-2,8)\},
\]
three \emph{integral} points: $F$ has an integral triple collision. The
project records and formalizes only the first two
(\texttt{collision\textsubscript{0}}, \texttt{collision\textsubscript{1}}).
Along the project's family, whose parameter (called $t$ in the project
README) is written $\tau$ here to keep it apart from the root $t=y+1/x$,
$F(\tau,-1/\tau,5/\tau^2)=F(0,2/\tau,-16/\tau^2)=(0,2/\tau,0)$
and the third point is $(-\tau,2/\tau,8/\tau^2)$. The torus action
$F(x/s,sy,s^2z)=(s^2P,sQ,R/s)$ at $s=-\tfrac12$ carries these three points
to $(2,-\tfrac12,\tfrac54)$, $(0,1,-4)$ and $(-2,1,2)$, the fiber over
$(0,1,0)$ displayed above; its two integral points have largest coordinate
$4$, against $16$ for the project's integral collision.
\end{remark}

There are three different distributions that must not be interchanged:
\begin{center}
\begin{tabular}{p{0.30\textwidth}p{0.22\textwidth}p{0.33\textwidth}}
\toprule
Source and target & Fiber sizes & Image fraction or measure\\
\midrule
$\F_2^3$ & $0,1,3$ & $6/8=3/4$\\
$(\Z/4\Z)^3$ & $0,2,6$ & $28/64=7/16$\\
$\Z_2^3$ & $0,1,2,3$ & $11/32$\\
$(\Z/2^m\Z)^3$, $m\ge3$ & $0,2,4,6$ & $11/32$\\
\bottomrule
\end{tabular}
\end{center}
The actual modular histograms for moduli $2,4,8$ are respectively
\[
 (N_0,N_1,N_3)=(2,5,1),\quad
 (N_0,N_2,N_6)=(36,26,2),\quad
 (N_0,N_2,N_4,N_6)=(336,112,48,16).
\]
Also, passing to a finite field $\F_{2^r}$ is not the same operation as
passing to the ring $\Z/2^r\Z$. The characteristic-two finite-field
formulas in \cite{finitefield} do not by themselves imply the dyadic law.

As a consistency check, the expected integral dyadic fiber size is
\[
 \frac7{32}+2\frac3{32}+3\frac1{32}=\frac12.
\]
This also follows from summing the contributions of 64 source balls,
each with an image of measure $1/128$. The corresponding expected
fiber size at an odd prime is one. These are exactly the factors suggested
by the determinant's local absolute value, but the argument here is a
finite lattice count rather than an appeal to a change-of-variables formula.

\section{Sparse rational images and a measure-zero congruence closure}\label{alg:sec:sparse}

\subsection{A closed subset of profinite measure zero}
\begin{theorem}\label{alg:thm:measurezero}
The compact image $F(\widehat\Z^3)$ has Haar measure zero in
$\widehat\Z^3$. Nevertheless, its intersection with $\Z^3$ contains
the Zariski-dense split-plane family of \cref{alg:cor:zariski}.
\end{theorem}
\begin{proof}
By \cref{alg:thm:closure}, the image is the product of the local images.
For $p>3$, \eqref{alg:eq:oddimage} gives $\rho_p<3/4$. The image is
contained in the cylinder imposing local membership at any chosen
$N$ primes greater than three, and that cylinder has measure at most
$(3/4)^N$. Letting $N$ tend to infinity proves measure zero.
The second assertion was already proved in \cref{alg:thm:closure}.
\end{proof}
This is a density statement in a compact topological space. It neither
asserts emptiness nor contradicts Zariski density in a plane. Topological
measure, polynomial equations, and integer counting measure detect different
notions of size.

\subsection{A quantitative bound for rationally soluble integral targets}
Let
\[
 \mathcal A(H)=\#\bigl(F(\Q^3)\cap[-H,H]^3\cap\Z^3\bigr),
 \qquad H\ge1.
\]
Unlike an integer-image sieve, a sieve for $F(\Q^3)$ cannot simply exclude
all residue classes outside $F(\F_p^3)$: a rational source may have a
$p$-adic denominator. The correct necessary reduction condition is that
the \emph{binary} inverse cubic have a projective root modulo $p$.

\begin{lemma}\label{alg:lem:reductionroot}
If $(A,B,C)\in\Z^3\cap F(\Q^3)$, then $G_{A,B,C}$ has a projective
root over $\F_p$ for every prime $p>2$.
\end{lemma}
\begin{proof}
The rational source gives a rational projective root by \cref{alg:thm:root}.
Represent it by coprime integers $(T,S)$. At every prime they are not
both zero after reduction, and the homogeneous equation remains valid.
The coefficient $-2$ ensures that the reduced binary polynomial is not
the zero polynomial for odd $p$. The reduced root need not remain simple,
which is why the stronger exclusion would be invalid.
\end{proof}

\begin{theorem}[An elementary quantitative sieve]\label{alg:thm:thinsieve}
Let $p_1,\ldots,p_N$ be distinct primes greater than three and
$M=\prod_{i=1}^N p_i$. Then
\begin{equation}\label{alg:eq:CRTbound}
 \mathcal A(H)\le
 \left(\prod_{i=1}^N\alpha_{p_i}\right)(2H+1+M)^3,
 \qquad
 \alpha_p=\frac23+\frac1{3p}+\frac1{3p^2}-\frac1{3p^3}<\frac34.
\end{equation}
In particular, $F(\Q^3)\cap\Z^3$ has zero density in integer boxes.
There is an absolute $c>0$ such that
\begin{equation}\label{alg:eq:quantzero}
 \mathcal A(H)\ll H^3\exp(-c\sqrt{\log H}).
\end{equation}
\end{theorem}
\begin{proof}
By \eqref{alg:eq:Iq} and \cref{alg:lem:reductionroot}, a fraction
$I(p)/p^3$ of target residues is forbidden at $p$. The allowed fraction
is exactly $\alpha_p$. The Chinese remainder theorem makes the number
of allowed classes modulo $M$ equal to $M^3\prod_i\alpha_{p_i}$.
Each class contains at most $((2H+1)/M+1)^3$ points of the integer cube.
This proves \eqref{alg:eq:CRTbound}.
For $p\ge5$, the positive correction
$1/(3p)+1/(3p^2)$ is at most $2/25<1/12$, proving
$\alpha_p<3/4$.

For density zero, first fix $N$, divide by $(2H+1)^3$, and let
$H\to\infty$. The upper limit is at most $(3/4)^N$; then let $N\to\infty$.
For the stated quantitative estimate, use the elementary Bertrand bound
that between $n$ and $2n$ there is a prime for every integer $n>1$.
Taking the primes in increasing order starting at five gives
$p_i<2^{i+2}$, and hence
\[
 M<2^{N(N+5)/2}.
\]
Choose the largest nonnegative integer $N$ with
$N(N+5)/2\le\log_2 H$. Then $M\le H$ and
$N=\sqrt{2\log H/\log2}+O(1)$. Substituting into
\eqref{alg:eq:CRTbound} gives
\[
 \mathcal A(H)\ll H^3
 \exp\left(-\log(4/3)\sqrt{\frac{2}{\log2}}\sqrt{\log H}\right),
\]
with an absolute multiplicative constant, proving the claim. Only the
quantitative refinement uses Bertrand's theorem; the zero-density
conclusion uses neither it nor Chebotarev.
\end{proof}
This weak but explicit bound is not advertised as an optimal thin-set
estimate. It is useful here because it is derived directly from the same
inverse cubic and makes the necessary rational-versus-integral sieve
correction transparent.

\section{Verification artifacts and the route to formal proofs}\label{alg:sec:verification}

\subsection{What was actually executed}
The accompanying archive contains three independent-purpose checkers and
an exact rational inverse utility. They were run in the preparation of this
article, with results recorded under \texttt{certificates/}.
\wtag{} In this report that directory is shipped as \rpath{data/}, with
the same file names; two of the scripts still write to
\rpath{../certificates/} next to their own directory
(\cref{alg:app:rerun}).

\paragraph{Symbolic identities.}
\texttt{verify\_symbolic.py}, run with SymPy 1.14.0, verifies 27 exact
identities. These include $\det JF=-2$, both inverse charts, the inverse
cubic and derivative identities, the discriminant, the symbolic
one-parameter factorization, all three symbolic family branches, and
characteristic-three identities. Every zero test is a polynomial-numerator
identity after exact rational simplification; no floating-point tolerance
is involved.

\paragraph{Finite local certificates.}
\texttt{verify\_finite.py} uses only the Python standard library. It checks
factored and expanded polynomial evaluators, all three one-parameter
branches for $2\le n\le100$ with exact fractions, the complete target-wise
modulus-eight dyadic certificate in two ways, and stabilization for every
target modulo $16$ and $32$. It also reproduces the known finite-field laws
for primes from $3$ through $31$ and checks good-prime lifting at moduli
$9,27,25$. Finally, it constructs exact CRT preimages for four composite
moduli. The finite range in these checks is not the theorem's range;
the universal conclusions are proved in the preceding sections.

\paragraph{The split plane.}
\texttt{verify\_split\_plane.py} uses only exact integers and fractions.
It checks all 1600 distinct-root fibers with $(a,b)\in[-20,20]^2$,
including the exact denominator gcd $g^2$. Its parameter counts through
$H=250$ provide a check on \cref{alg:thm:densityparams}, not evidence substituted
for the M\"obius proof.

\subsection{Trust boundaries}
A successful Python or SymPy computation is not a Lean kernel proof.
The new theorems in this article have not been compiled in Lean or Rocq,
and no claim is made that they are already part of ProveIt's trusted
theorem boundary. The dyadic total is a finite, reproducible
computer-assisted enumeration plus an ordinary mathematical proof of its
all-precision significance. A fully kernel-checked version should verify
that finite enumeration with a proved checker.

Conversely, the main integral-family and parameter-density theorems do
not depend on computational discovery being trustworthy. Their proofs
use explicit factorization, a complete inverse theorem, gcd identities,
and elementary counting, all written out in the article. The scripts
supply additional independent checks of the formulas.

\subsection{A dependency-conscious formalization plan}
A practical Lean extension should proceed in four layers.

\paragraph{Layer 1: exact algebra over a field.}
Reuse the repository's definition of $F$. Formalize
\eqref{alg:eq:targetwt}, \eqref{alg:eq:derivative}, and
\eqref{alg:eq:reconstruction} over a field with $2\ne0$. The central API should
state both directions of the preimage characterization and explicitly
separate $x=0$. Polynomial normalization proves the identities; the main
logical obligation is completeness, not evaluation of examples.

\paragraph{Layer 2: rational denominators and local selection.}
Formalize \cref{alg:thm:family,alg:thm:plane} over $\Q$ and $\Z$. The no-integer
point result needs only denominator arithmetic once completeness is
available. Before introducing $p$-adic numbers, prove the stronger
constructive congruence statement for every natural modulus by modular
inverses and CRT. For each prime, the rational selected branch is already
known explicitly; embedding it into $\Q_p$ should be a final transport
step rather than a search inside a local field.

\paragraph{Layer 3: the finite dyadic certificate.}
Represent source and target residues as triples of \texttt{ZMod 8}, prove
that the evaluator is reduction of the formal polynomial, and verify the
512-entry count. The semantic bridge is crucial: a theorem about an
uninterpreted list of integers is insufficient. The independent 64-ball
version can instead certify image incidence using the mod-two Jacobian.
No global search or oracle is needed at this layer.

\paragraph{Layer 4: analysis and external arithmetic.}
Formalize the inverse-ball lemma using completeness and contraction,
then connect Layer 3 to actual $\Z_2$ fibers. Haar measures of cylinders
finish the distribution theorem. The rational Hasse theorem is a separate
branch requiring a suitable Chebotarev theorem; it is not a dependency of
the explicit integral counterexamples. The parameter density and thin
sieve form further analytic-number-theory branches. Keeping these
branches separate prevents a deep external result from becoming an
unnecessary assumption of the elementary headline family.

\section{Further research questions}\label{alg:sec:questions}
The following questions are proposed research targets. They are not
represented as settled results, and their priority status would require
additional literature review before publication.

\subsection{A minimal congruence description at the prime two}\label{alg:sub:q-dyadic}
\begin{question}
Can the 512-entry dyadic certificate be compressed into a short list of
explicit congruence conditions on $(A,B,C)$, with a hand-checkable
case proof for each integral fiber size?
\end{question}
\noindent[Added 30 September 2026, batch~64: partly answered, for image
membership only, by Part~III: Proposition~\ref{alg:ns:prop:two} compresses
membership in $F(\Z_2^3)$ to the five-row table $E(A,B,C)$ of congruences
modulo $8$. The integral fiber sizes $j=1,2,3$ are not compressed, and
that proof is still the exhaustive $512$-case check.]

\cref{alg:thm:dyadic} proves that modulus eight suffices and modulus four
does not. Thus the remaining problem is logical and algebraic compression,
not an unbounded precision search. The eight parity rows suggest a small
branching structure. A compact description would improve both exposition
and formal verification.

\subsection{All locally soluble integral fibers, not just split integral roots}\label{alg:sub:q-split}
\begin{question}
Classify the integral targets with three rational preimages, and then
those with a single rational preimage and a quadratic pair, according to
local integral solubility at every prime.
\end{question}
\noindent[Added 30 September 2026, batch~58: the completely split sector
is classified in Part~II (Theorem~\ref{alg:dh:thm:local-complete}, with
Theorem~\ref{alg:dh:thm:finite}), in terms of the root numerators; the
rational-plus-quadratic sector remains open (Research
question~\ref{alg:dh:q:quadratic}).]
[Added 30 September 2026, batch~64: answered by Part~III on every slice
$C\ne0$, in both sectors and in the target coefficients:
Theorem~\ref{alg:ns:thm:main} is a complete local certificate (an integer
root of the monic inverse polynomial, $\gcd(3BC-4,27AC^2-4)$ a power of
two, and a table modulo~$8$), and Corollary~\ref{alg:ns:cor:nonsplit} the
complete Hasse test for a rational point with an irreducible quadratic
pair. The rational-plus-quadratic sector is no longer open.]

For arbitrary split rational fibers, \cref{alg:lem:denominator} already gives
a criterion once denominators are known. The challenge is to express it
in target coefficients without factoring every fiber individually.
The quadratic case is qualitatively different: the nonrational pair can
supply local integral points at primes where the rational point fails
integrality. A complete classification must record both factorization
and valuations of the reconstruction derivative.

\subsection{Target-height asymptotics for integral Hasse failures}\label{alg:sub:q-height}
\begin{question}
What is the asymptotic number of locally integral but globally nonintegral
targets on $C=2$ in boxes $|A|\le H_A$, $|B|\le H_B$?
\end{question}
\noindent[Added 30 September 2026, batch~58: for completely split fibers
and square boxes $|A|,|B|\le T$ Part~II proves
$\#\Hsplit_c(T)\sim\kappa(c)T^{2/3}$ for every fixed $c\ne0$, with
$\kappa(2)=27\Gamma(1/3)^2/(8\pi^2\Gamma(2/3))$
(Theorem~\ref{alg:dh:thm:headline-count}); fibers with one rational point
and a quadratic pair, unequal box sides, uniformity in $c$ and an error
term remain open.]
[Added 30 September 2026, batch~64: answered by Part~III for square boxes
in all sectors. On $C=2$ the locally integral targets, all integral
Hasse failures and the nonsplit ones each number
$\frac{27}{2\pi^2}T\log T+O(T)$ in $|A|,|B|\le T$
(Theorem~\ref{alg:ns:thm:count-main}), and the integral targets
$2^{5/3}T^{1/3}+O(1)$; the completely split sector ($T^{2/3}$) is not the
dominant one. Unequal box sides remain open (Research
question~\ref{alg:ns:q:unequal}), as do other slices, a second-order term
and a power-saving error.]

The density $27/(4\pi^2)$ counts ordered root parameters in a square.
Passing to target heights introduces a nonlinear region, a sixfold
symmetry away from the discriminant, and cusp-like ranges where one root
is small. These are genuine counting issues. Establishing a main term
with an explicit Euler product would sharpen the present Zariski-density
result substantially.

\subsection{Density in the entire three-dimensional target space}\label{alg:sub:q-density}
\begin{question}
Does the set of integral Hasse failures for this fixed $F$ form a
Zariski-dense subset of all of $\A^3$, rather than just the plane $C=2$?
\end{question}
\noindent[Added 30 September 2026, batch~58: answered affirmatively in
Part~II (Theorem~\ref{alg:dh:thm:headline-density}): the failures are
Zariski dense in every plane $C=c\ne0$, hence in $\A^3$; the plane $C=0$
has none, since $F(0,B,A-4B^2)=(A,B,0)$; the dense family persists over
every $\Z[S^{-1}]$ (Theorem~\ref{alg:dh:thm:S}).]

The present construction does not answer this: all its targets have the
same third coordinate. Rational weighted symmetries can move that
coordinate, but they can introduce new prime denominators and do not
preserve all local-integrality conditions automatically. A valid extension
must control those exceptional primes explicitly.

\subsection{Exact local laws over extensions of \texorpdfstring{$\Q_2$}{Q2}}\label{alg:sub:q-extensions}
\begin{question}
Determine integral fiber distributions over finite extensions of $\Q_2$,
including ramified extensions.
\end{question}
\noindent[Added 30 September 2026, batch~64: Part~III bears on this for
image membership only: over the valuation ring of any finite extension,
a solution modulo $8$ lifts (Lemma~\ref{alg:ns:lem:lift}), so membership
is decided in the finite quotient
(Theorem~\ref{alg:ns:thm:numberfields}). The fiber distributions, and the
least modulus deciding membership (Research
question~\ref{alg:ns:q:dyadicprec}), remain open.]

Over a discretely valued extension, the same contraction proof uses the
valuation of $2$ in that field. It therefore gives a finite precision
bound, but the number of residue classes and the rank patterns can
change. It would be particularly useful to find a uniform formula in
residue cardinality and ramification index, rather than a new exhaustive
table for each extension.

\subsection{Effective local witnesses for rational insolubility}
\begin{question}
For an integral target with irreducible inverse cubic, how small can one
guarantee a prime $p$ with no $\Q_p$-preimage, in terms of the target's
height and cubic discriminant?
\end{question}
\cref{alg:thm:rationalHasse} uses qualitative Chebotarev. An effective
specialized estimate could turn the abstract rational Hasse principle
into a practical proof-producing negative decision procedure. Unconditional
and conditional bounds should be separated, and the finite exceptional
set where the inverse polynomial is not an integral separable cubic
must be included in the implementation.

\subsection{Finer image-density estimates}
\begin{question}
Can one obtain the correct exponent or a main-term asymptotic for
$\mathcal A(H)$, rather than the elementary bound in
\cref{alg:thm:thinsieve}?
\end{question}
The projective root condition gives a uniform positive congruence loss,
but the simple CRT argument uses the moduli inefficiently. A large-sieve
approach or a direct height analysis of rational roots might yield a
substantially stronger bound. Such a result would concern the arithmetic
image of a concrete nonproper \'etale map, not merely its geometric degree.

\subsection{Minimal arithmetic complexity under stable equivalence}
\begin{question}
Among the integral stable representatives in ProveIt, which realizes
integral Hasse failures with the smallest target height, source denominator
height, degree, or ambient dimension?
\end{question}
\cref{alg:prop:stable} guarantees existence but not small formulas after
transport. Complexity should be reported as a vector of separate
measures; minimizing degree alone can raise dimension and coefficient
height. A search accompanied by exact certificates could produce useful
representatives without claiming a global optimum from a bounded search.

\subsection{A reusable formal library for finite arithmetic fibers}\label{alg:sub:q-library}
\begin{question}
Can the combination of a complete inverse certificate, denominator
covering, and the inverse-ball theorem become a generic proof-producing
workflow for polynomial maps with small geometric fibers?
\end{question}
The division of labor in \cref{alg:sec:verification} suggests a reusable
architecture: exact elimination produces a candidate inverse; algebraic
lemmas certify completeness; finite local tables are reflected; and an
analytic theorem upgrades them to all precisions. This would convert
many computational observations into theorems with an explicit trust
boundary, rather than a collection of map-specific numerical checks.

\section{Conclusion}
The original polynomial witness raises an arithmetic question quite
different from noninjectivity. For this particular map, rational local
solubility everywhere does force a rational point, but integral local
solubility everywhere does \emph{not} force an integral point. The failure
occurs in explicitly split three-point fibers, survives inversion of any
fixed finite set of primes, and fills a Zariski-dense subset of a target
plane. On that split-root parameter plane its exact density is
$27/(4\pi^2)$.

The local side is also completely explicit at the exceptional prime:
modulus eight determines the integral dyadic image and every integral
fiber size, with image measure $11/32$. This is an arithmetic refinement
of the existing cubic model, not a replacement for the known geometric
or finite-field results. The general inverse-ball argument explains why
a finite certificate can settle the infinite local question.

These results provide concrete new targets for ProveIt's formalization
program while keeping their status distinct: explicit mathematical
proofs and replayed exact computations are supplied here; new
kernel-checked formal proofs remain to be implemented.

% ==== Part II begins
%% ====================================================================
%% Batch-58 addition (manuscript 01, "Zariski-Dense Integral Hasse
%% Failures for the ProveIt Keller Map"). All labels of Part II carry the
%% prefix alg:dh: (dense Hasse). Section 14 is new ([write]); manuscript
%% Section k is printed as Section k+14 (k = 1..9), its research questions
%% as Section 24, its Conclusion as Section 25 and its Appendices A and B
%% as Sections 26 and 27.
%% ====================================================================
\clearpage
\phantomsection\part*{Part II.\quad Zariski-dense integral Hasse failures and split-fiber counting}
\addcontentsline{toc}{part}{Part II.\texorpdfstring{\quad}{ } Zariski-dense integral Hasse failures and split-fiber counting}
\fancyhead[L]{\small Part II: Zariski-dense integral Hasse failures}

\noindent\emph{Sections~1--13 and Appendices~A, B and C.1--C.3 form the
original report (Part~I) of 24~September 2026, unchanged except for dated notes
after Research questions~\ref{alg:sub:q-split}--\ref{alg:sub:q-density},
on the title page, and in Appendices~\ref{alg:app:replay}
and~\ref{alg:app:report}. Part~II was added on 30~September 2026 in
batch~58 of ProveIt's incoming-report intake, from one manuscript of
29~September 2026 written against Part~I as pinned. It answers Research
question~\ref{alg:sub:q-density} in full, and Research
questions~\ref{alg:sub:q-height} and~\ref{alg:sub:q-split} in the
completely split sector only. Section~\ref{alg:dh:sec:relation} states
exactly what is answered, fixes the notation against Part~I and lists
what is printed twice; Sections~\ref{alg:dh:sec:intro}--\ref{alg:dh:sec:repro}
are the manuscript.}

\noindent[Added 30 September 2026, batch~64: Part~III
(Section~\ref{alg:ns:sec:relation}) answers Research
question~\ref{alg:sub:q-split} in the rational-plus-quadratic sector too,
and Research question~\ref{alg:sub:q-height} for square boxes in all
sectors, where the split sector counted here is not the dominant one. It
added further dated notes to Parts~I and~II, listed in
Appendix~\ref{alg:ns:app:provenance}.]

\section{Relation to Part~I}\label{alg:dh:sec:relation}
\wtag{} This section was written when the manuscript was added to the
report.

\subsection{The manuscript}\label{alg:dh:sec:source}
Part~II is batch-58 manuscript~01, \emph{Zariski-Dense Integral Hasse
Failures for the ProveIt Keller Map: Complete split-fiber criteria and
sharp target-height asymptotics}, dated 29~September 2026, with the
author line ``AI-assisted research draft prepared for Vladimir
Reshetnikov'' (the same in its PDF metadata). It was delivered as the
archive \texttt{ProveIt\_Dense\_Integral\_Hasse} and is pinned to ProveIt
commit \dhsha; its provenance, shipped files and editorial changes are in
Appendix~\ref{alg:dh:app:provenance}. Its title-page abstract and status
note read as follows.
\begin{quote}\small
We study the arithmetic fibers of the three-variable polynomial map with
Jacobian determinant $-2$ already treated in ProveIt's Jacobian-conjecture
project. We resolve the repository report's question asking whether its
integral Hasse failures are Zariski dense in the entire target space.
An explicit polynomial family gives such failures densely in \emph{every}
integer plane $C=c\ne0$; the plane $C=0$, by contrast, has an integral
preimage over every integral target. The dense construction persists over
any prescribed ring of $S$-integers.

We also give a complete arithmetic parametrization of the distinct,
completely split rational fibers. Local solubility reduces to elementary
congruences at primes dividing $2c$ and a single greatest-common-divisor
condition elsewhere. The dyadic calculation has three exceptional low
valuation cases and a uniform higher-valuation formula. For each fixed
$c\ne0$, only finitely many completely split fibers meet $\Z^3$; all of
these exceptions admit a finite divisor enumeration. Combining these
facts with a two-dimensional lattice sieve and an explicit cusp estimate
proves the sharp asymptotic
\[
 \#\Hsplit_c(T)\sim \kappa(c)T^{2/3},
\]
where $\kappa(c)>0$ is given by a beta integral and an explicit Euler
product. In particular,
\[
 \kappa(2)=\frac{27\,\Gamma(1/3)^2}{8\pi^2\Gamma(2/3)}.
\]
Thus the target-height question is settled in the completely split sector,
\wtag{}~for square boxes $|A|,|B|\le T$,
not merely in a root-parameter square. All universal claims are proved in
the text; exact computational checks and constructive congruence witnesses
are supplied separately.

\textbf{Status and scope.} The map, the inverse cubic, and the previous
split-plane construction are inherited results. The density extension,
complete split-fiber local criteria, and sharp counting theorem are the
contributions of this draft relative to the inspected sources. Worldwide
priority is not certified. This is an unrefereed research manuscript,
not a Lean or Rocq formalization, and not a new disproof of the Jacobian
conjecture.
\end{quote}

\subsection{What Part~II answers, and what it does not}\label{alg:dh:sec:answers}
\begin{itemize}[itemsep=2pt,topsep=3pt]
\item \emph{Research question~\ref{alg:sub:q-density} is answered,
affirmatively and in full.} By Theorem~\ref{alg:dh:thm:headline-density}
the integral Hasse failures are Zariski dense in every plane $C=c\ne0$,
hence in $\A^3$. The plane $C=0$ has none: every integral target
$(A,B,0)$ has the integral preimage $(0,B,A-4B^2)$, which is Part~I's
boundary point \eqref{alg:eq:boundarypoint}. The dense family stays
globally insoluble over every $\Z[S^{-1}]$
(Theorem~\ref{alg:dh:thm:S}); this extends the Zariski density of Part~I's
Corollary~\ref{alg:cor:Sdensity} from $C=2$ to every nonzero slice
(that corollary also counts the parameters on $C=2$, which Part~II does
not do for other slices).
\item \emph{Research question~\ref{alg:sub:q-height} is answered only in
the completely split sector, and only for square boxes.} It asks for the
number of locally integral but globally nonintegral targets on $C=2$ in
boxes $|A|\le H_A$, $|B|\le H_B$.
Theorem~\ref{alg:dh:thm:headline-count} proves
$\#\Hsplit_c(T)\sim\kappa(c)T^{2/3}$ for every fixed $c\ne0$ in square
boxes $|A|,|B|\le T$, counting only fibers of three distinct rational
points; at $c=2$ the constant is \eqref{alg:dh:eq:c2constant}, and the
main term carries the explicit Euler product that the discussion of the
question asks for. Not answered: fibers with one rational point and an
irreducible quadratic pair (Research question~\ref{alg:dh:q:quadratic});
unequal box sides $H_A\ne H_B$; uniformity in $c$ (Research
question~\ref{alg:dh:q:uniform}); an error term (Research
question~\ref{alg:dh:q:remainder}). The manuscript's sentence ``This
settles the completely split sector of the prior report's Research
Question~12.3'' after Theorem~\ref{alg:dh:thm:headline-count}, and the
corresponding sentence of its abstract, are printed with the restriction
to square boxes added and marked \wtag.
[Added 30 September 2026, batch~64: this describes Part~II only. The
question is now answered for square boxes in all sectors by Part~III
(Theorem~\ref{alg:ns:thm:count-main}): on $C=2$ the integral Hasse
failures number $\frac{27}{2\pi^2}T\log T+O(T)$, almost all nonsplit, so
the split sector above is a lower-order part of that count.]
\item \emph{Research question~\ref{alg:sub:q-split} is answered in the
completely split sector.} Theorem~\ref{alg:dh:thm:local-complete} gives
the complete local criterion for every completely split integral target on
every slice $c\ne0$, in terms of the root numerators, and
Theorem~\ref{alg:dh:thm:finite} shows that on each slice only finitely
many of them have an integral point. (On $C=0$ every integral target has
an integral point.) A criterion in the target coefficients themselves,
without factoring, is Research question~\ref{alg:dh:q:coefficient}; the
rational-plus-quadratic sector remains open (Research
question~\ref{alg:dh:q:quadratic}).
[Added 30 September 2026, batch~64: both are answered by Part~III: the
coefficient criterion by Theorem~\ref{alg:ns:thm:main}, and the
rational-plus-quadratic sector by Corollary~\ref{alg:ns:cor:nonsplit}
(Section~\ref{alg:ns:sec:answers}).]
\item \emph{Research questions~\ref{alg:sub:q-dyadic}
and~\ref{alg:sub:q-library} are related, not answered.} Research
question~\ref{alg:dh:q:coefficient} (coefficient-only local certificates)
extends Research question~\ref{alg:sub:q-dyadic} from the prime two to all
primes, for split fibers; Research question~\ref{alg:dh:q:library}
(certified finite-fiber arithmetic) is essentially Research
question~\ref{alg:sub:q-library}. Both are printed with a cross-reference.
Research questions 12.5--12.8 are untouched.
\end{itemize}

\subsection{Consistency with Part~I}\label{alg:dh:sec:consistency}
Where the two Parts meet they agree. Each item was checked at the write,
by hand or by exact expansion.
\begin{itemize}[itemsep=2pt,topsep=3pt]
\item Part~II's three-parameter target \eqref{alg:dh:eq:Phi} at $c=2$ is
Part~I's two-parameter map \eqref{alg:eq:planemap}, written here
$\Phi_{\mathrm I}(a,b)=(abc',\,2(ab+ac'+bc'),\,2)$ with $c'=1-a-b$
(Part~I's third root $c$, renamed $c'$ here):
\[
 \Phi(2,a,b)=\bigl(ab-ab(a+b),\;2(a+b)-2(a^2+ab+b^2),\;2\bigr)
 =\Phi_{\mathrm I}(a,b).
\]
\item Densities. At $c=2$ coefficient integrality forces both root
numerators to be even, $a=2s$, $b=2t$, so every completely split fiber on
$C=2$ has integer roots, and the numerator lattice has index four in the
root lattice. Accordingly
$\rho(2)=\tfrac{9}{\pi^2}\cdot\tfrac{3}{16}=\tfrac{27}{16\pi^2}
=\tfrac14\cdot\tfrac{27}{4\pi^2}$ is one quarter of Part~I's parameter
density \eqref{alg:eq:paramdensity}, and Part~II's integer-root density
$\delta(2)$ of \eqref{alg:dh:eq:delta-integerroots} is $27/(4\pi^2)$
itself, as the manuscript notes.
\item At $c=2$ the two local criteria coincide. Part~II's odd-prime
condition \eqref{alg:dh:eq:gcd} reads: $\gcd(2s-2t,6s-2)=2\gcd(s-t,3s-1)$
has no odd prime factor. Its dyadic condition for $v_2(c)=1$ (exactly one
of $s,t,1-s-t$ odd) says that $\gcd(s-t,3s-1)$ is odd, since the three
roots have odd sum, and are all odd exactly when $s-t$ and $3s-1$ are
both even. Together these are Part~I's condition $\gcd(a-b,3a-1)=1$ of
Theorem~\ref{alg:thm:plane}.
\item The constant at $c=2$ is
$\kappa(2)=I\cdot8^{2/3}\rho(2)=4I\cdot\tfrac{27}{16\pi^2}
=\tfrac{27\,\Gamma(1/3)^2}{8\pi^2\Gamma(2/3)}\approx1.81235$.
\item The independent modulo-eight image of
Section~\ref{alg:dh:sec:certificates} has $176$ of $512$ target classes,
and $176/512=11/32$ is Part~I's dyadic image measure
(Theorem~\ref{alg:thm:dyadic}).
\item The plane $C=0$: Part~II's formula $F(0,B,A-4B^2)=(A,B,0)$ is
Part~I's boundary point \eqref{alg:eq:boundarypoint}; the collision fiber
of Remark~\ref{alg:rem:collisionfiber} contains its instance
$(0,-2,-16)$.
\item The table of integral split exceptions after
Theorem~\ref{alg:dh:thm:finite} lists none on $C=2$, as Part~I's
Theorem~\ref{alg:thm:plane} requires: every completely split fiber there
has distinct integer roots and no integral point.
\end{itemize}
The ten counts for $T=10^3$ and $T=10^4$ in the table of
Section~\ref{alg:dh:sec:count} were also recomputed at the write by an
independent enumeration, not the delivered program: distinct numerators
with integral coefficients and $|A|,|B|\le T$, the three points rebuilt
from the inverse chart and checked to map to the target in exact
fractions, and a target counted when no point is integral and the three
point denominators have greatest common divisor one
(Lemma~\ref{alg:lem:denominator}). All ten agree.

\subsection{Notation: Part~I against Part~II}\label{alg:dh:sec:notation}
Part~II keeps the manuscript's letters, with one exception, so several
letters mean different things in the two Parts.
Table~\ref{alg:dh:tab:notation} fixes them; within Part~II only the
Part~II column applies, and the last column prints the tempting false
reading.

\begingroup\small
\setlength{\tabcolsep}{4pt}
\begin{longtable}{@{}>{\raggedright}p{0.09\textwidth}>{\raggedright}p{0.27\textwidth}>{\raggedright}p{0.37\textwidth}>{\raggedright\arraybackslash}p{0.19\textwidth}@{}}
\caption{Symbols whose meaning differs between the Parts.}
\label{alg:dh:tab:notation}\\
\toprule
Symbol & Part I & Part II & Not to be read as\\
\midrule
\endfirsthead
\toprule
Symbol & Part I & Part II & Not to be read as\\
\midrule
\endhead
\bottomrule
\endfoot
$c$ & the third integer root, $a+b+c=1$, in \eqref{alg:eq:planemap}; the
third target coordinate is $C$ & the fixed third target coordinate,
$\target=(A,B,c)$ & a root\\
$a,b,r$ & $a,b$ integer roots on $C=2$; $r\in\{a,b,c\}$ an index &
numerators of the roots $a/c$, $b/c$, $r/c$, with $a+b+r=2$
\eqref{alg:dh:eq:numerator-roots}; in Section~\ref{alg:dh:sec:dense},
$s,t$ are integer roots and $a=cs$, $b=ct$. Also $k=4r+2$ in
\eqref{alg:dh:eq:integer-polynomial-family} and the slope $y=rx$ in
\eqref{alg:dh:eq:area} & roots\\
$\Phi$ & $\Phi(a,b)$ on $C=2$ \eqref{alg:eq:planemap} & $\Phi(c,s,t)$
\eqref{alg:dh:eq:Phi}; $\Phi(2,a,b)$ is Part~I's $\Phi(a,b)$
(Section~\ref{alg:dh:sec:consistency}) & different maps\\
$D_r$, $D_j$ & products of differences of roots & products of differences
of root numerators \eqref{alg:dh:eq:D}; $D$ also a signed divisor of $2c$
in Theorem~\ref{alg:dh:thm:finite} & the same normalization\\
$d$ & $d(t)=g'(t)$ \eqref{alg:eq:derivative}; $d_i$ common denominators
& $d=s-t$ \eqref{alg:dh:eq:duv}; a signed factor $D=dv$; a divisor in
\eqref{alg:dh:eq:exception-bound} & the derivative\\
$\rho$ & $\rho_p$, the Haar measure of $F(\Z_p^3)$
(Section~\ref{alg:sec:ball}) & $\rho(c)$, the density of locally
integral numerator pairs \eqref{alg:dh:eq:rho} & an image measure\\
$27/(4\pi^2)$ & density of ordered integer-root pairs on $C=2$
\eqref{alg:eq:paramdensity} & $\delta(2)$
\eqref{alg:dh:eq:delta-integerroots}; the numerator density is
$\rho(2)=27/(16\pi^2)$ & a density of targets\\
$T$ & the indeterminate of $g_{A,B,C}(T)$ & the indeterminate of $g$, and
the height bound $|A|,|B|\le T$; $L=T^{1/3}$ & \\
$H$ & parameter heights ($M(H)$, $\mathcal A(H)$); box sides $H_A,H_B$ in
Research question~\ref{alg:sub:q-height} & not used; the height bound is
$T$ & \\
$S$ & a finite set of primes; the second projective variable in
\eqref{alg:eq:binarycubic} & a finite set of primes
(Theorem~\ref{alg:dh:thm:S}); the manuscript's sum $S$ is printed
$\Sigma$ & a sum\\
$\Sigma$ & not used & $B_0+R_0$ in the proof of
Lemma~\ref{alg:dh:lem:dyadic-lift}; $s+t$ in
\eqref{alg:dh:eq:thirdbranch} & a summation sign\\
$h$ & $h=u^2z+y^2(1+3u)$ & the same; also the odd part of $c$
($c=2h$, $c=2^eh$) in Section~\ref{alg:dh:sec:dyadic} & \\
$I$ & $I(q)$ in \eqref{alg:eq:Iq}, a finite-field count & the beta area
$\Gamma(1/3)^2/(2\Gamma(2/3))$ \eqref{alg:dh:eq:rho}; $\Isplit_c$ the
integral split exceptions & \\
$\tau$ & the project's family parameter
(Remark~\ref{alg:rem:collisionfiber}) & the number of divisors in
\eqref{alg:dh:eq:exception-bound} & \\
$q$ & a finite field's cardinality & a prime in the proof of
Theorem~\ref{alg:dh:thm:S}; $n=2q+1$ in
\eqref{alg:dh:eq:integer-polynomial-family} & a field size\\
$R$ & the third coordinate of $F$ & the same; also the cutoff $R\ge4$ of
Lemma~\ref{alg:dh:lem:cusp} and the square $[-R,R]^2$ in
Lemma~\ref{alg:dh:lem:sieve} & \\
$x,y$ & source coordinates & the same; also the scaled numerators
$a=Lx$, $b=Ly$ in Section~\ref{alg:dh:sec:count} & \\
$e$ & an exponent, as in $p^e$ & $e=v_2(|c|)$ & \\
$f,g$ & $f$ a polynomial map (Lemma~\ref{alg:lem:denominator}); $g$ the
inverse cubic & the same $g$; also valuations $f,g$ in
Lemma~\ref{alg:dh:lem:odd-integrality} and
\eqref{alg:dh:eq:dyadic-pattern} & \\
$n,k$ & $n\ge2$ indexes the family $\target_n$
(Theorem~\ref{alg:thm:family}); $k$ a field or a precision & $n$ odd,
$k\equiv2\pmod4$, $k\ge6$ in Theorem~\ref{alg:dh:thm:family} & the
index of $\target_n$\\
$\sigma,\kappa,\lambda,\delta$ & not used & $\sigma(e)$
\eqref{alg:dh:eq:sigma}, $\kappa(c)$ \eqref{alg:dh:eq:mainasymptotic},
$\lambda(c)$ and $\delta(c)$ \eqref{alg:dh:eq:delta-integerroots} & \\
\end{longtable}
\endgroup

The one rename: the manuscript's $S$ for a sum of two half-numerators (in
the proof of Lemma~\ref{alg:dh:lem:dyadic-lift}) or of two integer roots
(before \eqref{alg:dh:eq:thirdbranch}) is printed $\Sigma$, because $S$ is
the finite set of primes of Theorem~\ref{alg:dh:thm:S} and of Part~I's
Section~\ref{alg:sec:stability}. No normalization was changed. Part~I's
Theorem~\ref{alg:thm:family} (the one-parameter family $\target_n$ on
$C=2$) and Part~II's Theorem~\ref{alg:dh:thm:family} (the two-parameter
family on every slice) are different theorems.

\subsection{What is printed twice, and why}\label{alg:dh:sec:duplicates}
The manuscript restates parts of Part~I to be self-contained. Its later
sections cite these restatements, and all of its labels are kept, so they
are printed where they stand, in a few lines each; they add nothing to
Part~I, and the manuscript does not claim them
(Section~\ref{alg:dh:sec:intro}).
\begin{itemize}[itemsep=2pt,topsep=3pt]
\item The map and its expansion \eqref{alg:dh:eq:F} are \eqref{alg:eq:F}.
\item The chart \eqref{alg:dh:eq:chart}--\eqref{alg:dh:eq:cubic} is
\eqref{alg:eq:wt}--\eqref{alg:eq:derivative} with $C$ written $c$.
Proposition~\ref{alg:dh:prop:inverse} is Part~I's
Theorem~\ref{alg:thm:root} without its projective form, with the same
proof; Proposition~\ref{alg:dh:prop:det} and its proof are Part~I's
Proposition~\ref{alg:prop:jacobian}. Both Parts reprove the determinant
that the project kernel-checks (Section~\ref{alg:dh:sec:formal}).
\item The lifting fact \eqref{alg:dh:eq:pdivc} and the modulus-eight
statement after it are special cases of Part~I's inverse-ball
Lemma~\ref{alg:lem:ball} (with Corollary~\ref{alg:cor:precision} and
Theorem~\ref{alg:thm:dyadic} at two); they are kept as a second route,
a direct contraction argument.
\item New in their generality, and printed in full: the numerator
inverse points \eqref{alg:dh:eq:numerator-inverse}, which extend Part~I's
\eqref{alg:eq:planepoints} from $C=2$ to every slice; the gcd criterion of
Proposition~\ref{alg:dh:prop:gcd}, which extends Part~I's
Theorem~\ref{alg:thm:plane} (Section~\ref{alg:dh:sec:consistency}); the
constructive congruence argument of Section~\ref{alg:dh:sec:certificates},
which is the method of Part~I's Corollary~\ref{alg:cor:congruence} applied
to every locally integral split target; and Theorem~\ref{alg:dh:thm:S},
proved by two assigned denominator primes on a dense two-parameter family,
whereas Part~I's Theorem~\ref{alg:thm:S} uses the one-parameter family
$\target_n$.
\end{itemize}

\subsection{The formal project and Part~II}\label{alg:dh:sec:formal}
The manuscript cites ProveIt's Lean file
\rpath{Algebra/JacobianConjecture/Lean/JacobianConjecture/Counterexample.lean}
\cite{dh:proveitformal} as the source of the map and names no
declaration. That file defines the map as \texttt{counterexample}
(namespace \texttt{LeanProofs.JacobianCounterexample}, lines 95--101) and
proves \texttt{jacobianDet\_counterexample}, $\det JF=-2$ over every
commutative ring (line 134); Part~II uses the map, reproves the
determinant (Proposition~\ref{alg:dh:prop:det}), and relies on no formal
statement. \emph{None of Part~II's results is formalized.} The project
has no $p$-adic, Hasse-principle or counting material, and the placement
of this report beside its Lean and Rocq developments confers no formal
status. The manuscript's formalization order (the last subsection of
Section~\ref{alg:dh:sec:certificates}) is a proposal.

\section{The problem, the results, and the attribution boundary}\label{alg:dh:sec:intro}

\subsection{The fixed polynomial map}
Throughout, the map is literally the one used in the repository, without
rescaling:
\begin{equation}\label{alg:dh:eq:F}
\begin{gathered}
 u=1+xy,\qquad h=u^2z+y^2(1+3u),\\
 F(x,y,z)=(P,Q,R)=(uh,\ y+3xh,\ x(5-3u-x^2z)).
\end{gathered}
\end{equation}
Equivalently,
\begin{align*}
 P&=(1+xy)^3z+y^2(1+xy)(4+3xy),\\
 Q&=y+3x(1+xy)^2z+3xy^2(4+3xy),\\
 R&=2x-3x^2y-x^3z.
\end{align*}
ProveIt provides independent Lean and Rocq/Coq developments for this map,
its determinant, and its collisions \cite{dh:proveitformal}. The present
work uses those definitions, but reproves the algebraic identities needed
below. In particular, no unverified assertion about the repository is
silently used as a mathematical hypothesis.

\begin{definition}
An \emph{integral Hasse failure} for $F$ is a target
$\target\in\Z^3$ such that
\[
 \target\in F(\Z_p^3)\quad\hbox{for every prime }p,
 \qquad \target\notin F(\Z^3),
\]
and the real fiber is nonempty. All failures constructed or counted in
this article have three rational preimages, so the real condition is
automatic. They are failures of an \emph{integral} local--global
principle, not failures of rational solubility.
\end{definition}

Write the target as $(A,B,c)$, reserving lowercase $c$ for its fixed
third coordinate. The notation $a,b,r$ below denotes \emph{numerators
of inverse roots}, not target coordinates. This distinction is especially
important because two different root parametrizations will be used.

\subsection{The specific repository question resolved here}
The report \emph{Arithmetic Local--Global Dichotomies for ProveIt's
Three-Variable Keller Map} \wtag{}~(Part~I of this report) proves Zariski density in the plane $C=2$ and
asks, in its Research Question~12.4, whether the integral failures are
dense in all of $\A^3$ \cite{dh:priorreport}. Its discussion correctly warns
that rational weighted symmetries can introduce new denominator primes.
We control those primes and prove a stronger statement.

\begin{theorem}[Density in every nonzero slice]\label{alg:dh:thm:headline-density}
For every nonzero integer $c$, the integral Hasse failures of $F$ are
Zariski dense in the target plane $C=c$. Consequently they are Zariski
dense in $\A^3_{\Q}$. For $C=0$, there are no integral Hasse failures:
\[
 F(0,B,A-4B^2)=(A,B,0)
 \qquad(A,B\in\Z).
\]
The dense family can be chosen to remain globally insoluble over any
prescribed $\Z[S^{-1}]$, while still being locally integral at every prime.
\end{theorem}
The proof is constructive and appears in Section~\ref{alg:dh:sec:dense}.
It does not use a density theorem for primes, Chebotarev, or a numerical
search.

\subsection{A sharp target-height theorem}
A fiber is called \emph{completely split} here if it consists of three
distinct rational points. For $c\ne0$, let
\begin{equation}\label{alg:dh:eq:Hdefinition}
\Hsplit_c(T)=\left\{(A,B)\in\Z^2:
\begin{array}{l}
 |A|,|B|\le T,\quad F^{-1}(A,B,c)\text{ is completely split},\\
 (A,B,c)\in F(\Z_p^3)\text{ for every }p,\\
 (A,B,c)\notin F(\Z^3)
\end{array}\right\}.
\end{equation}
Write $e=v_2(|c|)$ and define
\begin{equation}\label{alg:dh:eq:sigma}
\sigma(e)=
\begin{cases}
 3/8,&e=0,\\
 3/16,&e=1\text{ or }2,\\
 3/2^{2e-1},&e\ge3.
\end{cases}
\end{equation}
Set
\begin{equation}\label{alg:dh:eq:rho}
\rho(c)=\sigma(e)
 \prod_{\substack{p^\nu\Vert c\\p\text{ odd}}}\frac{3}{p^{2\nu}}
 \prod_{\substack{p\ge5\\p\nmid c}}\left(1-\frac1{p^2}\right),
\qquad
 I=\frac{\Gamma(1/3)^2}{2\Gamma(2/3)}.
\end{equation}
An empty product is one; divisibility by a negative integer is understood
through its absolute value. The infinite product converges to a positive
number.

\begin{theorem}[Sharp split-fiber counting]\label{alg:dh:thm:headline-count}
For every fixed nonzero integer $c$,
\begin{equation}\label{alg:dh:eq:mainasymptotic}
 \#\Hsplit_c(T)\sim \kappa(c)T^{2/3},\qquad
 \kappa(c)=I(2c^2)^{2/3}\rho(c)>0.
\end{equation}
The completely split fibers on $C=c$ that have an integral point form a
finite, effectively enumerable set.
\end{theorem}
This settles, \wtag{}~for square boxes $|A|,|B|\le T$, the completely
split sector of the prior report's
Research Question~12.3 about target-height asymptotics.
[\wtag{} That question, Research question~\ref{alg:sub:q-height} of
Part~I, asks for boxes $|A|\le H_A$, $|B|\le H_B$ on $C=2$; unequal box
sides are not treated here and remain open. The manuscript's sentence
has no box restriction; see Section~\ref{alg:dh:sec:answers}.] It does \emph{not}
count fibers having one rational point and a nonrational quadratic pair;
those fibers may have a different arithmetic behavior.
[Added 30 September 2026, batch~64: they do. Part~III counts them on
$C=\pm2$: in square boxes they number $\frac{27}{2\pi^2}T\log T+O(T)$
(Theorem~\ref{alg:ns:thm:count-main}), so on these slices the split sector
of this theorem is not the dominant part of the integral Hasse failures.
The theorem itself is unaffected.]

\subsection{What is old, what is proved here, and what remains unclaimed}
The inverse cubic and the geometric degree-three model are prior work.
They appear in the repository's arithmetic report, which attributes the
geometric fiber stratification to Gao and the cubic model to earlier
research notes \cite{dh:priorreport,gao}. We give a self-contained derivation
rather than claiming priority for that model. The $C=2$ integral-root
family and its root-parameter density $27/(4\pi^2)$ also belong to the
prior report.

The new work here consists of the all-slice construction, the treatment
of all rationally split fibers through numerator roots, the full
bad-prime analysis, the finiteness of globally integral split fibers,
and the target-height asymptotic including its cusp calculation.
The source review was focused, not exhaustive; the results should receive
independent mathematical review before being described as published
breakthroughs.

The repository was inspected at commit \dhsha.
The companion \rpath{SOURCES.md} records the exact files and boundaries.
\wtag{} It is shipped as \rpath{02-dense-hasse-SOURCES.md}, and the
manuscript's \texttt{STATUS.md} as \rpath{02-dense-hasse-STATUS.md}.
No assertion in this manuscript has been kernel-checked merely by being
placed next to a formal development.

\section{The complete inverse cubic, with a short determinant proof}\label{alg:dh:sec:inverse}

\subsection{A rational chart}
On the source chart $x\ne0$, put
\begin{equation}\label{alg:dh:eq:chart}
 w=\frac1x,\qquad t=y+\frac1x,\qquad c=R(x,y,z).
\end{equation}
Then
\begin{equation}\label{alg:dh:eq:source-chart}
 x=\frac1w,\qquad y=t-w,\qquad z=5w^2-3tw-cw^3.
\end{equation}
Direct substitution into \eqref{alg:dh:eq:F} gives
\begin{equation}\label{alg:dh:eq:target-chart}
 A=wt+t^2-ct^3,\qquad B=2w+4t-3ct^2,\qquad C=c.
\end{equation}
For instance, $u=t/w$ and $h=w^2+tw-ct^2w$, which immediately give
the first two coordinates.

Eliminating $w$ yields
\begin{equation}\label{alg:dh:eq:cubic}
 g_{A,B,c}(T)=cT^3-2T^2+BT-2A,
 \qquad g(t)=0,\qquad g'(t)=2w=\frac2x.
\end{equation}
The derivative condition is essential: only a simple root reconstructs
a finite source point.

\begin{proposition}[Complete inverse for $c\ne0$]\label{alg:dh:prop:inverse}
Over a field of characteristic different from two, the points of
$F^{-1}(A,B,c)$ with $c\ne0$ are in bijection with the simple roots of
$g_{A,B,c}$. The inverse is
\[
 w=g'(t)/2,\qquad
 (x,y,z)=(1/w,\ t-w,\ 5w^2-3tw-cw^3).
\]
For $c=0$, there is in addition exactly one source point on $x=0$,
namely $(0,B,A-4B^2)$.
\end{proposition}
\begin{proof}
The forward implication follows from the chart. Conversely, the second
formula in \eqref{alg:dh:eq:target-chart} becomes $B$ after setting $w=g'(t)/2$.
The first coordinate minus $A$ is $g(t)/2$, hence vanishes. This gives a
two-sided inverse on $x\ne0$. If $x=0$, evaluation of the original map
gives $(z+4y^2,y,0)$, proving the remaining assertion.
\end{proof}

\begin{proposition}[Polynomial determinant identity]\label{alg:dh:prop:det}
In $\Z[x,y,z]$, one has $\det JF=-2$.
\end{proposition}
\begin{proof}
The determinant of the coordinate change $(x,y,z)\mapsto(w,t,c)$ is
$x=1/w$. The determinant of \eqref{alg:dh:eq:target-chart}, as a map in
$(w,t,c)$, is
\[
 \det\begin{pmatrix}
 t&w+2t-3ct^2&-t^3\\
 2&4-6ct&-3t^2\\
 0&0&1
 \end{pmatrix}=-2w.
\]
The chain rule gives $-2$ after localizing at $x$. Since the polynomial
ring embeds in that localization, the identity holds globally.
\end{proof}

\begin{remark}
The map is etale over $\Z[1/2]$, but not over $\Z$ at the prime two.
None of the later arguments treats it as an everywhere etale integral
map. The dyadic exception is handled explicitly.
\end{remark}

\subsection{A minimal local lifting fact}
We will only need a standard elementary consequence of the determinant.
For odd $p$, every solution modulo $p$ lifts to a solution over $\Z_p$.
Here is the precise argument. At a residue solution represented by
$z_0\in\Z_p^3$, the matrix $JF(z_0)$ is invertible over $\Z_p$. The map
\[
 u\longmapsto u-JF(z_0)^{-1}(F(z_0+u)-\target)
\]
is a contraction on $p\Z_p^3$: its constant term lies in that ball and
the nonlinear remainder gains a factor of $p$ in differences. Completeness
produces a root. In particular,
\begin{equation}\label{alg:dh:eq:pdivc}
 p\mid c,\ p\text{ odd},\ A,B\in\Z
 \quad\Longrightarrow\quad (A,B,c)\in F(\Z_p^3),
\end{equation}
because $(0,B,A-4B^2)$ solves the equations modulo $p$.

For later computational checks, the corresponding safe statement at two
is that a solution modulo eight lifts. Indeed, $JF(z_0)^{-1}$ has entries
in $\frac12\Z_2$, and the same contraction acts on $4\Z_2^3$ when the
residual is in $8\Z_2^3$. Its nonlinear part contracts by at least $1/2$.
Thus local membership at two is equivalent to membership in the finite
image $F((\Z/8\Z)^3)$. The theoretical density calculation below does
not depend on enumerating that image.

\section{An arithmetic parametrization of every completely split fiber}\label{alg:dh:sec:param}

\subsection{Numerator roots rather than only integer roots}
Fix $c\in\Z\setminus\{0\}$. If the integer polynomial $g_{A,B,c}$
splits over $\Q$, the rational-root theorem implies that $ct$ is an
integer for every root $t$. Since the sum of the roots is $2/c$, we may
write their numerators as
\begin{equation}\label{alg:dh:eq:numerator-roots}
 a,\quad b,\quad r=2-a-b,\qquad a,b\in\Z.
\end{equation}
Define
\begin{equation}\label{alg:dh:eq:PQ}
 P_0=ab(2-a-b)=abr,\qquad
 Q_0=2(a+b)-a^2-ab-b^2=ab+ar+br.
\end{equation}
Then the target coefficients are
\begin{equation}\label{alg:dh:eq:coefficients}
 A=\frac{P_0}{2c^2},\qquad B=\frac{Q_0}{c},\qquad C=c.
\end{equation}
Consequently the exact integrality conditions are
\begin{equation}\label{alg:dh:eq:integrality-divisibility}
 2c^2\mid P_0,\qquad c\mid Q_0.
\end{equation}
Conversely, these divisibilities give an integral target whose inverse
cubic is
\begin{equation}\label{alg:dh:eq:factor}
 g(T)=c(T-a/c)(T-b/c)(T-r/c).
\end{equation}
We always exclude coincidences among $a,b,r$ when speaking of a completely
split fiber.

Put
\begin{equation}\label{alg:dh:eq:D}
 D_a=(a-b)(a-r),\qquad D_b=(b-a)(b-r),\qquad D_r=(r-a)(r-b).
\end{equation}
The three source points are explicitly
\begin{equation}\label{alg:dh:eq:numerator-inverse}
 X_j=\left(
 \frac{2c}{D_j},\quad \frac{2j-D_j}{2c},\quad
 \frac{D_j(10D_j-12j-D_j^2)}{8c^2}
 \right),\qquad j\in\{a,b,r\}.
\end{equation}
These follow from $g'(j/c)=D_j/c$ and Proposition~\ref{alg:dh:prop:inverse}.
Thus the parametrization is exhaustive, not merely a supply of examples.

\subsection{Odd primes in the leading coefficient}
The coefficient divisibilities have a particularly simple local form.

\begin{lemma}[Odd-prime integrality]\label{alg:dh:lem:odd-integrality}
Let $p$ be odd and $p^\nu\Vert c$. The conditions
\[
 v_p(P_0)\ge2\nu,\qquad v_p(Q_0)\ge\nu
\]
are equivalent to
\begin{equation}\label{alg:dh:eq:odd-residue}
 (a,b)\equiv(0,0),\ (0,2),\text{ or }(2,0)\pmod{p^\nu}.
\end{equation}
The three residue classes are disjoint and have total density
$3/p^{2\nu}$.
\end{lemma}
\begin{proof}
Modulo $p$, the monic polynomial with roots $a,b,r$ becomes
$U^3-2U^2=U^2(U-2)$. Therefore two roots are divisible by $p$ and the
third is a unit congruent to two. Call the two divisible roots $u,v$,
and put $f=v_p(u)$, $g=v_p(v)$, allowing infinite valuations for zero.
If $f\ne g$, then $v_p(Q_0)=\min(f,g)$ because the unit root multiplies
$u+v$, and $uv$ has strictly higher valuation. Thus both valuations are
at least $\nu$. If $f=g<\nu$, then $v_p(P_0)=2f<2\nu$, a contradiction.
The remaining case again has $f,g\ge\nu$. Conversely, two roots divisible
by $p^\nu$ plainly give the two required divisibilities. Their sum with
the third root is exactly two, proving \eqref{alg:dh:eq:odd-residue}.
\end{proof}
No extra local-solubility condition is needed at these odd primes,
by \eqref{alg:dh:eq:pdivc}. This is the precise control of the new denominator
primes that a rational rescaling alone would not provide.

\subsection{Odd primes away from the leading coefficient}
Let $p\nmid 2c$. The last two coordinates in
\eqref{alg:dh:eq:numerator-inverse} are automatically $p$-integral. The first is
$p$-integral exactly when $D_j$ is a unit. Among three roots, all three
products $D_j$ are divisible by $p$ exactly when all three roots have the
same residue modulo $p$. Since $a+b+r=2$, this gives the following.

\begin{proposition}[The single gcd obstruction]\label{alg:dh:prop:gcd}
For an integral, completely split target, local solubility at every odd
prime not dividing $c$ is equivalent to
\begin{equation}\label{alg:dh:eq:gcd}
 \gcd(a-b,3a-2)\text{ has no odd prime divisor outside }c.
\end{equation}
The prime three can never be an obstruction in \eqref{alg:dh:eq:gcd}.
For each $p\ge5$ not dividing $c$, the unique excluded residue class is
\[
 a\equiv b\equiv 2\cdot3^{-1}\pmod p,
\]
of density $1/p^2$ in the numerator-parameter plane.
\end{proposition}
\begin{proof}
The simultaneous congruences $a=b=r$ are equivalent to
$a-b\equiv0$ and $3a-2\equiv0$. If $p=3$, the second congruence is
impossible. At any other odd prime away from $c$, it is exactly one
class of pairs. Completeness of the inverse description ensures that
there is no further $p$-adic branch to consider.
\end{proof}

\section{The dyadic classification and its exact density}\label{alg:dh:sec:dyadic}
The prime two cannot be absorbed into the preceding gcd criterion.
This section proves the factor $\sigma(e)$ in \eqref{alg:dh:eq:sigma}, including
all exceptional low valuations.

\subsection{Odd \texorpdfstring{$c$}{c}}
If $c$ is odd, the three inverse roots lie in $\Z_2$. Write them as
$\alpha,\beta,\gamma$, with sum $2/c$. At a root $\alpha$,
\[
 w=\frac c2(\alpha-\beta)(\alpha-\gamma).
\]
For the corresponding source point to be integral, $x=1/w$ must be
integral and $y=\alpha-w$ must be integral. Together these force $w$ to
be a unit. Conversely, a unit $w$ makes all three source coordinates
integral. Hence one difference must be odd and the other must have
valuation exactly one.

The roots modulo two are either all even or have two odd roots and one
even root, because their sum is even. The first case fails. In the
second, the two odd roots must have residues one and three modulo four.
Their sum is zero modulo four, while $2/c\equiv2\pmod4$, so the even
root must be two modulo four. Multiplication by the odd number $c$
preserves the set $\{1,2,3\}$ modulo four. We have proved:
\begin{equation}\label{alg:dh:eq:odd-c-dyadic}
 c\text{ odd}:\qquad
 F^{-1}(A,B,c)(\Z_2)\ne\varnothing
 \iff \{a,b,r\}\pmod4=\{1,2,3\}.
\end{equation}
The product $abr$ is always even, so coefficient integrality introduces
no separate dyadic condition when $c$ is odd. There are six allowed
ordered pairs modulo four, giving $\sigma(0)=6/16=3/8$.

\subsection{One factor of two in \texorpdfstring{$c$}{c}}
Let $c=2h$ with $h$ odd. Coefficient integrality forces $a,b$ even,
because $Q_0$ is even only when $a,b$ are both even. Write
\begin{equation}\label{alg:dh:eq:half-numerators}
 a=2A_0,\qquad b=2B_0,\qquad r=2R_0,\qquad
 A_0+B_0+R_0=1.
\end{equation}
For $e=1$, these evenness conditions already suffice for dyadic
coefficient integrality. The inverse roots are $A_0/h,B_0/h,R_0/h$,
and $w$ at a root is $h$ times the product of the two root differences.
It is integral automatically, and its inverse is integral precisely
when both differences are odd. Since the root sum is odd, this happens
exactly when one of $A_0,B_0,R_0$ is odd and the other two are even.
The density is
\[
 \frac14\cdot\frac34=\frac3{16}=\sigma(1).
\]

\subsection{At least two factors of two}
Write $c=2^e h$, $h$ odd, and $e\ge2$. Again use
\eqref{alg:dh:eq:half-numerators}. Since
\[
 P_0=8A_0B_0R_0,\qquad
 Q_0=4(A_0B_0+A_0R_0+B_0R_0),
\]
the dyadic coefficient conditions become
\begin{equation}\label{alg:dh:eq:dyadic-integrality}
 v_2(A_0B_0R_0)\ge2e-2,\qquad
 v_2(A_0B_0+A_0R_0+B_0R_0)\ge e-2.
\end{equation}
For $e=2$, these are equivalent to exactly one of the three half-numerators
being odd. The other two are even, making their product divisible by
four; the second inequality is automatic. This has density $3/16$ in
$(a,b)$.

For $e\ge3$, there is again a unique odd half-numerator. Denote it by
$A_0$, and write the other two as $B_0,R_0$. Let their valuations be
$f,g\ge1$. If $f\ne g$, the valuation of the pair sum in
\eqref{alg:dh:eq:dyadic-integrality} is $\min(f,g)$. If $f=g$, the product
condition requires $f\ge e-1$. It follows that the allowed possibilities
are exactly
\begin{equation}\label{alg:dh:eq:dyadic-pattern}
\begin{split}
 &f,g\ge e-1;\qquad\text{or}\\
 &\{f,g\}=\{e-2,g'\}\quad\text{with }g'\ge e.
\end{split}
\end{equation}
The convention that zero has infinite valuation is harmless here.

\begin{lemma}[No further dyadic obstruction when $4\mid c$]\label{alg:dh:lem:dyadic-lift}
For a completely split integral target with $v_2(c)\ge2$, at least one
of its rational inverse points is in $\Z_2^3$.
\end{lemma}
\begin{proof}
Use the odd half-numerator $A_0$ just identified, and put
\[
 \Sigma=B_0+R_0,\qquad J=B_0R_0,\qquad
 A_0=1-\Sigma,\qquad K=(A_0-B_0)(A_0-R_0)=1-3\Sigma+2\Sigma^2+J.
\]
Here $K$ is odd. Its inverse point has coordinates
\begin{equation}\label{alg:dh:eq:dyadic-source}
 x=\frac{2^{e-1}h}{K},\qquad
 y=\frac{A_0-K}{2^{e-1}h},\qquad
 z=\frac{K(5K-3A_0-2K^2)}{2^{2e-2}h^2}.
\end{equation}
The first is integral. The identity
\[
 A_0-K=2\Sigma-2\Sigma^2-J
\]
shows the same for the second, since in the first pattern of
\eqref{alg:dh:eq:dyadic-pattern} one has $v_2(\Sigma)\ge e-1$, and in the second
one has $v_2(\Sigma)=e-2$, while always $v_2(J)\ge2e-2$.
For $e=2$ the two even half-numerators give $v_2(\Sigma)\ge1$ and
$v_2(J)\ge2$, which are also sufficient.

For the third coordinate, the exact cancellation is
\begin{align}\label{alg:dh:eq:dyadic-cancellation}
 5K-3A_0-2K^2
 ={}&J-16\Sigma^2+24\Sigma^3+12\Sigma J\\
 &{}-8\Sigma^4-8\Sigma^2J-2J^2.\notag
\end{align}
Every term has valuation at least $2e-2$ in either pattern. In particular,
when $v_2(\Sigma)=e-2$ and $e\ge3$, the potentially lowest quadratic term
$16\Sigma^2$ has valuation $2e$, not merely $2e-4$.
This factor of sixteen is the useful cancellation. The case $e=2$
follows from $v_2(\Sigma)\ge1$, $v_2(J)\ge2$. Thus $z$ is integral as well.
\end{proof}

To compute the density for $e\ge3$, fix which of the three half-numerators
is odd. The first pattern in \eqref{alg:dh:eq:dyadic-pattern} has measure
$2^{-2e+2}$ in the other two. The second has measure
\[
 2\cdot2^{-(e-1)}\cdot2^{-e}=2^{-2e+2}.
\]
The three choices of the odd root are disjoint. Finally the substitution
$(a,b)=(2A_0,2B_0)$ contributes $1/4$. The resulting density is
\[
 \frac14\cdot3\cdot2^{-2e+3}=\frac3{2^{2e-1}}=\sigma(e).
\]

\subsection{The complete local criterion}
\begin{theorem}[All local conditions for split fibers]\label{alg:dh:thm:local-complete}
Let $c\ne0$, let $a,b,r=2-a-b$ be distinct integers, and suppose
\eqref{alg:dh:eq:integrality-divisibility} holds. The associated target is in
$F(\Z_p^3)$ for every prime $p$ if and only if
\eqref{alg:dh:eq:gcd} holds and the following dyadic condition holds:
\begin{center}
\begin{tabular}{c p{0.71\textwidth}}
\toprule
$v_2(c)$ & Additional condition after coefficient integrality\\
\midrule
$0$ & The numerator residues $a,b,r$ modulo four are $1,2,3$.\\
$1$ & Among $a/2,b/2,r/2$, exactly one is odd.\\
$\ge2$ & No additional condition.\\
\bottomrule
\end{tabular}
\end{center}
The density of the combined coefficient-integrality and dyadic conditions
at two is exactly $\sigma(v_2(c))$.
\end{theorem}
\begin{proof}
At odd primes dividing $c$, use the lifting argument
\eqref{alg:dh:eq:pdivc}. At the other odd primes, use
Proposition~\ref{alg:dh:prop:gcd}. The three dyadic cases were proved above.
The inverse model is complete over $\Q_p$, so these conditions are both
necessary and sufficient, not just sufficient tests.
\end{proof}

\section{A dense polynomial family in every nonzero plane}\label{alg:dh:sec:dense}
The full numerator classification is useful for counting. For a simple
explicit construction, a different parametrization is shorter: require
two inverse roots to be integers $s,t$. The third is $2/c-s-t$.
The target is
\begin{equation}\label{alg:dh:eq:Phi}
 \Phi(c,s,t)=\left(
 st-\frac c2st(s+t),\quad
 2(s+t)-c(s^2+st+t^2),\quad c
 \right).
\end{equation}
This is integral for all integer parameters because $st(s+t)$ is even.
The parameters in Section~\ref{alg:dh:sec:param} are $a=cs$, $b=ct$.

Put
\begin{equation}\label{alg:dh:eq:duv}
 d=s-t,\qquad U=c(2s+t)-2,\qquad V=c(s+2t)-2.
\end{equation}
The first coordinates of the three inverse points are
\begin{equation}\label{alg:dh:eq:firstcoordinates}
 x_s=\frac2{dU},\qquad x_t=-\frac2{dV},\qquad
 x_0=\frac{2c}{UV}.
\end{equation}
The other coordinates follow from \eqref{alg:dh:eq:source-chart}. If desired,
the third branch can be written without denominators involving $c$:
with $\Sigma=s+t$ and $J=(2s+t)(s+2t)$,
\begin{equation}\label{alg:dh:eq:thirdbranch}
 y_0=2\Sigma-\frac c2J,\qquad
 z_0=\frac{UV}{8}\bigl(2J-36\Sigma^2+12c\Sigma J-c^2J^2\bigr).
\end{equation}
Thus primes dividing $c$ do not introduce hidden odd denominators in
these coordinates.

\begin{theorem}[A universal arithmetic family]\label{alg:dh:thm:family}
Let $c\ne0$, let $n$ be a positive odd integer, and let
$k\ge6$ satisfy $k\equiv2\pmod4$. Set
\begin{equation}\label{alg:dh:eq:family-parameters}
 s=nk,\qquad t=n(k-1).
\end{equation}
Then $\Phi(c,s,t)$ is an integral Hasse failure with exactly three
rational preimages. None of those points has an integral first coordinate.
\end{theorem}
\begin{proof}
Here $d=n$, and
\[
 \gcd(d,3cs-2)=\gcd(n,3cnk-2)=\gcd(n,2)=1.
\]
Therefore every odd prime is locally soluble: this is the root-difference
criterion away from $c$, and \eqref{alg:dh:eq:pdivc} at primes dividing $c$.
At two, if $c$ is odd, then $s\equiv2\pmod4$ and $t$ is odd, so the
three roots have residues $1,2,3$ modulo four. If $v_2(c)=1$, the two
integer roots $s,t$ have opposite parity and the third has the parity
needed for a unique odd root. If $v_2(c)\ge2$, apply
Lemma~\ref{alg:dh:lem:dyadic-lift}. These statements can also be read directly
from Theorem~\ref{alg:dh:thm:local-complete} after passing to the numerator roots
$a=cs$, $b=ct$, $r=2-cs-ct$.

The differences are nonzero. More quantitatively,
\[
 |U|\ge |c|n(3k-3)\ge15|c|n,\qquad
 |V|\ge |c|n(3k-4)\ge14|c|n.
\]
All three numbers in \eqref{alg:dh:eq:firstcoordinates} are nonzero and have
absolute value strictly less than one. Consequently no branch is
integral. Completeness follows from Proposition~\ref{alg:dh:prop:inverse},
and rationality supplies a real point.
\end{proof}

\Needspace{12\baselineskip}
\begin{example}[A small odd-slice failure]\label{alg:dh:ex:small}
Take $c=1,s=1,t=3$, which also satisfies the local conditions although it
is not the particular parametrization \eqref{alg:dh:eq:family-parameters}.
Then
\[
 \Phi(1,1,3)=(-3,-5,1),
\]
and its complete fiber is
\[
 \left(-\frac13,4,81\right),\qquad
 \left(\frac15,-2,-45\right),\qquad
 \left(\frac2{15},-\frac{19}2,-\frac{765}8\right).
\]
No point is integral. The first point is integral at every prime except
three, and the second at every prime except five. Hence these two branches
alone cover all local requirements. The obstruction is not a single bad
prime; the permissible branch changes with the prime.
\end{example}

\subsection{Dominance, not just infinitely many examples}
A direct determinant calculation gives
\begin{equation}\label{alg:dh:eq:Phi-Jacobian}
 \det\frac{\partial\Phi}{\partial(c,s,t)}
 =-\frac12(s-t)\bigl(c(2s+t)-2\bigr)\bigl(c(s+2t)-2\bigr)
 =-\frac12dUV.
\end{equation}
For fixed $c$, the same expression is the two-by-two determinant of the
first two target coordinates with respect to $s,t$. The change
$(n,k)\mapsto(nk,n(k-1))$ has determinant $n$. Thus the parametrization
in Theorem~\ref{alg:dh:thm:family} has a nonzero Jacobian polynomial.

For clarity, the elementary dominance principle being used is this:
a polynomial map in characteristic zero with nonzero square Jacobian
has algebraically independent coordinate polynomials. Indeed, a nonzero
polynomial relation of least degree would, after differentiation and
inversion of the Jacobian over the rational-function field, yield a
nonzero partial derivative relation of smaller degree, a contradiction.
A nonconstant relation has a nonzero partial derivative in characteristic
zero.

An infinite product of arithmetic progressions is Zariski dense: prove
this by induction on the number of variables, using that a nonzero
univariate polynomial has only finitely many zeros. Therefore restricting
$n$ to positive odd integers and $k$ to positive integers congruent to
two modulo four does not destroy density. This proves density in every
plane $C=c\ne0$. Taking infinitely many integer values of $c$ proves
density in $\A^3$. The formula at $c=0$ was already established.
This proves the first part of Theorem~\ref{alg:dh:thm:headline-density}.

One can make the entire family an ordinary integer-coefficient polynomial
map by writing $n=2q+1$, $k=4r+2$, where $q\ge0$ and $r\ge1$. Its target is
\begin{equation}\label{alg:dh:eq:integer-polynomial-family}
\begin{split}
 A&=n^2(2r+1)(4r+1)\bigl(2-cn(8r+3)\bigr),\\
 B&=2n(8r+3)-cn^2(48r^2+36r+7),\\
 C&=c.
\end{split}
\end{equation}
The determinant in variables $(c,q,r)$ is $-4n^2UV$, which is nonzero.

\subsection{Persistence over every prescribed ring of \texorpdfstring{$S$}{S}-integers}
\begin{theorem}\label{alg:dh:thm:S}
Fix $c\ne0$ and a finite set $S$ of primes. A Zariski-dense subset of
$C=c$ consists of targets from Theorem~\ref{alg:dh:thm:family} with no preimage
in $\Z[S^{-1}]^3$.
\end{theorem}
\begin{proof}
Choose distinct odd primes $\ell,q\notin S$ not dividing $c$, with
$q\ne3$. Impose
\[
 n\equiv0\pmod\ell,\qquad n\equiv1\pmod q,\qquad n\equiv1\pmod2,
\]
and
\[
 k\equiv2\pmod4,\qquad
 k\equiv 3^{-1}(1+2c^{-1})\pmod q.
\]
The Chinese remainder theorem supplies independent infinite arithmetic
progressions for $n$ and $k$, and we take positive values with $k\ge6$.
The prime $\ell$ divides $d=n$ but not the numerators two in the first
two fractions of \eqref{alg:dh:eq:firstcoordinates}; it excludes those branches
from $S$-integrality. At $q$, one has $U\equiv0$, $V=U-cn\not\equiv0$,
while $2c$ is a unit, so the third branch is excluded by $q$.
The original theorem still supplies local points at every prime.
The two arithmetic progressions remain Zariski dense in parameter space,
so the same dominance argument proves the asserted density.
\end{proof}

\section{Only finitely many split fibers can have an integral point}\label{alg:dh:sec:finite}
This finiteness is stronger than a density-zero assertion, and is the
reason that the later asymptotic need not subtract a second main term.

\begin{theorem}[Finite integral split exceptions]\label{alg:dh:thm:finite}
For fixed $c\ne0$, the set
\[
 \Isplit_c=\{(A,B)\in\Z^2:
 F^{-1}(A,B,c)\text{ consists of three distinct rational points and
 meets }\Z^3\}
\]
is finite. It has a complete divisor enumeration using the positive
divisors of $2|c|$.
\end{theorem}
\begin{proof}
Choose an integral source point in such a fiber and label its root
numerator $a$. By \eqref{alg:dh:eq:numerator-inverse}, its first coordinate is
$2c/D_a$, so
\begin{equation}\label{alg:dh:eq:D-divides}
 D_a=(a-b)(2a+b-2)\mid2c.
\end{equation}
Choose a nonzero signed divisor $D$ of $2c$ and a signed factorization
$D=dv$. The equations
\[
 a-b=d,\qquad 2a+b-2=v
\]
have the unique possible solution
\begin{equation}\label{alg:dh:eq:exception-parameters}
 a=(d+v+2)/3,\qquad b=(v+2-2d)/3.
\end{equation}
There are finitely many choices. Retain integral $a,b$, distinct roots,
integral coefficients \eqref{alg:dh:eq:coefficients}, and an integral point
\eqref{alg:dh:eq:numerator-inverse}. This gives an exhaustive finite procedure.
No height cutoff is used.
\end{proof}

A simple, nonoptimal upper bound for the number of candidates is
\begin{equation}\label{alg:dh:eq:exception-bound}
 4\sum_{d\mid2|c|}\tau(d)
 =4\prod_{p^\nu\Vert2c}\frac{(\nu+1)(\nu+2)}2.
\end{equation}
The actual number of targets can be much smaller because of integrality
conditions, repeated roots, and repetitions of the same target.
The exact program in the archive gives, for example,
\begin{center}
\begin{tabular}{c l}
\toprule
$c$ & Completely split targets on $C=c$ with an integral point\\
\midrule
$1$ & $(-1,-1,1)$\\
$2$ & none\\
$3$ & $(-1,-3,3)$\\
$4$ & $(0,-2,4)$\\
$8$ & none\\
$12$ & $(0,-2,12)$\\
$16$ & none\\
\bottomrule
\end{tabular}
\end{center}
These rows are finite exact enumerations justified by the theorem, not
claims inferred from a bounded brute-force search.

\begin{corollary}
On each nonzero integer slice, all but finitely many locally integral,
completely split targets are integral Hasse failures.
\end{corollary}

\section{The exact numerator-parameter density}\label{alg:dh:sec:density}
The finite-prime conditions of Sections~\ref{alg:dh:sec:param} and
\ref{alg:dh:sec:dyadic}, together with the one-class exclusions at all remaining
primes, yield the density $\rho(c)$ in \eqref{alg:dh:eq:rho}.
We need the density not just in squares, but in bounded regions suitable
for a height calculation.

\begin{lemma}[A lattice sieve in bounded regions]\label{alg:dh:lem:sieve}
Fix $c\ne0$. Let $\Omega\subset\R^2$ be bounded with boundary of
Lebesgue measure zero. As $L\to\infty$, the number of pairs
$(a,b)\in L\Omega\cap\Z^2$ satisfying coefficient integrality and local
solubility at every prime, with $a,b,2-a-b$ distinct, is
\begin{equation}\label{alg:dh:eq:sieve-asymptotic}
 \rho(c)\,\vol(\Omega)L^2+o(L^2).
\end{equation}
\end{lemma}
\begin{proof}
First impose the complete conditions at primes dividing $2c$, and the
one-class exclusions for primes $5\le p\le Y$ not dividing $c$.
All of these are residue conditions modulo a fixed integer. The Chinese
remainder theorem makes their density the product of the local densities:
\[
 \sigma(v_2(c))
 \prod_{\substack{p^\nu\Vert c\\p\text{ odd}}}\frac3{p^{2\nu}}
 \prod_{\substack{5\le p\le Y\\p\nmid c}}(1-p^{-2}).
\]
A Riemann-sum argument for each residue class gives this constant times
$\vol(\Omega)L^2+o(L^2)$. The distinctness exclusions are three affine
lines, contributing only $O_\Omega(L)$ points.

It remains to justify the infinite-prime limit. Enclose $\Omega$ in a
square $[-R,R]^2$. Apart from $a=b$, a prime causing a failure must divide
$a-b$, so it is at most $2RL$. For each $p\ge5$, the bad pairs are one
residue class and number at most $(2RL/p+1)^2$. Summing for $Y<p\le2RL$
and bounding prime sums by integer sums gives
\begin{equation}\label{alg:dh:eq:large-prime-tail}
 O_R\left(\frac{L^2}{Y}+L\log L+L\right).
\end{equation}
Thus, after division by $L^2$, the omitted tail has upper limit
$O_R(1/Y)$. Let $L$ tend to infinity and then $Y$ tend to infinity.
The Euler product converges since $\sum_p p^{-2}<\infty$.
\end{proof}

Using the classical identity $\zeta(2)=\pi^2/6$, the product away from
$c$ can be written without an infinite numerical computation:
\begin{equation}\label{alg:dh:eq:rho-finite}
 \rho(c)=\frac9{\pi^2}\,\sigma(v_2(c))
 \prod_{\substack{p^\nu\Vert c\\p\text{ odd}}}\frac3{p^{2\nu}}
 \prod_{\substack{p\mid c\\p\ge5}}(1-p^{-2})^{-1}.
\end{equation}
The prime three is absent from the sieve factor because three coincident
numerator residues would have sum zero rather than two.

\subsection{Density in the simpler two-integer-root family}
For comparison with the earlier $C=2$ report, suppose $s,t$ are the
integer roots used in \eqref{alg:dh:eq:Phi}. Their everywhere-local density in
large parameter squares is
\begin{equation}\label{alg:dh:eq:delta-integerroots}
 \delta(c)=\frac9{\pi^2}\,\lambda(c)
 \prod_{\substack{p\mid c\\p\ge5}}(1-p^{-2})^{-1},
\qquad
 \lambda(c)=\begin{cases}
 3/8,&c\text{ odd},\\
 3/4,&c\equiv2\pmod4,\\
 1,&4\mid c.
 \end{cases}
\end{equation}
At odd primes away from $3c$, the bad condition is
$s=t$, $3cs=2$ modulo $p$, again one class. For odd $c$, the dyadic
condition is that the three roots have residues $1,2,3$ modulo four.
For $c\equiv2\pmod4$, the only bad parity class is $s,t$ both odd.
For $4\mid c$, all integer-root parameters pass the dyadic test.
The same sieve proof gives \eqref{alg:dh:eq:delta-integerroots}.

In particular, $\delta(2)=27/(4\pi^2)$ recovers the prior split-plane
parameter density. This is a consistency check and an extension to other
slices, not a claim that the old $C=2$ density is new. The distinction
between $\delta(c)$ and $\rho(c)$ matters: the latter counts \emph{all}
rationally split fibers through their numerator lattice.

\section{Sharp target-height asymptotics}\label{alg:dh:sec:count}
We now prove Theorem~\ref{alg:dh:thm:headline-count}. The main issues are the
unbounded cusp in root-parameter space, the finite-to-one passage to
targets, and the removal of integral exceptions.

\subsection{The exact height domain}
Let $N_c(T)$ be the number of ordered numerator pairs satisfying all local
and integrality conditions, with distinct roots and
\begin{equation}\label{alg:dh:eq:height-domain}
 |ab(2-a-b)|\le2c^2T,\qquad
 |2(a+b)-a^2-ab-b^2|\le |c|T.
\end{equation}
Every locally integral, completely split target contributes exactly six
such ordered pairs: choose the first two of its three distinct numerator
roots in order. Therefore
\begin{equation}\label{alg:dh:eq:six-to-one}
 \#\Hsplit_c(T)=\frac16N_c(T)-\#\{(A,B)\in\Isplit_c:|A|,|B|\le T\}.
\end{equation}
The second term is $O_c(1)$ by Theorem~\ref{alg:dh:thm:finite}.

Set $L=T^{1/3}$. On any fixed bounded region in scaled coordinates
$a=Lx$, $b=Ly$, the first inequality in \eqref{alg:dh:eq:height-domain} tends to
\begin{equation}\label{alg:dh:eq:limit-domain}
 |xy(x+y)|\le2c^2.
\end{equation}
The second inequality is eventually automatic there, because its left
side is $O(L^2)$ and its right side is $|c|L^3$.
The limiting domain is unbounded, so a compact-region argument alone
would not prove the theorem.

\subsection{An explicit cusp bound}
\begin{lemma}[Uniform tail control]\label{alg:dh:lem:cusp}
For fixed $c\ne0$, the number of integer pairs satisfying
\eqref{alg:dh:eq:height-domain} and
$\max(|a|,|b|,|2-a-b|)>RL$ is
\begin{equation}\label{alg:dh:eq:cusp-tail}
 O_c(L^2/R+L^{3/2}),\qquad R\ge4.
\end{equation}
The estimate holds before imposing any congruence conditions.
\end{lemma}
\begin{proof}
The quadratic form gives a global bound. Completing the square,
\[
 Q_0=\frac43-
 \left((a-\tfrac23)^2+(a-\tfrac23)(b-\tfrac23)+(b-\tfrac23)^2\right).
\]
Hence every numerator is $O_c(T^{1/2})=O_c(L^{3/2})$.
Let $M$ be the largest absolute numerator. Since the three numerators
sum to two, when $M\ge4$ the second largest has absolute value at least
$M/4$. The product bound then forces the smallest absolute numerator to
be at most $8c^2T/M^2$.

Choose which numerator is largest and which is smallest, and choose the
sign of the largest. This costs only an absolute factor. For a fixed
integer largest numerator of absolute value $m$, there are
$O_c(T/m^2+1)$ possibilities for the smallest; the third is determined
by the sum. Summing from $RL$ to $O_c(L^{3/2})$ gives
\[
 O_c\left(T\sum_{m>RL}m^{-2}+L^{3/2}\right)
 =O_c(L^2/R+L^{3/2}).
\]
The $+1$ term is important: it includes lines on which a numerator is
zero, and shows directly that they are lower order.
\end{proof}
This argument also explains why the quadratic restriction changes no
leading constant: it is eventually inactive on each fixed scaled compact
set, and the part of the actual height domain outside such sets is
controlled by the lemma. One must not simply drop it from the exact
integer count: the cubic inequality alone contains infinite zero-product
lines. The $O_c(T^{1/2})$ term explicitly controls those lines after the
quadratic bound has been imposed.

\subsection{The beta integral}
The region
\[
 \calB=\{(x,y)\in\R^2:|xy(x+y)|\le1\}
\]
is divided by $x=0$, $y=0$, and $x+y=0$ into six cones. Permuting the
three linear forms $x,y,-x-y$, together with simultaneous negation,
identifies these cones by linear maps of determinant $\pm1$ that preserve
the absolute cubic. Each cone therefore has the same area as the positive
quadrant piece. On that piece put $y=rx$:
\begin{align}\label{alg:dh:eq:area}
 \vol\{x,y>0:xy(x+y)\le1\}
 &=\int_0^\infty\int_0^{[r(1+r)]^{-1/3}}x\,dx\,dr\notag\\
 &=\frac12\int_0^\infty r^{-2/3}(1+r)^{-2/3}\,dr\notag\\
 &=\frac12 B(1/3,1/3)
 =\frac{\Gamma(1/3)^2}{2\Gamma(2/3)}=I.
\end{align}
The integral converges both at zero and infinity. The beta--gamma identity
used here is standard \cite{dh:dlmf}, and the preceding integral is already
an independent exact description of the constant. Thus
\begin{equation}\label{alg:dh:eq:full-area}
 \vol\{(x,y):|xy(x+y)|\le2c^2\}=6I(2c^2)^{2/3}.
\end{equation}

\subsection{Completion of the asymptotic proof}
Fix $R$. On a bounded scaled region, the moving height domain converges
to \eqref{alg:dh:eq:limit-domain} except on its algebraic boundary, which has
area zero. Sandwiching it between fixed inner and outer regions and
applying Lemma~\ref{alg:dh:lem:sieve} gives the expected local density times
its limiting area. Equivalently, one can use the two fixed conditions
$|xy(x+y)|\le2c^2\pm\varepsilon$ inside a fixed square, then let
$\varepsilon$ tend to zero.

Lemma~\ref{alg:dh:lem:cusp} bounds the remaining lattice points by
$O_c(L^2/R+L^{3/2})$. The corresponding limiting continuous tail also
tends to zero, by the finite area just computed. Let $L\to\infty$ first
and then $R\to\infty$. We obtain
\[
 N_c(T)\sim6\rho(c)I(2c^2)^{2/3}T^{2/3}.
\]
Divide by six in \eqref{alg:dh:eq:six-to-one} and subtract the bounded set of
integral exceptions. This proves \eqref{alg:dh:eq:mainasymptotic}.

\subsection{Constants and reproducible counts}
The constant depends only on $|c|$. Some values are
\begin{center}
\begin{tabular}{r r r}
\toprule
$c$ & $\rho(c)$ & $\kappa(c)$\\
\midrule
$1$ & $0.3419589948$ & $1.4384663498$\\
$2$ & $0.1709794974$ & $1.8123540336$\\
$3$ & $0.1139863316$ & $2.0746274748$\\
$4$ & $0.1709794974$ & $4.5668459937$\\
$8$ & $0.0854897487$ & $5.7538653990$\\
\bottomrule
\end{tabular}
\end{center}
For $c=2$ the theorem specializes to
\begin{equation}\label{alg:dh:eq:c2constant}
 \#\Hsplit_2(T)\sim
 \frac{27\,\Gamma(1/3)^2}{8\pi^2\Gamma(2/3)}T^{2/3}.
\end{equation}
There are no integral exceptions on this slice.
The companion enumeration gives the following \emph{actual Hasse-failure}
counts, after subtracting the finite exceptions:
\begin{center}
\begin{tabular}{r r r r}
\toprule
$c$ & $T=10^3$ & $T=10^4$ & $T=10^5$\\
\midrule
$1$ & $103$ & $540$ & $2692$\\
$2$ & $121$ & $653$ & $3304$\\
$3$ & $134$ & $727$ & $3730$\\
$4$ & $282$ & $1576$ & $8118$\\
$8$ & $330$ & $1903$ & $9975$\\
\bottomrule
\end{tabular}
\end{center}
These finite data are not used to infer the asymptotic. Convergence of the
normalized count can be visibly slow because the cutoff in the cusp is
of order $T^{1/2}$, only a factor $T^{1/6}$ below the main scale
$T^{2/3}$. The theorem asserts an equivalent for fixed $c$, not a uniform
error estimate as $c$ grows.

\section{Constructive certificates and formalization boundaries}\label{alg:dh:sec:certificates}

\subsection{Every congruence has an explicit solution}
For any locally integral split target, the three rational points
\eqref{alg:dh:eq:numerator-inverse} suffice to construct a solution modulo every
positive integer $M$. Factor $M=\prod p^{e_p}$. At each $p$, choose a
branch whose coordinate denominators are coprime to $p$, reduce each
fraction using a modular inverse modulo $p^{e_p}$, and combine the
three coordinates by the Chinese remainder theorem.
Because the original polynomial has integer coefficients, the resulting
point solves $F(x,y,z)\equiv(A,B,c)\pmod M$.

This does not produce a single rational branch that is integral at all
primes. In Example~\ref{alg:dh:ex:small}, one branch is needed at three and a
different branch at five. Confusing a coherent choice of residues for
one modulus with a globally integral branch is precisely the invalid
local--global inference.

\subsection{The companion programs}
The archive contains three main programs.
\begin{description}[style=nextline,leftmargin=1.2em,labelindent=0pt]
\item[\rpath{code/split_fibres.py}]
Uses Python integers and exact fractions to reconstruct all three points,
decide the local criterion, enumerate all integral split exceptions for
fixed $c$, generate the universal family, and produce modular witnesses.
It has no third-party dependencies.
\item[\rpath{code/verify_results.py}]
Runs the exact symbolic identities, an independent image calculation
modulo eight, exhaustive finite checks of the bad-prime density factors,
odd-prime residue checks, and exact rational tests of the dense family.
It writes a structured verification report. Symbolic checks use SymPy;
vectorized finite residue checks use NumPy integers in ranges safely
within signed 64-bit arithmetic.
\item[\rpath{code/count_heights.py}]
Reproduces the target-height count table by exhaustive numerator
enumeration, exact divisibility and gcd tests, and subtraction of the
complete finite exception list. Floating-point arithmetic is used only
for displaying the theoretical constants and normalized counts, not
for membership decisions.
\end{description}
\wtag{} In this report the three programs are shipped as
\rpath{code/02-dense-hasse-split_fibres.py},
\rpath{code/02-dense-hasse-verify_results.py} and
\rpath{code/02-dense-hasse-count_heights.py}, and their recorded outputs
under \rpath{data/} with the same prefix; run them only on a copy with the
delivered names (Appendix~\ref{alg:dh:app:rerun}).

The executed verification contains twelve passing groups. In particular,
the modulo-eight image has 176 elements, and the complete bad-prime
numerator densities were independently checked for
$c=1,2,3,4,5,6,8,12,16$. Exact reconstruction and local tests were also
compared on a fixed-seed sample, and the universal family was tested on
108 signed-slice parameter choices. Finite tests supplement the proofs;
they do not establish universal statements on their own.

\subsection{An economical Lean/Rocq route}
A useful formalization order is: first the chart identities and the
complete simple-root inverse; then the numerator parametrization and
odd-prime valuation lemmas; then the explicit dense family and its
nonintegrality; then the dyadic cancellation. The constructive
congruence theorem can be formalized before introducing $p$-adic fields,
using reduced denominators, modular inverses, and CRT.

The counting theorem is a separate analytic branch. It requires lattice
point asymptotics for fixed congruence classes, the elementary large-prime
tail estimate, the cusp bound, and the beta integral. The finite integral
exception theorem needs only divisibility and the explicit inverse.
No Chebotarev theorem is needed anywhere in this paper.

The original repository declarations are not recompiled here, and the
new results have not been encoded as Lean or Rocq proofs. The package is
an English proof manuscript with reproducible exact checks.

\section{Further research questions of Part~II}\label{alg:dh:sec:questions}
The results above close the full-dimensional density question and provide
a sharp counting theorem for the split sector. They also distinguish
several genuinely different remaining problems.

\begin{question}[The rational-plus-quadratic sector]\label{alg:dh:q:quadratic}
For fixed $c\ne0$, classify integral Hasse failures whose inverse cubic
has one rational root and an irreducible quadratic factor over $\Q$.
What is their target-height exponent and leading constant?
\end{question}
The quadratic pair can supply local points at primes where the rational
branch is nonintegral. The split denominator-cover argument alone does
not classify this phenomenon. A useful first step is to express the
quadratic discriminant and the reconstruction valuations in terms of
the rational root and the remaining coefficients.
\wtag{} This is the second half of Part~I's Research
question~\ref{alg:sub:q-split}, with a counting question added.
[Added 30 September 2026, batch~64: the classification is answered by
Part~III for every $c\ne0$ (Theorem~\ref{alg:ns:thm:main} and
Corollary~\ref{alg:ns:cor:nonsplit}); the count for $c=\pm2$ in square
boxes is $\frac{27}{2\pi^2}T\log T+O(T)$
(Theorem~\ref{alg:ns:thm:count-main}). The exponent and constant for
other $c$ remain open (Research question~\ref{alg:ns:q:slices}).]

\begin{question}[A quantitative remainder]\label{alg:dh:q:remainder}
Can one prove an explicit error term
$\#\Hsplit_c(T)=\kappa(c)T^{2/3}+O_c(T^{2/3-\eta})$ with a specified
$\eta>0$, or identify a secondary term of order $T^{1/2}$?
\end{question}
The proof here deliberately separates the cutoff limits and gives an
$o(T^{2/3})$ remainder. A secondary term must account for the quadratic
height cutoff, affine shifts of the three cusp lines, lattice effects,
and congruence conditions simultaneously. The finite count table suggests
that these effects are numerically significant, but it does not identify
a secondary constant.

\begin{question}[Uniformity in the third coordinate]\label{alg:dh:q:uniform}
How far can the asymptotic be made uniform for $0<|c|\le T^\theta$?
What is the asymptotic when all three target coordinates are bounded
by a common height?
\end{question}
The factors $\sigma(v_2(c))$ and $3/p^{2v_p(c)}$ are explicit, but their
moduli grow with $c$. The fixed-$c$ sieve and cusp bounds cannot simply
be summed without uniform error control. The dyadic maximum near
$v_2(c)=3$ is an additional feature worth understanding in a global count.

\begin{question}[Coefficient-only local certificates]\label{alg:dh:q:coefficient}
Can the complete split local criterion be expressed directly in
$(A,B,c)$ by a short collection of valuations and congruences, without
first factoring the inverse cubic?
\end{question}
This paper supplies an exhaustive root-numerator parametrization, not a
factorization-free coefficient criterion. The modulo-eight target image
is finite, but the odd-prime root-collision condition also interacts
with the splitting assumption. A compact certificate format could be
more useful computationally than a single large formula.
\wtag{} This extends Part~I's Research question~\ref{alg:sub:q-dyadic},
which asks the same for the dyadic certificate alone.
[Added 30 September 2026, batch~64: answered by Part~III. The integral
roots of a completely split target's monic inverse polynomial are its
root numerators, so by Theorem~\ref{alg:ns:thm:main} such a target is
locally integral everywhere exactly when $\gcd(3Bc-4,27Ac^2-4)$ is a
power of two and the modulo-eight table $E(A,B,c)$ holds; no
factorization is needed (Section~\ref{alg:ns:sec:answers}). Applied to
completely split targets, this criterion reproduces all fifteen counts
of the table in Section~\ref{alg:dh:sec:count}
(Section~\ref{alg:ns:sec:consistency}).]

\begin{question}[Number fields and $S$-integer counting]\label{alg:dh:q:numberfield}
What replaces the constant $\kappa(c)$ over the integers of a number
field, and what is the exact counting constant after inverting a finite
set of primes?
\end{question}
The numerator-root method and denominator covering have natural ideal
versions. Archimedean regions, ideal classes, residue characteristic two,
and the unit group introduce new issues. The existence theorem for
$S$-integers proved here does not by itself give their height asymptotics.
[Added 30 September 2026, batch~64: Part~III's
Theorem~\ref{alg:ns:thm:numberfields} gives a complete local membership
certificate over every ring $\Oh_{K,S}$, for all fibers; the counting
constants asked for here remain open (Research
question~\ref{alg:ns:q:numberfield} asks the same for nonsplit fibers).]

\begin{question}[The finite integral exceptional set]\label{alg:dh:q:exceptions}
Can the divisor enumeration of $\Isplit_c$ be simplified to a closed
arithmetic classification of those $c$ for which it is empty, together
with an exact formula for its cardinality?
\end{question}
The proof reduces the problem to finitely many signed factorizations of
divisors of $2c$, congruences modulo three, and explicit integrality
conditions. This is a bounded arithmetic problem with no remaining
geometric elimination. A classification would refine, rather than
replace, the current effective procedure.

\begin{question}[Certified finite-fiber arithmetic]\label{alg:dh:q:library}
Can the inverse-certificate, local-selection, and finite-exception
arguments be packaged as a reusable Lean/Rocq library for polynomial
maps with small rational fibers?
\end{question}
The natural abstraction is a proved bijection between a polynomial
fiber and simple roots, followed by exact denominator data. This could
support both positive congruence witnesses and globally negative
integrality certificates without invoking an untrusted solver.
\wtag{} Cf.\ Part~I's Research question~\ref{alg:sub:q-library}, which
asks essentially the same question (complete inverse certificates,
denominator covering and the inverse-ball theorem as a generic
workflow).

\section{Conclusion of Part~II}\label{alg:dh:sec:conclusion}
The obstruction exhibited here is not geometric nonexistence: every
counted fiber has three rational points. It is the impossibility of
choosing one branch that is integral at all primes, even though an
integral branch exists separately at each prime.

For the fixed ProveIt map, this obstruction is Zariski dense in every
nonzero integer slice and hence in the full target space. In the split
sector it admits a complete arithmetic description, an exact dyadic
factor, a finite list of globally integral exceptions, and a sharp
$T^{2/3}$ target-height law. The resulting theorem resolves the concrete
full-dimensional question left in the repository report while keeping
the unsplit arithmetic sector, uniform counting, and formal verification
as explicitly separate next problems.
[Added 30 September 2026, batch~64: the unsplit sector is treated in
Part~III, which classifies it on every slice and shows that on $C=\pm2$ it
carries the dominant $T\log T$ term of the integral Hasse failures.]

\section{A compact proof and certificate ledger}\label{alg:dh:sec:ledger}
\wtag{} This is the manuscript's Appendix~A.

\begin{longtable}{>{\raggedright\arraybackslash}p{0.26\textwidth}>{\raggedright\arraybackslash}p{0.41\textwidth}>{\raggedright\arraybackslash}p{0.24\textwidth}}
\toprule
Claim & Main proof ingredient & Computational check\\
\midrule
\endhead
Complete inverse and determinant & Rational chart, chain rule, localization & Symbolic identities\\
Odd-prime integrality & Two divisible roots, valuation comparison & Finite residue checks\\
Odd-prime local obstruction & All three residues coincide & Exact rational comparison\\
Dyadic factor & Root parities and cancellation \eqref{alg:dh:eq:dyadic-cancellation} & Independent mod-eight image and bad-prime densities\\
Dense family & Explicit inverse first coordinates, gcd one, nonzero Jacobian & 108 exact signed-slice tests\\
$S$-integer density & Two assigned denominator primes, CRT & Proof, not claimed as a finite exhaustive test\\
Finite integral exceptions & $D_a\mid2c$ and signed factorizations & Complete divisor enumeration\\
Numerator density & CRT and \eqref{alg:dh:eq:large-prime-tail} & No numerical inference used\\
Sharp target asymptotic & Cusp bound, beta area, six-to-one map & Exact finite counts for comparison\\
\bottomrule
\end{longtable}

\Needspace{15\baselineskip}
\section{Reproduction commands of Part~II}\label{alg:dh:sec:repro}
\wtag{} This is the manuscript's Appendix~B.

From the extracted archive directory, install the optional verification
packages and run:
\begin{lstlisting}
python -m pip install -r requirements.txt
python code/verify_results.py
python code/split_fibres.py 1 1 3 --modulus 18144000
python code/count_heights.py
pdflatex -interaction=nonstopmode -halt-on-error article.tex
pdflatex -interaction=nonstopmode -halt-on-error article.tex
\end{lstlisting}
The exact-fraction utility \rpath{split_fibres.py} needs only the Python
standard library. The other packages are needed for symbolic verification,
fast residue enumeration, and numerical display of constants.
The retained outputs are under \rpath{data/}. This archive does not
modify the ProveIt repository.

\wtag{} In this report the files carry the prefix
\texttt{02-dense-hasse-}: the programs are under \rpath{code/}, the
\texttt{requirements.txt} is \rpath{data/02-dense-hasse-requirements.txt},
and the manuscript's \texttt{Makefile} is
\rpath{code/02-dense-hasse-Makefile}, which still names the delivered
paths. The manuscript is printed here as Part~II; its \texttt{article.pdf}
is not shipped. The commands above are the delivery's:
\texttt{count\_heights.py} imports \texttt{split\_fibres} by its
delivered name, and the scripts write under delivered names into
\rpath{../data/}, so they must be run on a copy
(Appendix~\ref{alg:dh:app:rerun}). On Windows use \texttt{py} or
\texttt{uv run} for \texttt{python}.

% ==== Part III begins
%% ====================================================================
%% Batch-64 addition (manuscript 02, "Beyond Split Fibers: Exact Local
%% Certificates and a T log T Law for Integral Hasse Failures"). All
%% labels of Part III carry the prefix alg:ns: (nonsplit). Section 28 is
%% new ([write]); manuscript Section k is printed as Section k+28
%% (k = 1..14; Section 41 is its research questions, Section 42 its
%% Conclusion) and its Appendices A--D as Sections 43--46.
%% ====================================================================
\clearpage
\phantomsection\part*{Part III.\quad Nonsplit fibers: a complete local certificate and a $T\log T$ law}
\addcontentsline{toc}{part}{Part III.\texorpdfstring{\quad}{ } Nonsplit fibers: a complete local certificate and a \texorpdfstring{$T\log T$}{T log T} law}
\fancyhead[L]{\small Part III: nonsplit fibers and a $T\log T$ law}

\noindent\emph{Part~III was added on 30~September 2026 in batch~64 of
ProveIt's incoming-report intake, from one manuscript of 30~September
2026 written against Parts~I and~II as pinned. Parts~I and~II are
unchanged except for dated notes and the title-page spacing (listed in
Appendix~\ref{alg:ns:app:provenance}). Part~III answers Research
question~\ref{alg:sub:q-split} on every slice $C\ne0$, and Research
question~\ref{alg:sub:q-height} for square boxes in all sectors: on
$C=2$ the integral Hasse failures number
$\frac{27}{2\pi^2}T\log T+O(T)$, almost all of them nonsplit, so the
completely split sector that Part~II counts ($T^{2/3}$) is not the
dominant one. Section~\ref{alg:ns:sec:relation} states exactly what is
answered, fixes the notation against Parts~I and~II and lists what is
printed twice; Sections~\ref{alg:ns:sec:intro}--\ref{alg:ns:app:dyadic}
are the manuscript.}

\section{Relation to Parts~I and~II}\label{alg:ns:sec:relation}
\wtag{} This section was written when the manuscript was added to the
report.

\subsection{The manuscript}\label{alg:ns:sec:source}
Part~III is batch-64 manuscript~02, \emph{Beyond Split Fibers: Exact
Local Certificates and a $T\log T$ Law for Integral Hasse Failures},
subtitled \emph{Arithmetic of ProveIt's three-variable Keller map} and
dated 30~September 2026. It was delivered as the archive
\texttt{ProveIt\_Beyond\_Split\_Fibers} and is pinned to ProveIt commit
\nssha; its provenance, shipped files and editorial changes are in
Appendix~\ref{alg:ns:app:provenance}. Its title-page abstract and
evidence note read as follows.
\begin{quote}\small
We resolve the rational-plus-quadratic local-solubility question left open
in ProveIt's arithmetic-fiber report. For its explicit three-variable
polynomial map $F$ with Jacobian determinant $-2$, an integral target
$(A,B,C)$ is locally integral at every prime precisely when a monic inverse
cubic has an integer root, the integer
$\gcd(3BC-4,27AC^2-4)$ has no odd prime factor, and an explicit table of
congruences modulo $8$ holds. Thus no separate analysis of the quadratic
factor's splitting fields is needed once a rational root is known.

On the slice $C=2$, this criterion becomes a primitive parametrization:
$A=r^3-r^2+rb$, $B=2b$, and $\gcd(3r-1,3b-1)=1$.
We prove that nonsplit integral Hasse failures in $|A|,|B|\le T$ number
$\frac{27}{2\pi^2}T\log T+O(T)$; the same formula counts all integral
Hasse failures on the slice. This differs fundamentally from the
$T^{2/3}$ scale previously obtained for completely split fibers.
We also determine the entire integral image on $C=2$: it has two
polynomial parametrizations with disjoint integer images, containing
$2^{5/3}T^{1/3}+O(1)$ targets in the same box. A number-field and
$S$-integer extension, exact algorithms, finite certificates, and further
research questions complete the study.

\textbf{Evidence and novelty boundary.}
All principal claims have proofs below. The computations are reproducible
checks and finite certificates, not Lean or Rocq formalizations. The
classification and counting theorems are candidate new contributions
relative to the pinned repository and primary sources reviewed here;
worldwide priority has not been established. The map, inverse-cubic
geometry, and earlier split-fiber results are explicitly attributed.
This article does not claim a new disproof of the Jacobian conjecture.
\end{quote}

\subsection{What Part~III answers, and what it does not}\label{alg:ns:sec:answers}
\begin{itemize}[itemsep=2pt,topsep=3pt]
\item \emph{Research question~\ref{alg:sub:q-split} is answered, on every
slice $C\ne0$ and in both sectors.} Theorem~\ref{alg:ns:thm:main} is a
complete local certificate for every integral target: an integer root of
the monic inverse polynomial $H$, the gcd $G$ of $U=3BC-4$ and
$V=27AC^2-4$ a power of two, and the modulus-eight table $E$.
Corollary~\ref{alg:ns:cor:nonsplit} turns it into a complete Hasse test
for fibers with one rational point and an irreducible quadratic pair,
the second sector of the question, which Part~II left open. (Part~II's
Theorem~\ref{alg:dh:thm:local-complete} had classified the completely
split sector in terms of root numerators.) On $C=0$ every integral target
has an integral point.
\item \emph{Research question~\ref{alg:sub:q-height} is answered for square
boxes, in all sectors.} On $C=2$ (and, by reflection, $C=-2$), in
$|A|,|B|\le T$, the locally integral targets, all integral Hasse failures
and the nonsplit ones each number $\frac{27}{2\pi^2}T\log T+O(T)$
(Theorem~\ref{alg:ns:thm:count-main}), and the integral targets
$2^{5/3}T^{1/3}+O(1)$ (Corollary~\ref{alg:ns:cor:integral-count}). The
completely split sector counted by Part~II,
$\#\Hsplit_2(T)\sim\kappa(2)T^{2/3}$
(Theorem~\ref{alg:dh:thm:headline-count}), is therefore \emph{not} the
dominant term: almost every integral Hasse failure on $C=2$ is nonsplit
(Corollary~\ref{alg:ns:cor:density}). The constant $27/(2\pi^2)$ is twice
the Euler product $\prod_{p\ne3}(1-p^{-2})=27/(4\pi^2)$, the main term
with an explicit Euler product that the discussion of the question asks
for. Not answered: unequal box sides $H_A\ne H_B$ (Research
question~\ref{alg:ns:q:unequal}); a second-order term or a power saving
(Research questions~\ref{alg:ns:q:second} and~\ref{alg:ns:q:power}); other
slices (Research question~\ref{alg:ns:q:slices}).
\item \emph{Part~II's Research question~\ref{alg:dh:q:coefficient} is
answered.} For a completely split integral target the numerators of the
roots are integer roots of $H$ (Section~\ref{alg:ns:sec:consistency}), so
the root gate holds automatically and local solubility everywhere is
``$G$ is a power of two and $E(A,B,C)$ holds'': a criterion in
$(A,B,c)$ that needs no factorization of the inverse cubic.
\item \emph{Part~II's Research question~\ref{alg:dh:q:quadratic} is
answered in its classification half on every slice, and in its counting
half on $C=\pm2$} (order $T\log T$, constant $27/(2\pi^2)$). The exponent
and constant on other slices remain open (Research
question~\ref{alg:ns:q:slices}).
\item \emph{Partly answered or related.} Research
question~\ref{alg:sub:q-dyadic}: membership in the dyadic image is
compressed to the five-row table $E$ in the target coefficients
(Proposition~\ref{alg:ns:prop:two}); the integral fiber \emph{sizes}
$j=1,2,3$ are not compressed, and the proof is still the exhaustive
$512$-case check, not a hand case proof. Research
question~\ref{alg:sub:q-extensions} and Part~II's Research
question~\ref{alg:dh:q:numberfield}: Theorem~\ref{alg:ns:thm:numberfields}
gives a membership certificate over every ring $\Oh_{K,S}$, with the
uniform dyadic modulus $8R$; fiber distributions, minimal moduli and
height counts over number fields remain open (Research
questions~\ref{alg:ns:q:dyadicprec} and~\ref{alg:ns:q:numberfield}).
Research question~\ref{alg:sub:q-library} and Part~II's Research
question~\ref{alg:dh:q:library} are related to Research
question~\ref{alg:ns:q:software}. Research questions 12.4 (answered by
Part~II), 12.6--12.8 and Part~II's Research
questions~\ref{alg:dh:q:remainder}, \ref{alg:dh:q:uniform}
and~\ref{alg:dh:q:exceptions} are untouched.
\end{itemize}
Part~II's results stand as proved: its split-sector theorems are not
superseded, and its $T^{2/3}$ law is the exact order of the split
sector, which Part~III bounds only by $O(T)$
(Lemma~\ref{alg:ns:lem:splitbound}). What changes is the scope of the
sentence ``Research question~12.3 is answered only in the completely split
sector'': dated notes in Parts~I and~II now point here.

\subsection{Consistency with Parts~I and~II}\label{alg:ns:sec:consistency}
Where the Parts meet they agree. Each item was checked at the write, by
hand or by an exact computation of its own (not the delivered programs).
\begin{itemize}[itemsep=2pt,topsep=3pt]
\item \emph{Part~II's numerators are the roots of $H$.} Since
$H(Ct)=C^2g(t)$, the roots of the monic $H$ are $C$ times the roots of
Part~I's $g_{A,B,C}$; on $C=c$ these are exactly Part~II's root
numerators $a,b,r$, with $a+b+r=2$ \eqref{alg:dh:eq:numerator-roots}.
\item \emph{Odd primes.} At an odd prime $p\nmid c$, Part~II's obstruction
\eqref{alg:dh:eq:gcd} is $a\equiv b\equiv r\equiv2\cdot3^{-1}$, that is,
$\overline H=(S-2/3)^3$, which is exactly $p\mid U$ and $p\mid V$ in the
proof of Theorem~\ref{alg:ns:thm:odd}.
\item \emph{The slice $C=2$.} Part~III's polynomial $f$ with $H(2T)=8f(T)$
is $g_{A,B,2}/2$, so its integer roots are Part~I's roots $a,b,c$ of
\eqref{alg:eq:planemap}; $d=f'(r)$ is Part~I's $D_r$, Part~I's
$d(t)=g'(t)$ is $2d$ there, and \eqref{alg:ns:eq:slice-inverse} is
Part~I's \eqref{alg:eq:planepoints}. For a completely split target,
Theorem~\ref{alg:ns:thm:param} reduces to Part~I's criterion
\eqref{alg:eq:alllocalcriterion}: by \eqref{alg:ns:eq:gcd-derivative} its
condition is $\gcd(3r-1,d)=1$, and with $u=r-s$, $v=r-t$ one has
$3r-1=u+v$ and $d=uv$; a prime divides both $u+v$ and $uv$ exactly when
it divides $u$ and $v$, so the condition is $\gcd(u,v)=1$, which is
Part~I's $g=1$.
\item \emph{The dyadic table.} The program of
Section~\ref{alg:ns:app:dyadic} was rerun at the write: the image of $F$
on $(\Z/8\Z)^3$ has $176$ classes and equals $E$ pointwise, and the
residue fiber histogram is $336,112,48,16$ for sizes $0,2,4,6$, which is
Part~I's \eqref{alg:eq:dyadicfinite} at $m=3$;
$176/512=11/32$ is \eqref{alg:eq:rho2}.
\item \emph{Constants.} Part~III's $\alpha=27/(4\pi^2)$ is Part~I's
parameter density \eqref{alg:eq:paramdensity}, from the same M\"obius sum
over $d$ prime to $3$; $27/(2\pi^2)=1.3678359792\ldots$ and
$2^{5/3}(10^5)^{1/3}=147.36\ldots$ (against $I_2(10^5)=148$).
\item \emph{The split counts.} By Lemma~\ref{alg:ns:lem:split-nointeger}
every locally integral completely split target on $C=2$ is a Hasse
failure, so the manuscript's $\Sigma(T)$ is Part~II's $\#\Hsplit_2(T)$.
Its recorded values $121$, $653$, $3\,304$ at $T=10^3,10^4,10^5$ (table
in Section~\ref{alg:ns:sec:computation}) are Part~II's recorded exact
counts for $c=2$ (\rpath{data/02-dense-hasse-height_counts.txt}).
Moreover, the coefficient criterion of Theorem~\ref{alg:ns:thm:main} with
the integrality test of Corollary~\ref{alg:ns:cor:integral-test}, applied
by an independent enumeration of completely split targets on $C=c$,
reproduces all fifteen counts of Part~II's table in
Section~\ref{alg:dh:sec:count} ($c=1,2,3,4,8$ and $T=10^3,10^4,10^5$).
\item \emph{Integral points.} Part~II's table of integral split exceptions
after Theorem~\ref{alg:dh:thm:finite} lists none on $C=2$, as
Lemma~\ref{alg:ns:lem:split-nointeger} requires; the integral image on
$C=2$ (Theorem~\ref{alg:ns:thm:curves}) consists of nonsplit fibers and
the repeated-root target $(0,0,2)$. The formulas $F(1/2,-2,4)=(0,4,2)$,
\eqref{alg:ns:eq:family-point} for $2\le b<40$,
\eqref{alg:ns:eq:integralpoint} for $|r|\le30$ and $F(1,0,0)=(0,0,2)$
were evaluated exactly.
\item \emph{The plane $C=0$}: \eqref{alg:ns:eq:boundary} is Part~I's
boundary point \eqref{alg:eq:boundarypoint} and the display in Part~II's
Theorem~\ref{alg:dh:thm:headline-density}.
\end{itemize}

\subsection{Notation: Parts~I and~II against Part~III}\label{alg:ns:sec:notation}
Part~III keeps the manuscript's letters, with one exception, so several
letters mean different things in the three Parts.
Table~\ref{alg:ns:tab:notation} fixes them; within Part~III only the
Part~III column applies, and the last column prints the tempting false
reading. Table~\ref{alg:dh:tab:notation} compares Parts~I and~II.

\begingroup\small
\setlength{\tabcolsep}{4pt}
\begin{longtable}{@{}>{\raggedright}p{0.09\textwidth}>{\raggedright}p{0.33\textwidth}>{\raggedright}p{0.36\textwidth}>{\raggedright\arraybackslash}p{0.15\textwidth}@{}}
\caption{Symbols of Part~III whose meaning differs in Parts~I and~II.}
\label{alg:ns:tab:notation}\\
\toprule
Symbol & Parts I and II & Part III & Not to be read as\\
\midrule
\endfirsthead
\toprule
Symbol & Parts I and II & Part III & Not to be read as\\
\midrule
\endhead
\bottomrule
\endfoot
$H$ & I: parameter heights ($M(H)$, $\mathcal A(H)$), box sides
$H_A,H_B$ in Research question~\ref{alg:sub:q-height}; II: $\Hsplit_c(T)$
is a set & the monic inverse polynomial $H(S)$ \eqref{alg:ns:eq:HUV};
the counts $H_2(T)$, $H_2^{\NS}(T)$, $H_2^{\spc}(T)$ of Hasse failures on
$C=2$ & a height\\
$G$ & I: the binary cubic $G_{A,B,C}(T,S)$ \eqref{alg:eq:binarycubic} &
the positive integer $\gcd(|U|,|V|)$ \eqref{alg:ns:eq:HUV} & a cubic\\
$U,V$ & I: $u=1+xy$ (lower case); II: an indeterminate $U$ in the proof
of Lemma~\ref{alg:dh:lem:odd-integrality}, and auxiliary integers $U,V$
in the proof of Theorem~\ref{alg:dh:thm:S} & $U=3BC-4$, $V=27AC^2-4$ & \\
$E$ & not used & the dyadic condition $E(A,B,C)$ (table before
Theorem~\ref{alg:ns:thm:main}) & \\
$S$ & a finite set of primes; I: also the second projective variable in
\eqref{alg:eq:binarycubic}; II: a sum, printed $\Sigma$ & the
indeterminate of $H(S)$; a finite set of finite places in
Section~\ref{alg:ns:sec:numberfields}; the manuscript's split count
$S(T)$ is printed $\Sigma(T)$ & a count\\
$\Sigma$ & II: a sum of two half-numerators or roots & $\Sigma(T)$, the
locally integral completely split targets on $C=2$ in the box
($=\#\Hsplit_2(T)$) & a sum\\
$a,b,r$ & I: integer roots $a,b,c$ on $C=2$, $r$ an index; II: root
numerators with $a+b+r=2$ & $a$: an integer root of $H$ (a root
numerator, $a=Ct$); $b=B/2$ and $r$ an integer root of $f$ on $C=2$
\eqref{alg:ns:eq:param}; $\target$ is the target & roots of $g$\\
$C$, $c$ & I: $C$ the third coordinate, $c$ a root; II: $c$ the third
coordinate & $C$ the third coordinate (the questions write $c$ for a
fixed slice); $c_0$ an implied constant in Lemma~\ref{alg:ns:lem:range} &
a root\\
$d,D$ & I: $d(t)=g'(t)$, $D_r$ products of root differences; II: $d=s-t$,
$D=dv$, $D_r$ for numerators & $D=H'(a)$ \eqref{alg:ns:eq:D};
$d=f'(r)$ \eqref{alg:ns:eq:dDelta}, so $D=4d$ and $d$ is Part~I's $D_r$
on $C=2$; $d$ also a divisor in Lemmas~\ref{alg:ns:lem:primitive}
and~\ref{alg:ns:lem:mean} & Part~I's $d(t)$\\
$\Delta$ & I: the discriminant of the inverse cubic & $\Delta_a$
\eqref{alg:ns:eq:delta-H} and $\Delta$ \eqref{alg:ns:eq:dDelta},
discriminants of the \emph{quadratic} factor, $\Delta_{2r}=4\Delta$ on
$C=2$ & the cubic's\\
$f,g$ & I: $g$ the inverse cubic; II: also valuations & $g$ the same
(proof of Proposition~\ref{alg:ns:prop:inverse}); $f(T)=H(2T)/8=g(T)/2$
on $C=2$ \eqref{alg:ns:eq:f} & \\
$T$ & I: the indeterminate of $g$; II: also the height bound & the height
bound; also the indeterminate of $f$ and $g$ & \\
$t$ & I: $t=y+1/x$; II: an integer root & $t=y+1/x$; a root $t$ in
Lemmas~\ref{alg:ns:lem:split-nointeger} and~\ref{alg:ns:lem:splitbound};
the parameter of the cusp curve & \\
$P,Q,R$ & the coordinates of $F$; II: also the cutoff $R$ &
the coordinates; $P(T)$ the marked-root count
(Section~\ref{alg:ns:sec:count}); $R=\lfloor\sqrt{T/8}\rfloor$ in
Lemma~\ref{alg:ns:lem:range}; $R=\Oh_{K,S}$ in
Section~\ref{alg:ns:sec:numberfields} & \\
$L,I,N$ & I: $L=2H+1$, $I(q)$ \eqref{alg:eq:Iq}, $N$ a number of primes;
II: $L=T^{1/3}$, the beta area $I$, $\Isplit_c$ & $L_2(T)$, $I_2(T)$,
$N(T)$ target counts on $C=2$; $L$ an interval length in
Lemma~\ref{alg:ns:lem:primitive}; $\mathcal I_r(T)$ an interval & \\
$\alpha$ & I: $\alpha_p$ in Theorem~\ref{alg:thm:thinsieve}, automorphisms
$\alpha,\beta$; II: roots $\alpha,\beta,\gamma$ & $\alpha=27/(4\pi^2)$
\eqref{alg:ns:eq:mean}, Part~I's density \eqref{alg:eq:paramdensity} & \\
$e$ & I: $e=v_p(\det J)$ in Lemma~\ref{alg:lem:ball}; II: $e=v_2(|c|)$ &
$e=v(2)$, the valuation of $\det J$ up to a unit, as in Part~I &
Part~II's $e$\\
$q,\tau,\delta$ & I: $q$ a field size, $\tau$ the project's family
parameter; II: $q$ a prime, $\tau$ a divisor count, $\delta(c)$ a density
& $q=|r|$; $\tau$ the divisor function, as in Part~II; $\delta$ a Newton
correction in Lemma~\ref{alg:ns:lem:lift} & \\
\end{longtable}
\endgroup

The one rename: the manuscript's split count $S(T)$ (Sections
\ref{alg:ns:sec:mult}, \ref{alg:ns:sec:computation}
and~\ref{alg:ns:app:small}) is printed $\Sigma(T)$, as Part~II printed its
sum $S$ as $\Sigma$, because $S$ is also the indeterminate of $H$ and the
finite set of places. The manuscript's implied constant $C$ in the proof
of Lemma~\ref{alg:ns:lem:range} is printed $c_0$, because $C$ is the third
coordinate. No normalization was changed.

\subsection{What is printed twice, and why}\label{alg:ns:sec:duplicates}
The manuscript restates parts of Parts~I and~II to be self-contained.
Its later sections cite these restatements, and all of its labels are
kept, so they are printed where they stand; the manuscript attributes
them (its attribution table in Section~\ref{alg:ns:sec:intro}).
\begin{itemize}[itemsep=2pt,topsep=3pt]
\item The map \eqref{alg:ns:eq:F} is \eqref{alg:eq:F}, and
\eqref{alg:ns:eq:boundary} is \eqref{alg:eq:boundarypoint}.
\item Proposition~\ref{alg:ns:prop:inverse} is Part~I's
Theorem~\ref{alg:thm:root} (and Part~II's
Proposition~\ref{alg:dh:prop:inverse}) in the monic coordinate $a=Ct$:
Part~I's $w=g'(t)/2=1/x$ is $D/(2C)$. The proof here is a direct
substitution.
\item Lemma~\ref{alg:ns:lem:rootgate} is the Chebotarev mechanism of
Part~I's Lemma~\ref{alg:lem:cubicHasse} and
Theorem~\ref{alg:thm:rationalHasse}, applied to integral local points; the
manuscript lists it as inherited.
\item Lemma~\ref{alg:ns:lem:hensel-odd} at odd primes and
Lemma~\ref{alg:ns:lem:lift} at $p=2$ are cases of Part~I's inverse-ball
Lemma~\ref{alg:lem:ball} (with Corollary~\ref{alg:cor:precision}), and
of Part~II's lifting fact \eqref{alg:dh:eq:pdivc}.
Lemma~\ref{alg:ns:lem:lift} is stated over any complete discrete
valuation ring of characteristic zero, which
Theorem~\ref{alg:ns:thm:numberfields} uses; it is kept as a second
route, a Newton iteration.
\item Modulus-eight sufficiency and the measure $11/32$ are Part~I's
Theorem~\ref{alg:thm:dyadic}; the residue histogram in the Remark after
Proposition~\ref{alg:ns:prop:two} is \eqref{alg:eq:dyadicfinite}
(pointer added there). The table $E$ itself is new.
\item Lemma~\ref{alg:ns:lem:split-nointeger} is Part~I's
Theorem~\ref{alg:thm:plane} (``It has no integral point''), with the
same proof; the manuscript does not cite it, and a pointer was added
there. \eqref{alg:ns:eq:slice-inverse} is Part~I's
\eqref{alg:eq:planepoints} for split fibers.
\item \eqref{alg:ns:eq:prior-split} is Part~II's
Theorem~\ref{alg:dh:thm:headline-count} at $c=2$; the manuscript does not
use it, and Lemma~\ref{alg:ns:lem:splitbound} is a weaker bound proved
independently.
\item New in their generality: Theorems~\ref{alg:ns:thm:main},
\ref{alg:ns:thm:odd}, the table $E$ of
Proposition~\ref{alg:ns:prop:two}, Corollary~\ref{alg:ns:cor:nonsplit},
Theorem~\ref{alg:ns:thm:param} (which extends Part~I's criterion
\eqref{alg:eq:alllocalcriterion} from split targets to every locally
integral target on $C=2$), Corollary~\ref{alg:ns:cor:family},
Proposition~\ref{alg:ns:prop:repeated}, Theorem~\ref{alg:ns:thm:curves},
the counting results of Sections~\ref{alg:ns:sec:count}
and~\ref{alg:ns:sec:mult}, and Theorem~\ref{alg:ns:thm:numberfields},
which extends the method of Part~I's
Theorem~\ref{alg:thm:rationalHasse} from rational to $S$-integral points.
\end{itemize}

\subsection{The formal project and Part~III}\label{alg:ns:sec:formalproject}
The manuscript cites ProveIt's Lean file
\rpath{Algebra/JacobianConjecture/Lean/JacobianConjecture/Counterexample.lean}
at its pin \nssha; the file is identical at Part~II's pin and at the
placement commit, so it is cited here as \cite{dh:proveitformal}. The file
defines the map as \texttt{counterexample} (namespace
\texttt{LeanProofs.JacobianCounterexample}, lines 95--101) and proves
\texttt{jacobianDet\_counterexample}, $\det JF=-2$ over every commutative
ring (line 134), and the collision \texttt{collision} (line 165). Part~III
uses the map, checks the determinant symbolically, and relies on no formal
statement. \emph{None of Part~III's results is formalized.} The project
has no $p$-adic, Hasse-principle or counting material, and the placement
of this report beside its Lean and Rocq developments confers no formal
status. The manuscript's formalization plan
(Section~\ref{alg:ns:sec:formal}) is a proposal; it keeps Chebotarev as an
explicitly imported theorem.

\section{The question, the answer, and its scope}\label{alg:ns:sec:intro}

\subsection{A specific unresolved question in ProveIt}
The repository's report \emph{Arithmetic Local--Global Dichotomies for
ProveIt's Three-Variable Keller Map} develops split integral Hasse
failures and exact local fiber laws \wtag{}~(Parts~I and~II of this
report). Its Part~II classifies
completely split fibers and counts them in fixed slices. It explicitly
leaves open the case of one rational point accompanied by an irreducible
quadratic pair, and the corresponding target-height asymptotics. Part~I,
Section~\ref{alg:sub:q-split}, asks for local integral classification beyond split fibers;
Section~\ref{alg:sub:q-height} asks for target-height asymptotics. The appended Part~II
records the nonsplit case as still open and asks for coefficient-level
certificates \wtag{}~(Research questions~\ref{alg:dh:q:quadratic}
and~\ref{alg:dh:q:coefficient}).

The difficulty is genuine: a rational point can fail integrality at a
prime where the two nonrational geometric points become locally rational
and provide integral alternatives. Testing only the denominator of the
rational point therefore gives the wrong local answer. Nevertheless,
for this particular map the entire odd-prime phenomenon collapses to
two coefficient congruences.

We use the snapshot \nssha. The report is located at
\begin{quote}\small\raggedright
\rpath{SetTheory/Cardinals/docs/reports/jacobian-conjecture/}\allowbreak
\rpath{arithmetic-local-global-fibers/article.tex}.
\end{quote}
\wtag{} That is the present report, whose Parts~I and~II are unchanged
from this pin to the placement commit \texttt{3025c15df}; the pinned
commit is abbreviated \texttt{6d04e1e38} elsewhere in this report.
Its mathematical assertions are treated as research claims rather than
as formal consequences of residing in a repository. The proofs of our
main theorems are independent of its unformalized counting arguments.

\subsection{The map and the terminology}
Throughout, let
\begin{equation}\label{alg:ns:eq:F}
\begin{gathered}
u=1+xy,\qquad h=u^2z+y^2(1+3u),\\
F(x,y,z)=(uh,\ y+3xh,\ x(5-3u-x^2z)).
\end{gathered}
\end{equation}
Thus its coordinate polynomials are
\begin{align*}
P&=(1+xy)^3z+y^2(1+xy)(4+3xy),\\
Q&=y+3x(1+xy)^2z+3xy^2(4+3xy),\\
R&=2x-3x^2y-x^3z.
\end{align*}
The identity $\det JF=-2$ and explicit collisions are formalized in the
repository \cite{dh:proveitformal} \wtag{}~(declarations
\texttt{jacobianDet\_counterexample} and \texttt{collision};
Section~\ref{alg:ns:sec:formalproject}). The map is attributed there and by Gao to
Alp\"oge; Gao establishes its complex fiber geometry \cite{gao}.

\begin{definition}
An integral target $\target=(A,B,C)\in\Z^3$ is \emph{everywhere locally
integral} if $\target\in F(\Z_p^3)$ for every rational prime $p$.
It is an \emph{integral Hasse failure} if it is everywhere locally integral
but does not belong to $F(\Z^3)$. A fiber is \emph{nonsplit} in this article
if its inverse cubic factors over $\Q$ as a linear factor times an
irreducible quadratic. A \emph{completely split} fiber has three distinct
rational points unless repeated roots are explicitly mentioned.
\end{definition}
Every everywhere locally integral target studied here has a rational
preimage, and hence a real preimage. Adding an archimedean-solubility
condition would not change the definition. An integral Hasse failure is
not a rational Hasse failure.

\subsection{Main classification}
Put
\begin{equation}\label{alg:ns:eq:HUV}
\begin{aligned}
H(S)&=S^3-2S^2+BCS-2AC^2,\\
U&=3BC-4,\qquad V=27AC^2-4,\qquad G=\gcd(|U|,|V|).
\end{aligned}
\end{equation}
For integer coefficients $U\ne0$, so $G$ is positive. Let $E(A,B,C)$
be the following condition. All congruences in the first column are
modulo $8$.
\begin{center}
\begin{tabular}{@{}cl@{}}
\toprule
$C$ & Condition $E(A,B,C)$\\
\midrule
$0$ & no condition on $A,B$\\
$4$ & $B$ is even\\
$2$ or $6$ & $A$ and $B$ are even\\
odd & $[A\text{ even and }B\equiv2\pmod4]$\\
 & \quad or $[A\text{ odd and }BC\equiv2A+1\pmod8]$\\
\bottomrule
\end{tabular}
\end{center}

\begin{theorem}[Complete local certificate]\label{alg:ns:thm:main}
For every $(A,B,C)\in\Z^3$, the following are equivalent:
\begin{enumerate}[label=(\roman*)]
\item $(A,B,C)\in F(\Z_p^3)$ for every prime $p$;
\item $H$ has an integer root, $G$ is a power of $2$ (including $1$),
and $E(A,B,C)$ holds.
\end{enumerate}
If $C=0$, these conditions always hold and the target already has an
integral preimage. If $C\ne0$, global integral solubility is decided by
the finitely many simple integer roots of $H$ and the reconstruction
formulas in Section~\ref{alg:ns:sec:inverse}.
\end{theorem}
The rational-root condition cannot be dropped. The theorem is a finite
coefficient certificate \emph{together with} an ordinary rational-root
test; it is not a bounded congruence characterization of rational
solubility.

\subsection{Main counting and integral-image results}
Let $L_2(T)$ count everywhere locally integral targets $(A,B,2)$ with
$|A|,|B|\le T$. Let $I_2(T)$ count their globally integral subset, and
let $H_2(T)$ and $H_2^{\NS}(T)$ count all and nonsplit integral Hasse
failures, respectively. Targets, not roots or ordered parameter tuples,
are counted.

\begin{theorem}[The dominant nonsplit law]\label{alg:ns:thm:count-main}
As $T\to\infty$,
\begin{align}
L_2(T)&=\frac{27}{2\pi^2}T\log T+O(T),\label{alg:ns:eq:L-main}\\
H_2(T)&=\frac{27}{2\pi^2}T\log T+O(T),\label{alg:ns:eq:H-main}\\
H_2^{\NS}(T)&=\frac{27}{2\pi^2}T\log T+O(T).\label{alg:ns:eq:NS-main}
\end{align}
The entire integral image on $C=2$ is parametrized without duplication by
\begin{equation}\label{alg:ns:eq:curves-main}
(A,B)=(-2r^3+r^2+\varepsilon r,\ -6r^2+4r+2\varepsilon),
\qquad r\in\Z,\quad\varepsilon\in\{-1,1\}.
\end{equation}
Consequently
\begin{equation}\label{alg:ns:eq:I-main}
I_2(T)=2^{5/3}T^{1/3}+O(1).
\end{equation}
The same conclusions, after reflecting the second coordinate, hold on
$C=-2$.
\end{theorem}
The previous split theorem \wtag{}~(Part~II,
Theorem~\ref{alg:dh:thm:headline-count} at $c=2$) gives
\begin{equation}\label{alg:ns:eq:prior-split}
H_2^{\spc}(T)\sim
\frac{27\Gamma(1/3)^2}{8\pi^2\Gamma(2/3)}T^{2/3}.
\end{equation}
[\wtag{} In Part~II's notation $H_2^{\spc}(T)=\#\Hsplit_2(T)$, and the
constant is \eqref{alg:dh:eq:c2constant}; the manuscript's
$\Sigma(T)$ below counts the same targets
(Section~\ref{alg:ns:sec:consistency}).]
We do not claim that result anew and do not need its proof to establish
\eqref{alg:ns:eq:L-main}--\eqref{alg:ns:eq:NS-main}: an elementary $O(T)$ bound for
split fibers suffices. The new main term comes from the rational-plus-
quadratic sector, not from refining the split-family constant.

\subsection{Attribution and what is not claimed}
\begin{center}\small
\begin{tabular}{@{}p{.31\textwidth}p{.29\textwidth}p{.34\textwidth}@{}}
\toprule
Ingredient or result & Status here & Scope\\
\midrule
Map, determinant, collisions & Prior; independently checked & ProveIt formal source; Gao\\
Inverse cubic and missing cusp & Prior geometry; rederived & All relevant odd characteristics\\
Rational Hasse argument & Inherited mechanism & Chebotarev plus cubic Galois groups\\
Modulus-$8$ sufficiency and image size & Prior dyadic phenomenon & Explicit coefficient table reverified\\
$G$-criterion for all rooted cubics & Main refinement & Includes the previously open nonsplit case\\
Primitive parametrization on $C=2$ & Main refinement & All locally integral targets\\
$T\log T$ and two integral curves & Main counting results & Square target boxes on $C=\pm2$\\
General number-field certificate & Extension & Finite local tests outside a finite set\\
\bottomrule
\end{tabular}
\end{center}
We make no claim of a sharp second term in \eqref{alg:ns:eq:H-main}, a uniform
asymptotic for every varying nonzero $C$, or a new solution of the plane
Jacobian problem. The limited literature search did not locate these
specific arithmetic and counting statements in the inspected primary
sources; that finding is not a global priority proof.

\section{Inverse cubics and exact reconstruction}\label{alg:ns:sec:inverse}

\subsection{The boundary \texorpdfstring{$C=0$}{C=0}}
Direct substitution gives
\begin{equation}\label{alg:ns:eq:boundary}
F(0,B,A-4B^2)=(A,B,0).
\end{equation}
Thus no integral target on $C=0$ is an integral Hasse failure. Formula
\eqref{alg:ns:eq:boundary} is also the reason primes dividing $C$ will cause no
odd-prime obstruction.

\subsection{The affine inverse chart}
\begin{proposition}[Inverse chart]\label{alg:ns:prop:inverse}
Let $k$ be a field of characteristic different from $2$, and let
$A,B,C\in k$ with $C\ne0$. The points of $F^{-1}(A,B,C)(k)$ are in
bijection with the simple roots $a\in k$ of $H$. If
\begin{equation}\label{alg:ns:eq:D}
D=H'(a)=3a^2-4a+BC\ne0,
\end{equation}
then the corresponding point is
\begin{equation}\label{alg:ns:eq:reconstruct}
\boxed{\quad
x=\frac{2C}{D},\qquad
 y=\frac{2a-D}{2C},\qquad
 z=\frac{D(10D-12a-D^2)}{8C^2}.
\quad}
\end{equation}
For a point in the fiber, the inverse coordinate is $a=C(y+1/x)$.
\end{proposition}
\begin{proof}
Since $C=x(2-3xy-x^2z)\ne0$, we have $x\ne0$. Put
$t=y+1/x$ and $g(T)=CT^3-2T^2+BT-2A$.
Substitution of \eqref{alg:ns:eq:F}, followed by collecting terms, gives
\begin{equation}\label{alg:ns:eq:inverse-identities}
g(t)=0,\qquad g'(t)=\frac2x.
\end{equation}
These identities can also be derived by eliminating $z$ from the third
coordinate, using $z=2/x^2-3y/x-C/x^3$, and then substituting in the first
two coordinates. As $H(Ct)=C^2g(t)$, differentiation yields
$H'(Ct)=Cg'(t)=2C/x\ne0$.

Conversely, if $H(a)=0$ and $D\ne0$, set $t=a/C$ and $x=2C/D$.
Then $y=t-1/x=(2a-D)/(2C)$. Solving the third coordinate for $z$ gives
exactly \eqref{alg:ns:eq:reconstruct}. Substitution in the first two coordinates
reduces their residuals to multiples of $H(a)$, so they equal $A,B$.
All operations divide only by $2$, $C$, and $D$, which are nonzero.
The recovered value $C(y+1/x)=a$ proves injectivity of the construction.
\end{proof}
The same root chart appears in the prior arithmetic report, expressed
with $t$ rather than the monic coordinate $a$ \wtag{}~(Part~I,
Theorem~\ref{alg:thm:root}; Part~II,
Proposition~\ref{alg:dh:prop:inverse}). Gao's degree-
three and missing-fiber geometry is consistent with this description
\cite{gao}. We use the monic version because an integral coefficient
polynomial has integral rational roots.

\begin{corollary}[Exact integral test]\label{alg:ns:cor:integral-test}
For $C\ne0$ and a simple integer root $a$ of $H$, the reconstructed
point is integral if and only if
\begin{equation}\label{alg:ns:eq:divisibility}
D\mid2C,\qquad
2C\mid2a-D,\qquad
8C^2\mid D(10D-12a-D^2).
\end{equation}
An integral target belongs to $F(\Z^3)$ if and only if either $C=0$ or
at least one simple integer root satisfies \eqref{alg:ns:eq:divisibility}.
\end{corollary}
\begin{proof}
Each rational coordinate in \eqref{alg:ns:eq:reconstruct} is integral precisely
under the corresponding divisibility condition. Every rational root of
$H$ is an integer by the rational-root theorem, and
Proposition~\ref{alg:ns:prop:inverse} accounts for the entire rational fiber.
\end{proof}

\subsection{Why local solubility supplies a rational root}
\begin{lemma}[The cubic rational-root gate]\label{alg:ns:lem:rootgate}
If $(A,B,C)\in\Z^3$ is everywhere locally integral, then $H$ has an
integer root. For $C\ne0$, a rational root of $H$ guarantees that $H$
has a simple rational root.
\end{lemma}
\begin{proof}
When $C=0$, $H=S^2(S-2)$. Suppose $C\ne0$. If $H$ were irreducible over
$\Q$, its Galois group would be a transitive subgroup of $S_3$, hence
$C_3$ or $S_3$. Both contain a $3$-cycle. By the Chebotarev density
theorem \cite[Theorem~8.31 and Example~8.36]{milne}, there are infinitely
many good primes, excluding $2C$ and the discriminant, for which $H$
has irreducible reduction. At such a prime it has no root in $\Q_p$:
any $\Q_p$-root of the monic integral polynomial would be $p$-integral
and reduce to a root. Proposition~\ref{alg:ns:prop:inverse} rules out a local
preimage, a contradiction.

A reducible cubic has a rational root, which is integral here. If a
rational root is repeated but not triple, the remaining linear factor
has a simple rational root. The triple case would be
$H=(S-2/3)^3$, impossible for a monic integer polynomial. This proves
the last assertion.
\end{proof}
The argument uses qualitative Chebotarev, not an unproved hypothesis.
Its purpose is only to justify the rational-root gate. The subsequent
local tests and the entire counting proof are elementary.

\section{Every odd prime is controlled by one gcd}\label{alg:ns:sec:odd}

\subsection{From integral points to simple roots modulo a prime}
\begin{lemma}\label{alg:ns:lem:hensel-odd}
For every odd prime $p$,
\[
\target\in F(\Z_p^3)
\quad\Longleftrightarrow\quad
\overline{\target}\in F(\F_p^3).
\]
If $p\nmid C$, the right side holds precisely when $\overline H$ has a
simple root in $\F_p$. If $p\mid C$, it always holds.
\end{lemma}
\begin{proof}
Reduction gives necessity. At any residue preimage, $\det JF=-2$ is a
unit modulo $p$, so the multivariable Hensel argument lifts it to the
prescribed target over $\Z_p$. When $p\nmid C$, apply
Proposition~\ref{alg:ns:prop:inverse} over $\F_p$. When $p\mid C$, use
\eqref{alg:ns:eq:boundary} modulo $p$ and then lift.
\end{proof}

\begin{lemma}[The only obstruction for a rooted cubic]\label{alg:ns:lem:triple}
Let $q(S)$ be a monic cubic over a field and suppose it has a root in
that field. It has no simple root in the field if and only if
$q(S)=(S-a)^3$ for some field element $a$.
\end{lemma}
\begin{proof}
If a root is simple there is nothing to prove. Otherwise divide by its
square. The remaining factor is linear. Its root is simple unless it
coincides with the repeated root, which is exactly the triple case.
The converse is immediate in every characteristic.
\end{proof}

\subsection{The coefficient criterion}
\begin{theorem}[Odd-prime classification]\label{alg:ns:thm:odd}
Assume $H$ has an integer root. For each odd prime $p$,
\begin{equation}\label{alg:ns:eq:oddcriterion}
(A,B,C)\in F(\Z_p^3)
\quad\Longleftrightarrow\quad
p\nmid\gcd(3BC-4,27AC^2-4).
\end{equation}
\end{theorem}
\begin{proof}
The integer root reduces to a root modulo every prime. If $p\mid C$,
Lemma~\ref{alg:ns:lem:hensel-odd} gives local solubility, and
$3BC-4\equiv-4\not\equiv0\pmod p$ gives the right side.

Suppose $p\nmid C$. By Lemmas~\ref{alg:ns:lem:hensel-odd} and
\ref{alg:ns:lem:triple}, failure occurs exactly when $\overline H$ is a cube
of a linear factor. In characteristic $3$, a linear cube has no
$S^2$ term, whereas $H$ has coefficient $-2\ne0$; moreover
$3BC-4\not\equiv0\pmod3$.

For $p\ne2,3$, the $S^2$ coefficient forces the triple root to be
$2/3$. Comparing the other coefficients with $(S-2/3)^3$ gives
\[
BC=\frac43,\qquad AC^2=\frac4{27}
\quad\text{in }\F_p.
\]
These are exactly $p\mid U$ and $p\mid V$.
\end{proof}

\begin{remark}[The geometric source of the gcd]
The locus $U=V=0$ is the missing cusp in the inverse-cubic geometry;
in characteristic zero it is parametrized by
$(A,B,C)=(4/(27t^2),4/(3t),t)$, $t\ne0$ \cite{gao}.
The new arithmetic point is that, \emph{after a rational root is known},
every odd local failure is reduction into this one locus. Merely avoiding
the cusp is not sufficient for an arbitrary irreducible cubic.
\end{remark}

\subsection{Why the quadratic pair need not be analyzed separately}
Suppose a rational root $a$ is known. At a good odd prime, its reduction
may be repeated. If it is double rather than triple, the remaining
linear factor has a simple residue root and lifts to a local integral
point. That point need not arise from a globally rational geometric
point. Thus quadratic splitting does enter the mechanism, but does not
survive as an independent condition in \eqref{alg:ns:eq:oddcriterion}.

This explains why tracking discriminants, Legendre symbols, and the
valuation of each reconstructed nonrational point would overcomplicate
this problem. Such data are useful for finer fiber distributions, but
not for the yes/no local-solubility classification under the root gate.

\section{The prime two and a finite coefficient table}\label{alg:ns:sec:two}

\subsection{A lifting lemma with a nonunit determinant}
\begin{lemma}[Uniform inverse correction]\label{alg:ns:lem:lift}
Let $\Oh$ be a complete discrete valuation ring of characteristic zero,
with valuation $v$, and put $e=v(2)$. For the map $F$, suppose
$a_0\in\Oh^3$ and $b\in\Oh^3$ satisfy
\[
\min_i v(F_i(a_0)-b_i)>2e.
\]
Then there is $a\in\Oh^3$ with $F(a)=b$. In particular, over $\Z_2$
a solution modulo $8$ lifts to a solution over $\Z_2$.
\end{lemma}
\begin{proof}
At each integral point, $J^{-1}=-(1/2)\operatorname{adj}(J)$ has entries
of valuation at least $-e$. If the current error has minimum valuation
$k>2e$, take the Newton correction
$\delta=-J^{-1}(F(a)-b)$. Then $v(\delta_i)\ge k-e>e\ge0$.
The integral polynomial Taylor expansion has linear term canceled and
all remaining terms of degree at least two in $\delta$. Hence the new
error has minimum valuation at least $2k-2e$.

Iterating, the excess $k-2e$ at least doubles, and corrections tend to
zero in valuation. Completeness gives an integral limit, and polynomial
continuity gives $F(a)=b$. For $\Z_2$, an error divisible by $8$ has
$k\ge3>2$.
\end{proof}

\subsection{The complete residue certificate}
\begin{proposition}[Dyadic coefficient criterion]\label{alg:ns:prop:two}
For every integer target,
\begin{equation}\label{alg:ns:eq:dyadic}
(A,B,C)\in F(\Z_2^3)\quad\Longleftrightarrow\quad E(A,B,C).
\end{equation}
The set of admissible targets modulo $8$ has cardinality $176$.
\end{proposition}
\begin{proof}
Lemma~\ref{alg:ns:lem:lift} reduces the statement to the finite image of $F$
on $(\Z/8\Z)^3$. Substituting the $512$ input triples and collecting
outputs gives the table defining $E$ in Section~\ref{alg:ns:sec:intro}.
Section~\ref{alg:ns:app:dyadic} \wtag{}~(the manuscript's Appendix~D)
gives a full exhaustive certificate. It uses
only integer arithmetic and finite loops, and establishes both
inclusions, not just the number of admissible targets.

The certificate has been executed in the supplied verification package. The
finite identity it checks, together with the lifting lemma, proves the
statement at every $2$-adic precision.
\end{proof}

The size $176$ also follows immediately from the coefficient table:
$64$ targets for $C=0$, $32$ for $C=4$, $16$ for each of $C=2,6$, and
$12$ for each of the four odd values of $C$. The corresponding normalized
image measure is $176/512=11/32$, agreeing with the prior dyadic law
\wtag{}~(Part~I, \eqref{alg:eq:rho2} in Theorem~\ref{alg:thm:dyadic}). This measure is not claimed as a new result.

\begin{remark}[Finite fibers are not $2$-adic fibers]
The counts of residue targets with respectively $0,2,4,6$ preimages
modulo $8$ are $336,112,48,16$. These are finite-ring fiber cardinalities.
They must not be confused with cardinalities of fibers over $\Z_2$.
Our argument needs only membership in the image.
\wtag{} These counts are Part~I's \eqref{alg:eq:dyadicfinite} at $m=3$,
$2^9\cdot(21,7,3,1)/32$, and its table in the proof of
Theorem~\ref{alg:thm:dyadic} (total row); Part~I also gives the
$2$-adic fiber law \eqref{alg:eq:dyadiclaw} that this remark
distinguishes them from.
\end{remark}

\section{Complete classification, including nonsplit fibers}\label{alg:ns:sec:complete}

\subsection{Proof of the main certificate}
\begin{proof}[Proof of Theorem~\ref{alg:ns:thm:main}]
If a target is locally integral everywhere, Lemma~\ref{alg:ns:lem:rootgate}
gives the integer root. Theorem~\ref{alg:ns:thm:odd} then says no odd prime
divides $G$, while Proposition~\ref{alg:ns:prop:two} gives $E$.

Conversely, an integer root and a power-of-two gcd imply local integral
solubility at every odd prime by Theorem~\ref{alg:ns:thm:odd}. Condition $E$
supplies the prime $2$. If $C=0$, formula~\eqref{alg:ns:eq:boundary} gives a
global integral point; the conditions are indeed automatic, since
$H=S^2(S-2)$ and $G=4$.
\end{proof}

\subsection{An explicit rational-plus-quadratic formulation}
For an integer root $a$, factor
\begin{equation}\label{alg:ns:eq:factor-H}
H(S)=(S-a)\bigl(S^2+(a-2)S+a^2-2a+BC\bigr).
\end{equation}
The discriminant of the quadratic is
\begin{equation}\label{alg:ns:eq:delta-H}
\Delta_a=-3a^2+4a+4-4BC=(3a-2)^2-4D.
\end{equation}
Here and below, ``a square in $\Z$'' means the square of an integer;
negative integers are not squares.

\begin{corollary}[Complete nonsplit Hasse test]\label{alg:ns:cor:nonsplit}
Suppose $C\ne0$, $H(a)=0$ for an integer $a$, and $\Delta_a$ is not a
square in $\Z$. Then the fiber has exactly one rational point, given by
\eqref{alg:ns:eq:reconstruct}. It is an integral Hasse failure if and only if
$G$ is a power of two, $E(A,B,C)$ holds, and at least one of the three
divisibilities in \eqref{alg:ns:eq:divisibility} fails.
\end{corollary}
\begin{proof}
The discriminant condition says the quadratic factor is irreducible
over $\Q$. In particular $D\ne0$, since otherwise $a$ would also be a
root of that factor. Proposition~\ref{alg:ns:prop:inverse} gives the unique
rational point. Apply Theorem~\ref{alg:ns:thm:main} and
Corollary~\ref{alg:ns:cor:integral-test}.
\end{proof}
This resolves the rational-plus-quadratic local-solubility question for
\emph{all} integral target slices, not only $C=2$. It also supplies the
requested coefficient certificate at every odd prime. What remains
restricted to $C=2$ is the sharp counting law.

\subsection{The three tests are genuinely different}
The following examples isolate the conditions. In the first column,
all targets have $C=2$.
\begin{center}
\begin{tabular}{@{}llll@{}}
\toprule
Target & Rational-root gate & Odd gcd test & Dyadic test\\
\midrule
$(2,0,2)$ & fails & passes & passes\\
$(8,4,2)$ & passes & fails at $5$ & passes\\
$(1,2,2)$ & passes & passes & fails\\
$(0,4,2)$ & passes & passes & passes\\
\bottomrule
\end{tabular}
\end{center}
For $(2,0,2)$ the relevant cubic is $T^3-T^2-2$, which has no rational
root and already has no root modulo $3$. The gcd equals $4$; avoiding
the cusp alone does not help. For $(8,4,2)$ the root is $2$, and the
cubic becomes $(T-2)^3$ modulo $5$. For $(1,2,2)$ the inverse cubic is
$(T-1)(T^2+1)$, but the target fails the parity condition at $2$.
Finally, $(0,4,2)$ will be our smallest displayed nonsplit integral
Hasse failure, with unique rational point $(1/2,-2,4)$.

\section{A primitive parametrization on the slice \texorpdfstring{$C=2$}{C=2}}\label{alg:ns:sec:slice}

\subsection{All locally integral targets}
\begin{theorem}[Primitive root coordinates]\label{alg:ns:thm:param}
An integer target $(A,B,2)$ is everywhere locally integral if and only if
there exist $r,b\in\Z$ such that
\begin{equation}\label{alg:ns:eq:param}
A=r^3-r^2+rb,\qquad B=2b,\qquad
\gcd(3r-1,3b-1)=1.
\end{equation}
Every integer root of the corresponding cubic satisfies the gcd
condition, not just one chosen root.
\end{theorem}
\begin{proof}
The dyadic table forces $B$ and $A$ to be even. Write $B=2b$ and define
\begin{equation}\label{alg:ns:eq:f}
f(T)=T^3-T^2+bT-A.
\end{equation}
The relation with the monic inverse polynomial is $H(2T)=8f(T)$.
The rational-root gate is therefore equivalent to an integer root $r$
of the monic polynomial $f$, or equivalently the first equation of
\eqref{alg:ns:eq:param}.

Now
\[
G=4\gcd(3b-1,27A-1),
\]
and the exact congruence
\begin{equation}\label{alg:ns:eq:congruence-cube}
27A-1\equiv(3r-1)^3\pmod{3b-1}
\end{equation}
shows that its odd prime support is the common odd prime support of
$3r-1$ and $3b-1$. In fact the difference in
\eqref{alg:ns:eq:congruence-cube} is $9r(3b-1)$.
Modulo $2$, $A=r^3-r^2+rb\equiv rb$. The gcd of $3r-1$ and $3b-1$
is even exactly when both $r$ and $b$ are odd, exactly when $A$ is odd.
Thus the odd gcd test and the dyadic test combine into
$\gcd(3r-1,3b-1)=1$. The argument applies to any integer root of $f$.
The converse follows from the same identities and
Theorem~\ref{alg:ns:thm:main}.
\end{proof}

\subsection{Reconstruction and the quadratic discriminant}
Put
\begin{equation}\label{alg:ns:eq:dDelta}
d=3r^2-2r+b=f'(r),\qquad
\Delta=(3r-1)^2-4d=-3r^2+2r+1-4b.
\end{equation}
Then
\begin{equation}\label{alg:ns:eq:factor-f}
f(T)=(T-r)\bigl(T^2+(r-1)T+r^2-r+b\bigr).
\end{equation}
Whenever $d\ne0$, the point belonging to $r$ is particularly simple:
\begin{equation}\label{alg:ns:eq:slice-inverse}
\boxed{\quad
(x,y,z)=\left(\frac1d,\ r-d,\ 5d^2-3rd-2d^3\right).
\quad}
\end{equation}
Indeed, in \eqref{alg:ns:eq:reconstruct} take $a=2r$, $C=2$, and $D=4d$.
The second and third coordinates are integers automatically. Thus the
entire global integrality obstruction is $d\ne\pm1$.

For reference, the primitive gcd may also be tested with the derivative:
\begin{equation}\label{alg:ns:eq:gcd-derivative}
\gcd(3r-1,d)=\gcd(3r-1,3b-1),
\end{equation}
because $3d\equiv3b-1\pmod{3r-1}$ and $3$ is coprime to $3r-1$.
This is an equality of positive gcds, not merely an equality of their
prime supports.

\subsection{An infinite nonsplit Hasse family}
\begin{corollary}\label{alg:ns:cor:family}
For every integer $b\ge2$, the target
\begin{equation}\label{alg:ns:eq:family}
(0,2b,2)
\end{equation}
is an integral Hasse failure with exactly one rational preimage,
\begin{equation}\label{alg:ns:eq:family-point}
\left(\frac1b,-b,5b^2-2b^3\right).
\end{equation}
\end{corollary}
\begin{proof}
Take $r=0$. The primitive gcd is $\gcd(-1,3b-1)=1$, so all local
conditions hold. The remaining quadratic is $T^2-T+b$, whose
discriminant $1-4b$ is negative. Formula~\eqref{alg:ns:eq:slice-inverse} gives
\eqref{alg:ns:eq:family-point}, and its first coordinate is not integral.
\end{proof}
This family makes branch switching visible. At an odd prime dividing
$b$, the rational point is not integral, but the reduction of
$T(T^2-T+b)$ is $T^2(T-1)$: its root $1$ is simple. The quadratic branch
therefore supplies an integral local point. For $b=2$, the nonrational
pair has discriminant $-7$, and the rational point is $(1/2,-2,4)$.

\section{Repeated roots and the entire integral image}\label{alg:ns:sec:integral}

\subsection{The unique repeated-root exception}
\begin{proposition}\label{alg:ns:prop:repeated}
The only everywhere locally integral target on $C=2$ whose inverse
cubic has a repeated root is $(A,B,2)=(0,0,2)$. Its inverse polynomial
is $T^2(T-1)$; its unique rational point in the fiber is $(1,0,0)$.
\end{proposition}
\begin{proof}
A repeated root of a monic cubic over $\Q$ is rational, hence integral.
Write
\[
f(T)=(T-s)^2(T-r),\qquad r=1-2s,\qquad b=2s-3s^2.
\]
Then
\[
3b-1=-(3s-1)^2,
\]
so Theorem~\ref{alg:ns:thm:param}, applied to the root $s$, requires
$|3s-1|=1$. The only integer solution is $s=0$; the apparent alternative
$s=2/3$ is not integral. Therefore $r=1$, $b=A=0$. The repeated root
$0$ does not correspond to a point of the inverse chart; the simple root
$1$ gives $(1,0,0)$. A triple root would require $3s=1$ and is impossible.
\end{proof}
The repeated root is counted once when marking distinct integer roots,
not twice according to multiplicity. This distinction will produce the
constant corrections in the exact counting identities.

\subsection{Split distinct fibers have no integral point}
\begin{lemma}\label{alg:ns:lem:split-nointeger}
On $C=2$, a completely split fiber with three distinct rational points
has no integral point.
\end{lemma}
\begin{proof}
The three roots $r,s,t$ of the monic polynomial $f$ are integers and
satisfy $r+s+t=1$. If the point belonging to $r$ were integral, then
\[
f'(r)=(r-s)(r-t)=\pm1
\]
by \eqref{alg:ns:eq:slice-inverse}. Both integer differences must be $1$ or
$-1$. Since $s\ne t$, the other roots must be $r-1,r+1$. Their sum
would be $3r=1$, impossible for integer $r$.
\end{proof}
This proof is specific to the sum of roots on $C=2$; it should not be
exported without modification to an arbitrary slice.
\wtag{} This is Part~I's Theorem~\ref{alg:thm:plane} (``It has no
integral point''), reproved here by the same argument; the manuscript
does not cite it. The three roots $r,s,t$ of $f$ are Part~I's integer
roots $a,b,c$ of \eqref{alg:eq:planemap}, and $f'(r)$ is Part~I's $D_r$
(Section~\ref{alg:ns:sec:duplicates}). Part~II's
Theorem~\ref{alg:dh:thm:finite} gives the finitely many exceptions on
other slices.

\subsection{Two polynomial curves}
\begin{theorem}[Exact integral-image curves]\label{alg:ns:thm:curves}
The targets in $F(\Z^3)$ with third coordinate $2$ are precisely
\begin{equation}\label{alg:ns:eq:curves}
\begin{aligned}
A&=-2r^3+r^2+\varepsilon r,\\
B&=-6r^2+4r+2\varepsilon,
\qquad r\in\Z,\quad\varepsilon\in\{-1,1\}.
\end{aligned}
\end{equation}
The corresponding integral preimage is
\begin{equation}\label{alg:ns:eq:integralpoint}
(x,y,z)=(\varepsilon,\ r-\varepsilon,\ 5-\varepsilon(3r+2)).
\end{equation}
Distinct pairs $(r,\varepsilon)$ give distinct targets.
\end{theorem}
\begin{proof}
Any integral target is locally integral, so $B=2b$ and the inverse cubic
$f$ is available. Its point corresponds to a simple integer root $r$.
By \eqref{alg:ns:eq:slice-inverse}, integrality is equivalent to
$d=\varepsilon\in\{-1,1\}$. Solving for $b$ gives
$b=\varepsilon-3r^2+2r$. Substitution into
$A=r^3-r^2+rb$ and $B=2b$ gives \eqref{alg:ns:eq:curves}, and
\eqref{alg:ns:eq:slice-inverse} becomes \eqref{alg:ns:eq:integralpoint}. Conversely,
these formulas substitute directly into $F$ and give integral points.

If two different parameter pairs gave the same target, they could not
have the same $r$, since $B$ determines $\varepsilon$ once $r$ is fixed.
There would therefore be two different simple rational roots. A cubic
with two distinct simple rational roots must have three distinct
rational roots. Lemma~\ref{alg:ns:lem:split-nointeger} contradicts the existence
of the integral point. This also handles possible intersections of the
two curves. The repeated-root target has only one simple root, so is
not an exception to injectivity.
\end{proof}

\begin{corollary}[Exact integral-image scale]\label{alg:ns:cor:integral-count}
In the square target box,
\[
I_2(T)=2^{5/3}T^{1/3}+O(1).
\]
Except for $(0,0,2)$, every integral target on this slice has a nonsplit
fiber.
\end{corollary}
\begin{proof}
For each fixed $\varepsilon$, the first coordinate in
\eqref{alg:ns:eq:curves} is $-2r^3+O(r^2)$. The inequality $|A|\le T$ therefore
restricts $r$ to endpoints
$\pm(T/2)^{1/3}+O(1)$. More explicitly, for all sufficiently large
positive or negative $r$ the cubic is strictly monotone on that tail,
and its real threshold roots have these estimates. The number of
integer $r$ is $2(T/2)^{1/3}+O(1)$.
For these $r$, $|B|=O(T^{2/3})<T$ for sufficiently large $T$, so the
second-coordinate bound removes no additional asymptotic range.
There are two disjoint integer-parametrized images, proving the formula. A rational cubic
factorization other than nonsplit would be completely split or repeated;
Lemma~\ref{alg:ns:lem:split-nointeger} and Proposition~\ref{alg:ns:prop:repeated}
leave only the stated exception.
\end{proof}

\section{Marked-root counting and the \texorpdfstring{$T\log T$}{T log T} term}\label{alg:ns:sec:count}

\subsection{Why a logarithm should appear}
The parametrization is nearly quadratic in its center but linear in
$b$. For a nonzero root $r$, keeping $A$ bounded forces $b$ into an
interval of length about $2T/|r|$ centered at $-r^2+r$. The bound on
$B=2b$ permits $|r|$ up to a constant multiple of $\sqrt T$.
Summing $T/|r|$ through this range produces $T\log T$. The constant is
not an arbitrary density: it is the mean primitive density in the
arithmetic progressions $3r\pm1$.

\subsection{The exact parameter region}
Let $P(T)$ denote the number of pairs $(r,b)\in\Z^2$ satisfying
\eqref{alg:ns:eq:param} and $|A|,|B|\le T$. These pairs mark integer roots,
so they are not yet target counts. We may first assume $T$ is an integer;
replacing a real $T$ by $\lfloor T\rfloor$ does not change the conclusions.

For $r\ne0$, the permitted real interval for $b$ is
\begin{equation}\label{alg:ns:eq:interval}
\mathcal I_r(T)=
[-T/2,T/2]\cap
[-r^2+r-T/|r|,\,-r^2+r+T/|r|].
\end{equation}
For $r=0$, every $|b|\le T/2$ is primitive, yielding
$2\lfloor T/2\rfloor+1=O(T)$ pairs.

\begin{lemma}[Range and boundary reduction]\label{alg:ns:lem:range}
Only $|r|=O(\sqrt T)$ contribute. If
$R=\lfloor\sqrt{T/8}\rfloor$, then, for sufficiently large $T$ and
$3\le|r|\le R$, the second interval in \eqref{alg:ns:eq:interval} is contained
in the first. The total contribution from $|r|>R$, together with
$|r|<3$, is $O(T)$.
\end{lemma}
\begin{proof}
Nonemptiness implies
\[
r^2-r\le T/2+T/|r|.
\]
In terms of $q=|r|$, the left side is at least $q^2-q$, so $q=O(\sqrt T)$.
For example the explicit bound $q\le\lfloor\sqrt{2T}\rfloor+3$ is safe
for every integer $T\ge1$: beyond it,
$q^2-q-T/2>T/q$.

If $3\le q\le R$, the absolute center is at most $T/8+\sqrt{T/8}$,
and the radius is at most $T/3$. Their sum is smaller than $T/2$ once
$T$ is sufficiently large. Thus the interval is uncut there and has
length $2T/q$.

For $R<q\le c_0\sqrt T$ \wtag{}~(the manuscript writes $C\sqrt T$; $c_0$ is the implied constant of $q=O(\sqrt T)$ above), forget the gcd condition and bound the number
of possible integers by $2T/q+1$. The harmonic sum over this fixed-ratio
range is $O(1)$, so its contribution is $O(T)$. Each of the finitely
many roots $|r|<3$ also contributes $O(T)$.
\end{proof}

\subsection{Primitive points in each interval}
\begin{lemma}[A uniform interval count]\label{alg:ns:lem:primitive}
For fixed $r\in\Z$ and any interval of length $L$, the number of
integers $b$ in the interval with $\gcd(3r-1,3b-1)=1$ is
\begin{equation}\label{alg:ns:eq:primitivecount}
L\frac{\varphi(|3r-1|)}{|3r-1|}+O(\tau(|3r-1|)),
\end{equation}
with an absolute implied constant. Here $\tau$ is the divisor function.
\end{lemma}
\begin{proof}
Write $n=|3r-1|$, so $3\nmid n$. Inclusion--exclusion gives
\[
\sum_{b\in\mathcal I\cap\Z}
\sum_{d\mid\gcd(n,3b-1)}\mu(d)
=\sum_{d\mid n}\mu(d)
\#\{b\in\mathcal I\cap\Z:3b\equiv1\pmod d\}.
\]
Each congruence is one residue class, hence contributes $L/d+O(1)$.
The main term is $L\sum_{d\mid n}\mu(d)/d=L\varphi(n)/n$, and summing
the errors gives the assertion.
\end{proof}

\begin{lemma}[The arithmetic mean]\label{alg:ns:lem:mean}
For each sign,
\begin{equation}\label{alg:ns:eq:mean}
\sum_{1\le r\le X}\frac{\varphi(3r\pm1)}{3r\pm1}
=\alpha X+O(\log(2X)),
\qquad
\alpha=\prod_{p\ne3}(1-p^{-2})=\frac{27}{4\pi^2}.
\end{equation}
Consequently
\begin{equation}\label{alg:ns:eq:harmonicmean}
\sum_{1\le r\le X}\frac{\varphi(3r\pm1)}{r(3r\pm1)}
=\alpha\log X+O(1).
\end{equation}
\end{lemma}
\begin{proof}
Expand $\varphi(n)/n=\sum_{d\mid n}\mu(d)/d$. Since $3$ never divides
$3r\pm1$, only $3\nmid d$ occur. For each such $d$, the condition
$3r\pm1\equiv0\pmod d$ specifies one residue class, with
$X/d+O(1)$ representatives up to $X$. Thus
\[
\sum_{r\le X}\frac{\varphi(3r\pm1)}{3r\pm1}
=X\!\sum_{\substack{d\le3X+1\\3\nmid d}}\frac{\mu(d)}{d^2}
+O\!\left(\sum_{d\le3X+1}\frac1d\right).
\]
The tail of the absolutely convergent main sum is $O(1/X)$.
Its Euler product is
\[
\alpha=\frac{1}{\zeta(2)(1-3^{-2})}=\frac{27}{4\pi^2}.
\]
This proves \eqref{alg:ns:eq:mean}. Partial summation gives
\eqref{alg:ns:eq:harmonicmean}; the error integral
$\int_1^X O(\log(2t))t^{-2}\,dt$ is bounded.
\end{proof}
\wtag{} The constant $\alpha$ is the density $27/(4\pi^2)$ of Part~I's
\eqref{alg:eq:paramdensity}, obtained there from the same M\"obius sum
$\sum_{3\nmid d}\mu(d)/d^2$; the counted quantities differ (here a mean
of $\varphi(3r\pm1)/(3r\pm1)$ over roots, there a density of ordered
root pairs). Part~I's $\alpha_p$ in Theorem~\ref{alg:thm:thinsieve}
is unrelated (Table~\ref{alg:ns:tab:notation}).

\begin{proposition}[Marked-root asymptotic]\label{alg:ns:prop:P}
\begin{equation}\label{alg:ns:eq:P}
P(T)=\frac{27}{2\pi^2}T\log T+O(T).
\end{equation}
\end{proposition}
\begin{proof}
Apply Lemmas~\ref{alg:ns:lem:range} and \ref{alg:ns:lem:primitive} on the full intervals:
\[
P(T)=2T\!\sum_{3\le|r|\le R}
\frac{\varphi(|3r-1|)}{|r|\,|3r-1|}+O(T).
\]
The accumulated divisor errors are harmless: since the integers
$|3r-1|$ lie below $O(\sqrt T)$ and each is attained at most a bounded
number of times,
\[
\sum_{|r|\le R}\tau(|3r-1|)
\ll\sum_{n\ll\sqrt T}\tau(n)
\ll\sqrt T\log(2T)=O(T).
\]
The last inequality follows directly from
$\sum_{n\le Y}\tau(n)=\sum_{d\le Y}\lfloor Y/d\rfloor\ll Y\log(2Y)$.

Positive roots produce $3r-1$, negative roots produce $3|r|+1$.
Lemma~\ref{alg:ns:lem:mean} therefore makes the displayed sum
$2\alpha\log R+O(1)$. Finally $\log R=\tfrac12\log T+O(1)$, giving
$P(T)=2\alpha T\log T+O(T)$ and the stated constant.
\end{proof}
There is no hidden equidistribution hypothesis in this computation. The
primitive density is obtained by a uniformly controlled divisor sum.

\section{From roots to targets: exact multiplicity corrections}\label{alg:ns:sec:mult}

\subsection{An elementary bound for split targets}
Let $\Sigma(T)$ be the number of locally integral targets in the box with
three distinct rational roots, and let $N(T)$ count locally integral
nonsplit targets. By Lemma~\ref{alg:ns:lem:split-nointeger}, every target
counted by $\Sigma(T)$ is already an integral Hasse failure.

\begin{lemma}\label{alg:ns:lem:splitbound}
$\Sigma(T)=O(T)$.
\end{lemma}
\begin{proof}
For a split cubic write the integer roots as $r,s,t$ with sum $1$.
The second elementary symmetric function is $b$, so
\[
r^2+s^2+t^2=(r+s+t)^2-2b=1-2b\le T+1.
\]
Thus $r,s$ range in a disk inside a square of side $O(\sqrt T)$, and
$t=1-r-s$ is fixed by them. There are $O(T)$ ordered pairs, hence no
more than $O(T)$ targets. Imposing the $A$ bound and the primitive gcd
can only reduce the count.
\end{proof}
This deliberately elementary estimate is enough for our main theorem.
The stronger prior formula \eqref{alg:ns:eq:prior-split} is compatible with,
but not required by, the argument.

\subsection{Exact identities}
For every integer $T\ge1$, Proposition~\ref{alg:ns:prop:repeated} yields
\begin{align}
P(T)&=3\Sigma(T)+N(T)+2,\label{alg:ns:eq:multiplicity-P}\\
L_2(T)&=\Sigma(T)+N(T)+1=P(T)-2\Sigma(T)-1.\label{alg:ns:eq:multiplicity-L}
\end{align}
The $2$ in the first identity is the number of distinct integer roots
$0,1$ of $T^2(T-1)$; the $1$ in the second counts its single target.
Every other locally integral cubic is either nonsplit, with one integer
root, or split distinct, with three.

All integral targets except this repeated-root target are nonsplit.
Therefore the exact Hasse counts are
\begin{align}
H_2^{\NS}(T)&=P(T)-3\Sigma(T)-I_2(T)-1,\label{alg:ns:eq:exact-NS}\\
H_2(T)&=P(T)-2\Sigma(T)-I_2(T)-1.\label{alg:ns:eq:exact-H}
\end{align}
These formulas are also used by the supplied counting program. They
prevent the common error of dividing every marked-root count by three.

\begin{proof}[Proof of Theorem~\ref{alg:ns:thm:count-main}]
Insert Proposition~\ref{alg:ns:prop:P}, Lemma~\ref{alg:ns:lem:splitbound}, and
Corollary~\ref{alg:ns:cor:integral-count} into
\eqref{alg:ns:eq:multiplicity-L}, \eqref{alg:ns:eq:exact-NS}, and \eqref{alg:ns:eq:exact-H}.
The curve statement is Theorem~\ref{alg:ns:thm:curves}. Finally direct
substitution in \eqref{alg:ns:eq:F} gives
\[
F(-x,-y,z)=(P(x,y,z),-Q(x,y,z),-R(x,y,z)).
\]
This bijection preserves square target heights and all integral and
local-solubility conditions, transferring the results to $C=-2$.
\end{proof}

\subsection{Density consequences and the distinction of sample spaces}
\begin{corollary}\label{alg:ns:cor:density}
Among everywhere locally integral targets on $C=2$ in growing square
boxes, the proportion that are integral Hasse failures tends to $1$.
Among those Hasse failures, the proportion that are nonsplit also tends
to $1$. More precisely,
\[
\frac{I_2(T)}{L_2(T)}
\sim\frac{2^{5/3}\,2\pi^2}{27}\,
\frac{T^{-2/3}}{\log T}.
\]
In contrast, the proportion of everywhere locally integral targets
among \emph{all} integer targets in the square tends to $0$.
\end{corollary}
\begin{proof}
Use $L_2(T)\sim(27/(2\pi^2))T\log T$,
$I_2(T)\sim2^{5/3}T^{1/3}$, and $\Sigma(T)=O(T)$. The ambient box has
$(2\lfloor T\rfloor+1)^2$ targets, which is of order $T^2$.
\end{proof}
Thus ``almost every'' here has a specified conditioning set. No positive
ambient density is asserted. The arithmetic image can be thin in the
ambient plane while local-to-global failure is overwhelmingly typical
inside its locally soluble part.

\section{Number fields and rings of \texorpdfstring{$S$}{S}-integers}\label{alg:ns:sec:numberfields}

The collapse of odd local conditions is not special to rational primes.
It uses a rooted cubic over a finite residue field and the invertibility
of $2$. We record a general extension, including the effect of dyadic
ramification.

Let $K$ be a number field and $S$ a finite set of finite places. Put
$R=\Oh_{K,S}$. For a finite place $v\notin S$, write $\Oh_v$ for the
valuation ring of the completion. All target coordinates now lie in
$R$. The ring $R/8R$ is finite; if $2$ is a unit in $R$, it is the zero
ring and its image condition is vacuous.

\begin{theorem}[Finite certificate over $S$-integers]\label{alg:ns:thm:numberfields}
Let $C\ne0$. A target $\target\in R^3$ belongs to $F(\Oh_v^3)$ for
every finite $v\notin S$ if and only if all three conditions hold:
\begin{enumerate}[label=(\alph*)]
\item $H$ has a simple root in $K$;
\item for every $v\notin S$ of odd residue characteristic, at least one
of $U,V$ is a $v$-adic unit;
\item the residue of $\target$ in $(R/8R)^3$ lies in
$F((R/8R)^3)$.
\end{enumerate}
Condition (b) is equivalently that the ideal $(U,V)R$ has no prime
support outside $S$ and the dyadic places. For $C=0$, every target has
the global $R$-point \eqref{alg:ns:eq:boundary}.
\end{theorem}
\begin{proof}
Suppose local solubility holds outside $S$. If $H$ is irreducible over
$K$, its transitive cubic Galois group has a $3$-cycle. Chebotarev gives
infinitely many residue fields of irreducible reduction, avoiding the
finite sets $S$, the dyadic places, the primes of $C$, and the bad
reduction primes. No local root is possible there. Hence $H$ is
reducible. A reducible cubic has a simple $K$-root unless it is a triple
linear factor. The triple case has no point over any extension field
by Proposition~\ref{alg:ns:prop:inverse}, contradicting local solubility.
This proves (a).

A root in $K$ of the monic polynomial $H\in R[T]$ belongs to the
integrally closed ring $R$. Reducing it at any $v\notin S$ gives a
residue root. At odd residue characteristic the proof of
Theorem~\ref{alg:ns:thm:odd} works verbatim over the residue field: places
dividing $C$ are automatically soluble, residue characteristic $3$
permits no triple root, and in every other odd characteristic the only
obstruction is $U=V=0$. This proves the equivalence of the odd local
conditions with (b), in both directions.

It remains to handle dyadic places not inverted in $R$. Necessity of
(c) follows by reducing the local solutions modulo $8\Oh_v$ and
combining their residues by the Chinese remainder theorem. Indeed,
$R/8R$ is the product of the corresponding finite local quotients.
Conversely, (c) supplies an integral approximate solution at each such
place. If $e=v(2)$, its error has valuation at least $3e>2e$, so
Lemma~\ref{alg:ns:lem:lift} gives a genuine local point. Thus (a)--(c) are
sufficient as well.
\end{proof}

\begin{remark}[Why ``simple'' is explicit in the extension]
Over $\Z$, the triple polynomial $(S-2/3)^3$ cannot have integral
coefficients. Over $S$-integers with $3$ inverted it can. The word
``simple'' in condition (a) prevents importing that integer-specific
shortcut into the more general theorem. Equivalently, one could allow
any $K$-root and separately exclude the triple case.
\end{remark}

The dyadic modulus $8R$ is a uniform sufficient modulus, including
ramified dyadic fields. We do not claim that it is minimal. The compact
rational coefficient table $E$ is special to $\Z/8\Z$; for a general
number field the finite image in (c) must be computed in its actual
quotient ring, not by applying that table componentwise to coordinates
of an arbitrary integral basis.

Once a simple root is known, global $R$-membership is likewise finite:
reconstruct all simple $K$-roots by \eqref{alg:ns:eq:reconstruct} and test
membership of their coordinates in $R$. The theorem therefore supplies
a complete local/global decision procedure, although no optimal
complexity bound is claimed for number-field factorization.

\section{Algorithms and reproducible checks}\label{alg:ns:sec:computation}

\subsection{An exact classification algorithm}
For integer input $(A,B,C)$, the following procedure has no prime-search
loop and needs no factorization of $G$.
\begin{enumerate}[label=\arabic*.]
\item If $C=0$, return the point $(0,B,A-4B^2)$.
\item Compute $G=\gcd(|3BC-4|,|27AC^2-4|)$ and repeatedly divide out $2$.
Reject local solubility if the remaining odd part exceeds $1$, or if
$E(A,B,C)$ fails.
\item Compute the rational roots of the monic cubic $H$. They are
integers. If none exist, reject local solubility.
\item Reconstruct all simple roots exactly. Return ``integral image''
if a reconstructed point is integral and ``integral Hasse failure''
otherwise. The quadratic discriminant distinguishes the nonsplit case.
\end{enumerate}
The order is chosen for speed: the inexpensive gcd and dyadic tests can
reject many inputs before polynomial factorization. Correctness follows
from Theorem~\ref{alg:ns:thm:main}; the root gate is never silently omitted.
The supplied \rpath{code/classify.py} implements this procedure with
exact rational arithmetic and SymPy's rational-root computation.

\subsection{Exact counting without target-box scanning}
For each root $r$, compute the integer endpoints of \eqref{alg:ns:eq:interval}.
For each squarefree divisor $d$ of $|3r-1|$, the permitted residue class
is $b\equiv3^{-1}\pmod d$. Inclusion--exclusion counts all primitive
$b$ in the interval exactly by floor differences. This evaluates $P(T)$
without scanning $O(T^2)$ target pairs.

To count $\Sigma(T)$, enumerate ordered roots $r<s<t$ with $r+s+t=1$ and
$r^2+s^2+t^2\le T+1$. Compute $A=rst$, $B=2(rs+rt+st)$, check the box
and gcd, and count each unordered triple once. The integral curves give
$I_2(T)$ directly. Equations~\eqref{alg:ns:eq:exact-NS} and \eqref{alg:ns:eq:exact-H}
then give the exact Hasse counts.

This is a practical exact algorithm, not a claimed optimal complexity
result. The delivered split enumerator uses $O(T)$ candidate pairs;
trial division of $|3r-1|$ is used in the marked-root counter. Both can
be improved by sieving, but no asymptotic theorem here depends on such
an optimization.

\subsection{Recorded exact counts}
The following outputs were generated by the included programs. The last
column is illustrative floating-point arithmetic; all counts are exact.
\begin{center}\small
\begin{tabular}{@{}r r r r r r@{}}
\toprule
$T$ & $P(T)$ & $\Sigma(T)$ & $I_2(T)$ & $H_2^{\NS}(T)$ &
$H_2^{\NS}/(T\log T)$\\
\midrule
$10^2$ & $633$ & $21$ & $14$ & $555$ & $1.205167$\\
$10^3$ & $9\,462$ & $121$ & $32$ & $9\,066$ & $1.312438$\\
$10^4$ & $126\,107$ & $653$ & $68$ & $124\,079$ & $1.347171$\\
$10^5$ & $1\,575\,664$ & $3\,304$ & $148$ & $1\,565\,603$ & $1.359865$\\
\bottomrule
\end{tabular}
\end{center}
The predicted leading constant is
$27/(2\pi^2)\approx1.3678359792$.
For comparison, the total Hasse counts in these rows are
$576$, $9\,187$, $124\,732$, and $1\,568\,907$.
The proximity of the ratios is evidence of consistency, not a proof of
the asymptotic or its error term.

\subsection{Independent verification layers}
The package records the following completed checks.
\begin{enumerate}[label=(\roman*)]
\item Exact symbolic checks of the determinant, inverse identities,
reconstruction, factorization, discriminant, and the coefficient-gcd
identity, using SymPy.
\item Exhaustive comparison of the dyadic coefficient table with all
$512$ targets modulo $8$.
\item Exhaustive image and root-cusp comparison over
$p=3,5,7,11,13,17,19,23,29,31$, totaling $82\,142$ target triples.
\item Direct substitution and inverse recovery for $1\,331$ integer
input points in $[-5,5]^3$.
\item Independent target-box enumeration at $T=5,10,25,50$, checking
the exact multiplicity corrections, integral curves, and nonsplit
counts; and $100$ exact instances of Corollary~\ref{alg:ns:cor:family}.
\end{enumerate}
The finite dyadic enumeration is a complete finite certificate for the
residue-table claim. The other finite tests do not establish statements
at untested primes or arbitrarily large heights; those statements have
separate proofs in the article.

\section{Formalization plan and dependency audit}\label{alg:ns:sec:formal}

\subsection{What is and is not kernel checked}
The source repository contains a Lean declaration for the displayed map
and a theorem computing its formal multivariate-polynomial Jacobian
determinant \cite{dh:proveitformal}. This article did not run that Lean build.
Its own classification, counting, and number-field results are not
formalized in Lean or Rocq, and no such status is inferred from the
existing project.

The arithmetic proofs naturally decompose into a mostly elementary
formalization target. First, formalize the inverse chart as identities
in a localization, with the conditions $2C D\ne0$ explicit. Second,
formalize the finite-field cubic lemma and coefficient comparison.
Third, implement the dyadic table as a decidable finite statement and
prove the Newton lifting lemma. Fourth, formalize the primitive
parametrization and the two integral curves. These four stages contain
the main classification innovation. Chebotarev is needed only for the
necessity of the global rational-root gate, and should remain an
explicit imported theorem or an explicitly stated dependency until an
appropriate formal library result is available.

For the counting theorem, the critical analytic inputs are modest:
M\"obius inversion, the Euler product for $\zeta(2)$, an elementary
divisor-sum bound, and partial summation. A clean formal proof should
keep $P,S,N,L,I,H$ distinct rather than hiding multiplicity corrections
inside a single counting function.

\subsection{Dependency graph}
\begin{center}\small
\begin{tabular}{@{}p{.35\textwidth}p{.59\textwidth}@{}}
\toprule
Claim & Essential dependencies\\
\midrule
Odd gcd criterion & Inverse chart; unit-Jacobian lifting; rooted cubic lemma\\
Dyadic table & $512$-case finite identity; nonunit Newton lifting\\
Complete local classification & Odd criterion; dyadic table; cubic Chebotarev gate\\
Primitive parametrization & Classification; exact coefficient congruence\\
Integral curves & Slice reconstruction; distinct-root derivative product\\
$T\log T$ count & Primitive parametrization; divisor inversion; harmonic mean;
 multiplicity correction\\
Number-field extension & Same rooted cubic lemma; Chebotarev over $K$;
 integrally closed rings; dyadic lifting and CRT\\
\bottomrule
\end{tabular}
\end{center}
No proof of the main counting constant uses numerical extrapolation,
the previously claimed split asymptotic, or a conjectural equidistribution
statement.

\section{Further research questions of Part~III}\label{alg:ns:sec:questions}
The results above settle the nonsplit local-solubility question and the
leading square-box asymptotic on $C=2$. They also suggest several
substantial extensions. The questions below specify what remains open
in this article; they are not all claimed to be previously published
open problems.

\begin{question}[A genuine second-order term]\label{alg:ns:q:second}
Does the limit
\[
\lim_{T\to\infty}\frac{H_2^{\NS}(T)-\frac{27}{2\pi^2}T\log T}{T}
\]
exist, and can it be expressed by an explicit convergent arithmetic
series and a boundary integral?
\end{question}
The present proof absorbs the clipped intervals and finitely many small
roots into $O(T)$, precisely the scale of the desired term. Improving it
requires retaining their contributions rather than merely sharpening
the divisor error. The summatory mean in Lemma~\ref{alg:ns:lem:mean} is already
strong enough to make such a refinement plausible.

\begin{question}[Power-saving remainder]\label{alg:ns:q:power}
After a second main term is identified, can the remainder be bounded by
$O(T^{1/2+\epsilon})$, or is a larger fluctuation forced by the changing
integer boundaries and split-root subtraction?
\end{question}
The known split contribution is of order $T^{2/3}$, so the answer for
all Hasse failures and for nonsplit failures may require explicitly
subtracting different $T^{2/3}$ terms. A proposed error smaller than
that scale must account for this distinction.
\wtag{} Part~II's Research question~\ref{alg:dh:q:remainder} asks the
analogous question for the split count $\#\Hsplit_c(T)$ itself.

\begin{question}[Other fixed slices]\label{alg:ns:q:slices}
For each fixed nonzero integer $c$, determine the order and leading
constant for nonsplit integral Hasse failures on $C=c$ in square boxes.
Which $c$ share the $T\log T$ law, and what local factors replace
$27/(2\pi^2)$?
\end{question}
The complete local certificate already applies. The difficulty is the
height geometry and the divisibility conditions in the monic-root
parametrization. The special slice $c=2$ removes denominator congruences
and makes two reconstructed coordinates automatically integral; that
simplification cannot simply be assumed for other $c$.
\wtag{} This is the part of Part~II's Research
question~\ref{alg:dh:q:quadratic} (exponent and leading constant of the
rational-plus-quadratic sector) that Part~III leaves open: its
classification half is answered on every slice by
Corollary~\ref{alg:ns:cor:nonsplit}, its counting half on $C=\pm2$ by
Theorem~\ref{alg:ns:thm:count-main}. The split sector of every slice is
Part~II's Theorem~\ref{alg:dh:thm:headline-count}.

\begin{question}[Unequal boxes and phase changes]\label{alg:ns:q:unequal}
Determine the asymptotic in $|A|\le H_A$, $|B|\le H_B$ uniformly in
natural ranges of the ratio $H_A/H_B^{3/2}$.
\end{question}
The root interval has radius $H_A/|r|$ but the admissible centers extend
to $|r|\asymp\sqrt{H_B}$. The logarithmic range can shrink or disappear
as these two scales cross. A uniform result would explain the mechanism
behind the square-box logarithm rather than only one specialization.
\wtag{} This is what remains of Part~I's Research
question~\ref{alg:sub:q-height}, which asks for boxes $|A|\le H_A$,
$|B|\le H_B$; Part~III answers it for square boxes in all sectors.

\begin{question}[Quadratic fields and branch-switching statistics]\label{alg:ns:q:quadfields}
Among nonsplit Hasse failures, what is the distribution of the quadratic
field $\Q(\sqrt\Delta)$ and of the primes at which the unique rational
point is nonintegral but a quadratic point is locally integral?
\end{question}
The gcd criterion removes these data from the membership test, but not
from finer arithmetic statistics. A field-conditioned count would
involve square classes of $-3r^2+2r+1-4b$ in the primitive region.

\begin{question}[Minimal dyadic precision over number fields]\label{alg:ns:q:dyadicprec}
For a finite extension of $\Q_2$, what is the least ideal modulus that
determines membership in $F(\Oh_v^3)$? How does it depend on ramification
and residue degree?
\end{question}
Theorem~\ref{alg:ns:thm:numberfields} supplies the uniform modulus $8\Oh_v$,
not an optimal one. A classification of the finite image could generalize
the small rational table without an exponential enumeration of the
entire quotient cube.
\wtag{} Cf.\ Part~I's Research question~\ref{alg:sub:q-extensions},
which asks for the integral fiber distributions over finite extensions
of $\Q_2$; this question asks only for the membership modulus.

\begin{question}[Beyond this cubic map]\label{alg:ns:q:beyond}
Which polynomial maps with an inverse polynomial of degree $d>3$ admit
an analogous collapse from local factorization data to a small set of
coefficient obstructions after one rational branch is known?
\end{question}
Lemma~\ref{alg:ns:lem:triple} is unusually strong in degree three. In degree
four a repeated rational factor can leave an irreducible quadratic with
no residue root, so the same proof fails. Identifying the correct
replacement would separate genuinely cubic behavior from a reusable
arithmetic method.

\begin{question}[Counting over number fields and $S$-integers]\label{alg:ns:q:numberfield}
For a fixed number field and a specified archimedean height, determine
the asymptotic number of locally integral but globally nonintegral
nonsplit targets on a fixed slice. How do units and the ideal class group
alter the logarithmic factor?
\end{question}
The certificate in Section~\ref{alg:ns:sec:numberfields} settles the local
membership problem but does not settle height counting. The meaning of
a box and the treatment of units must be fixed before formulating a
leading constant.
\wtag{} Part~II's Research question~\ref{alg:dh:q:numberfield} asks the
same for split fibers (the constant $\kappa(c)$ over number fields and
after inverting finitely many primes).

\begin{question}[Proof-producing arithmetic software]\label{alg:ns:q:software}
Can the rational-root, gcd, dyadic-residue, and reconstruction certificates
be emitted by a small program and checked inside ProveIt's Lean project,
leaving only the general rational-root necessity theorem dependent on
Chebotarev?
\end{question}
Such a checker would certify individual global and integral-Hasse
examples without trusting a computer algebra system. For positive local
certificates, the root and gcd witnesses are especially compact.
\wtag{} Cf.\ Part~I's Research question~\ref{alg:sub:q-library} and
Part~II's Research question~\ref{alg:dh:q:library}, which ask for a
reusable library; this question asks for per-target certificates.

\section{Conclusion of Part~III}\label{alg:ns:sec:conclusion}
The unresolved nonsplit sector is not a minor correction to the split
case. It changes the dominant arithmetic scale. Once a rational root
is available, every odd local obstruction is detected by reduction to
the missing cubic cusp; the entire infinite family of odd-prime tests
reduces to the gcd of two explicit target polynomials. The dyadic test
is finite and uniform.

On $C=2$, these facts give a primitive integer parametrization. Its
harmonic-width geometry yields the coefficient
$27/(2\pi^2)$ in a $T\log T$ law. Global integrality, meanwhile, forces
a derivative to be a unit and collapses the image to two curves with
only $T^{1/3}$ targets. This establishes a sharp qualitative separation:
locally integral targets are sparse in the ambient plane, but almost
all of them are nonsplit integral Hasse failures.

\section{Package contents and replay instructions of Part~III}\label{alg:ns:app:package}
\wtag{} This is the manuscript's Appendix~A.
The archive contains this article as \rpath{article.tex} and
\rpath{article.pdf}, together with:
\begin{center}\small
\begin{tabular}{@{}p{.34\textwidth}p{.60\textwidth}@{}}
\toprule
File & Purpose\\
\midrule
\rpath{code/arithmetic.py} & Exact map, residue test, reconstruction,
primitive interval counts, and target counts\\
\rpath{code/classify.py} & Exact integer-target classifier with JSON output\\
\rpath{code/verify.py} & Dyadic certificate, finite-field checks, inverse
checks, independent small-box enumeration\\
\rpath{code/symbolic.py} & Exact symbolic identity checks\\
\rpath{code/count.py} & Reproduce the counting table\\
\rpath{data/*.json} & Recorded outputs of the verification and counting runs\\
\rpath{README.md} & Build and execution commands\\
\rpath{STATUS.md} & Proof, novelty, formalization, and computation boundaries\\
\rpath{SOURCES.md} & Pinned source and attribution ledger\\
\rpath{SHA256SUMS} & Integrity checks for the delivered files\\
\bottomrule
\end{tabular}
\end{center}
From the package root, run
\Needspace{9\baselineskip}
\begin{lstlisting}[language=bash]
python code/verify.py
python code/symbolic.py
python code/count.py
python code/classify.py 0 4 2
latexmk -pdf -interaction=nonstopmode -halt-on-error article.tex
\end{lstlisting}
The exact arithmetic and finite verification programs need only Python's
standard library. Symbolic verification and the general classifier use
SymPy. Counts and proofs do not depend on floating-point calculations.
The package uses Python 3.10 or later; the recorded run used Python 3.13
and SymPy 1.14.0.

\wtag{} In this report the files carry the prefix
\texttt{03-nonsplit-hasse-}: the programs and the \texttt{Makefile} are
under \rpath{code/}, the recorded outputs and \texttt{requirements.txt}
under \rpath{data/}, and \texttt{STATUS.md} and \texttt{SOURCES.md} at
the report root. The manuscript is printed here as Part~III; its
\texttt{article.pdf}, \texttt{README.md} and \texttt{SHA256SUMS} are not
shipped. The commands above are the delivery's: \texttt{classify.py},
\texttt{count.py}, \texttt{symbolic.py} and \texttt{verify.py} import
\texttt{arithmetic} by its delivered name, and \texttt{verify.py},
\texttt{symbolic.py} and \texttt{count.py} write \texttt{verification.json},
\texttt{symbolic.json} and \texttt{counts.json} into \rpath{../data/}
relative to their own directory, so they must be run on a copy
(Appendix~\ref{alg:ns:app:rerun}). On Windows use \texttt{py} or
\texttt{uv run} for \texttt{python}.

\section{Explicit finite bounds used by the programs of Part~III}\label{alg:ns:app:bounds}
\wtag{} This is the manuscript's Appendix~B.
The marked-root enumerator loops through
$|r|\le\lfloor\sqrt{2T}\rfloor+3$. To verify that this is safe, put
$q=|r|$. A feasible pair would satisfy
\[
q(q^2-q-T/2)\le T.
\]
For $q\ge\sqrt{2T}+2$, the expression $q^2-q-T/2$ is at least
$3T/2+3\sqrt{2T}+2>T$, so feasibility is impossible. The chosen integer
bound is slightly looser, not tighter, than needed.

For split counting, $r^2+s^2+t^2\le T+1$ bounds each root. Enumerating
$r<s<t$ with $t=1-r-s$ is unique; it does not divide an ordered count by
six near repeated roots. For integral curves, the inequality
$|-6r^2+4r+2\varepsilon|\le T$ implies a quadratic bound on $|r|$;
$\lfloor\sqrt{T+1}\rfloor+3$ is again deliberately safe. The sharp
$T^{1/3}$ scale is proved from the first coordinate, not imposed as an
unverified truncation in the program.

\section{Small-box cross-checks of Part~III}\label{alg:ns:app:small}
\wtag{} This is the manuscript's Appendix~C.
These values were obtained both from the root-parametrized algorithm
and by independently scanning all integer target pairs in the indicated
boxes and computing their roots.
\begin{center}\small
\begin{tabular}{@{}r r r r r r r@{}}
\toprule
$T$ & $P$ & $\Sigma$ & $I_2$ & $L_2$ & $H_2^{\NS}$ & $H_2$\\
\midrule
$5$ & $14$ & $1$ & $4$ & $11$ & $6$ & $7$\\
$10$ & $33$ & $2$ & $5$ & $28$ & $21$ & $23$\\
$25$ & $104$ & $5$ & $8$ & $93$ & $80$ & $85$\\
$50$ & $261$ & $10$ & $12$ & $240$ & $218$ & $228$\\
\bottomrule
\end{tabular}
\end{center}
For example, $P=14$ at $T=5$ is not a target count: its three-root
fiber contributes three markings, and the repeated-root target
contributes two. The corrections give $L_2=14-2\cdot1-1=11$.

\section{The complete dyadic finite certificate of Part~III}\label{alg:ns:app:dyadic}
\wtag{} This is the manuscript's Appendix~D.
This short standalone program verifies the exact equality between the
residue image and the coefficient table. No symbolic algebra package is
used. There are $512$ source triples and $512$ target triples; the set
comprehension computes the former image and the final loop checks every
possible target. Polynomial evaluation is performed over the integers
and then reduced modulo $8$, which agrees with evaluation in the quotient
ring.

\begin{lstlisting}[language=Python]
from itertools import product

def F(x, y, z):
    u = 1 + x*y
    h = u*u*z + y*y*(1 + 3*u)
    return u*h, y + 3*x*h, x*(5 - 3*u - x*x*z)

def E(A, B, C):
    c = C % 8
    if c == 0: return True
    if c == 4: return B % 2 == 0
    if c in (2, 6): return A % 2 == B % 2 == 0
    return ((A % 2 == 0 and B % 4 == 2)
            or (A % 2 == 1 and (B*C - 2*A - 1) % 8 == 0))

image = {tuple(v % 8 for v in F(*a))
         for a in product(range(8), repeat=3)}
assert len(image) == 176
for b in product(range(8), repeat=3):
    assert (b in image) == E(*b)
\end{lstlisting}

The recorded run passed every assertion. The same comparison appears in
\rpath{code/verify.py}, which also records the finite fiber histogram.
For readers replacing executable enumeration by a proof-assistant
certificate, the final assertion is a decidable statement on a finite
type. Its formalization is separate from the valuation-theoretic
argument lifting the residue image to $\Z_2$.

The cardinality check alone would not prove the coefficient table:
two different subsets can both have $176$ elements. The final loop
certifies the stronger pointwise equivalence used in the theorem.

%% End of Part III. The appendices below belong to Part I, apart from
%% the Part II and Part III provenance and rerun subsections at the end
%% of Appendix C.
\clearpage
\fancyhead[L]{\small Arithmetic local--global dichotomies}
\phantomsection\addcontentsline{toc}{part}{Appendices}
\appendix
\section{Replay guide and artifact inventory}\label{alg:app:replay}
The archive is self-contained except for standard TeX packages and the
optional SymPy dependency. From its root directory:
\begin{lstlisting}[language=bash]
python3 code/verify_finite.py --output certificates/finite_checks.json
python3 code/verify_split_plane.py
python3 code/verify_symbolic.py
python3 code/rational_inverse.py 0 -4 2
pdflatex -interaction=nonstopmode -halt-on-error arithmetic_fibers.tex
pdflatex -interaction=nonstopmode -halt-on-error arithmetic_fibers.tex
\end{lstlisting}
The first two commands use only Python's standard library. The last two
Python utilities need SymPy; the executed symbolic version was 1.14.0.
The PDF build uses ordinary packages supplied in a full TeX Live
installation, and does not require shell escape or external figures.

\begin{center}\small
\begin{tabular}{p{0.43\textwidth}p{0.47\textwidth}}
\toprule
Path & Purpose\\
\midrule
\texttt{arithmetic\_fibers.tex} & Complete article source, bibliography included\\
\texttt{arithmetic\_fibers.pdf} & Compiled article\\
\texttt{code/verify\_finite.py} & Exact local tables, CRT construction, finite checks\\
\texttt{code/verify\_split\_plane.py} & Split-plane gcd and parameter counts\\
\texttt{code/verify\_symbolic.py} & Universal algebraic identity checks\\
\texttt{code/rational\_inverse.py} & Complete rational inverse utility\\
\texttt{certificates/*.json} & Machine-readable executed results\\
\texttt{certificates/*.txt} & Human-readable execution transcripts\\
\texttt{README.md} & Build instructions and statement-status summary\\
\texttt{SOURCE\_REVIEW.md} & Pinned source paths and novelty boundaries\\
\bottomrule
\end{tabular}
\end{center}
No external article PDFs, font files, or copies of third-party research
repositories are redistributed in this archive.

\wtag{} In ProveIt the article is \rpath{article.tex} (this text, rewritten
as described in \cref{alg:app:report}), \rpath{article.pdf} is a build of
it, \texttt{certificates/} is shipped as \rpath{data/}, the delivered
README is replaced by the report README, and the delivered PDF
\texttt{arithmetic\_fibers.pdf} is not shipped. The commands above are the
delivery's; \cref{alg:app:rerun} gives commands that work with the shipped
names.
[Added 30 September 2026, batch~58: Part~II's delivered replay commands
are printed in Section~\ref{alg:dh:sec:repro}; commands that work with
its shipped names are in Appendix~\ref{alg:dh:app:rerun}.]
[Added 30 September 2026, batch~64: Part~III's delivered replay commands
are printed in Section~\ref{alg:ns:app:package}; commands that work with
its shipped names are in Appendix~\ref{alg:ns:app:rerun}.]

\section{A characteristic-three simplification}
The separate characteristic-three finite-field count has a direct
algebraic explanation. Reducing \eqref{alg:eq:F} modulo $3$ gives
\[
 Q=y,\qquad P=z+y^2-y^3R,\qquad R=-x-x^3z.
\]
For a target $(A,B,C)$, set
$\zeta=A-B^2+B^3C$. The first two equations force $y=B$, $z=\zeta$;
the remaining equation is
\[
 \zeta x^3+x+C=0.
\]
Its derivative with respect to $x$ is one. If $\zeta=0$, it has exactly
one geometric root; if $\zeta\ne0$, it has three distinct geometric roots.
Over a finite field its solution set is either empty or a translate of
the kernel of the additive map $x\mapsto\zeta x^3+x$, whose kernel
has one or three elements. This confirms the possible rational fiber
sizes $0,1,3$, without carrying characteristic-zero cusp formulas
through an invalid division by three. It is an explanatory calculation,
not a new finite-field histogram claim.

\section{Provenance of this report}\label{alg:app:report}

\subsection{One manuscript}\label{alg:app:onesource}
\wtag{} This report is built from \emph{one} manuscript: manuscript~03 of
batch~36 of ProveIt's incoming reports, the archive
\texttt{ProveIt\_Arithmetic\_Local\_Global} (inner directory
\texttt{proveit\_arithmetic\_fibers}, main file
\texttt{arithmetic\_fibers.tex}, a 26-page PDF dated 24~September 2026,
author line ``Research draft prepared for Vladimir Reshetnikov'', PDF
metadata ``Research draft prepared with ChatGPT''), pinned to ProveIt
commit \nolinkurl{e21766d04c2b8a9b2cdba0cd43e563024b3bd1b9}. The archive
arrived in commit \texttt{e13affd32} and was placed in \texttt{1a1396d4d}
as a new report. Nothing was merged and nothing was selected away: every
result, proof, example, table, question and limitation of the manuscript
is printed. The manuscript contributed everything except the text marked
\wtag.
[Added 30 September 2026, batch~58: this subsection describes Part~I. The
report is now built from \emph{two} manuscripts; Part~II is batch-58
manuscript~01, pinned to \dhsha, and its provenance, shipped files and
editorial changes are in Appendix~\ref{alg:dh:app:provenance}.]
[Added 30 September 2026, batch~64: the report is now built from
\emph{three} manuscripts; Part~III is batch-64 manuscript~02, pinned to
\nssha, and its provenance is in Appendix~\ref{alg:ns:app:provenance}.]

The editorial changes are these.
\begin{enumerate}[label=(\arabic*)]
\item Every label carries the prefix \texttt{alg:}. The source's 77 labels
are kept, unchanged after the prefix. Six were added:
\texttt{alg:sub:attribution}, \texttt{alg:rem:collisionfiber},
\texttt{alg:app:report}, \texttt{alg:app:onesource},
\texttt{alg:app:pinned} and \texttt{alg:app:rerun}.
\item Attribution was made precise: a bracketed \wtag{} clause in the
abstract, the status note on the title page, and
\cref{alg:sub:attribution}, which names the formalized statements of the
project and says where the complex stratification and the finite-field
laws come from.
\item \Cref{alg:rem:collisionfiber} records an observation made at the
write, not a claim of the manuscript. It is the only new numbered item and
the last one of its section, so no source number moved.
\item Notes marked \wtag{} were added in \cref{alg:sec:verification} and
in Appendix~\ref{alg:app:replay} (shipped names), and this appendix is
new.
\end{enumerate}
No mathematical statement of the manuscript was changed and no symbol was
renamed. The manuscript's $t=y+1/x$ is not the project's $t$ (which is
$xy$ in \rpath{Algebra/JacobianConjecture/Research/README.md} and the
family parameter in the project README, which
\cref{alg:rem:collisionfiber} writes $\tau$), and its $p$, $q$ are a prime and a
field size, not the project's $p=xy^2$, $q=x^2yz$. The PDF bookmark of the
title page was fixed (no duplicate page anchor), and its top space was
reduced from 0.65 to 0.3~inch so that the added status note fits.

\subsection{Repository statements at the pin}\label{alg:app:pinned}
\wtag{} The source read the root README, the project README and
\rpath{Algebra/JacobianConjecture/Research/README.md} at its pin (its
delivered \rpath{SOURCE_REVIEW.md}). Between the pin and the placement
commit the project did not change, so its statements about ProveIt are
current. They are accurate: the formal developments, the integral stable
shears of \cref{alg:prop:stable} (unitriangular source and target maps
over $\Z$), and the refusal to upgrade the project's SymPy frontier
certificates to kernel proofs. The manuscript claims no gap in the
project; the question \eqref{alg:eq:localglobal-question} is its own. The
two external sources were checked at the write
(\cref{alg:sub:attribution}); Milne was not.

\subsection{Rerunning the checks}\label{alg:app:rerun}
\wtag{} Rerun on a copy with the delivered layout, never in the report:
\rpath{verify_split_plane.py} and \rpath{verify_symbolic.py} always
write \rpath{split_plane_checks.json} and \rpath{symbolic_checks.json}
into \rpath{../certificates/} relative to their own directory, and fail
after passing their checks if it does not exist.
\begin{lstlisting}[language=bash]
mkdir <scratch>          # outside the repository
cp -r code <scratch>/code
cp -r data <scratch>/certificates
cd <scratch>
python code/verify_finite.py --output certificates/finite_checks.json
python code/verify_split_plane.py
python code/verify_symbolic.py
python code/rational_inverse.py 0 -4 2
\end{lstlisting}
The last two need SymPy (the recorded runs used 1.14.0 and Python 3.13.5;
no requirements file was delivered). At the write the four commands were
run so with Python 3.13.5 and SymPy 1.14.0: every regenerated file equals
the shipped one apart from Windows line endings.
\rpath{data/build_review.json} is a static build record, not regenerated.

\subsection{Part~II: the second manuscript}\label{alg:dh:app:provenance}
\wtag{} [Added 30 September 2026, batch~58.] Part~II is built from
\emph{one} manuscript: manuscript~01 of batch~58 of ProveIt's incoming
reports, the archive \texttt{ProveIt\_Dense\_Integral\_Hasse} (inner
directory of the same name, main file \texttt{article.tex}, a 20-page
US-Letter PDF dated 29~September 2026, author line and PDF metadata
``AI-assisted research draft prepared for Vladimir Reshetnikov''), pinned
to ProveIt commit \dhsha. The archive arrived in commit
\texttt{b2b626d85} and was placed in \texttt{b30441a8c} as an addition
to this report. The report is therefore built from two manuscripts:
Part~I (batch~36, manuscript~03; Appendix~\ref{alg:app:onesource}) and
Part~II. Nothing was merged across the Parts and nothing was selected
away: every result, proof, example, table, question and limitation of the
manuscript is printed in
Sections~\ref{alg:dh:sec:intro}--\ref{alg:dh:sec:repro}, and it
contributed everything there except the text marked \wtag.
Section~\ref{alg:dh:sec:relation} is new.

Part~I is unchanged from Part~II's pin to the placement commit (the only
change to this directory in between is the placed files), so the
manuscript's references to Sections 2, 4 and 12 of Part~I and to Research
questions 12.2--12.4 are current. The manuscript also read the
project's \rpath{Algebra/JacobianConjecture/Research/README.md} and the
README of the neighbouring report \rpath{weighted-keller-rigidity}; it
uses no result of either (\rpath{02-dense-hasse-SOURCES.md}).

The editorial changes are these.
\begin{enumerate}[label=(\arabic*)]
\item \emph{Labels.} The manuscript's 70 labels carry the prefix
\texttt{alg:dh:}, unchanged after it. Twenty were added for Part~II: the
section, its six subsections and the notation table
(\rpath{alg:dh:sec:relation}, \rpath{alg:dh:sec:source},
\rpath{alg:dh:sec:answers}, \rpath{alg:dh:sec:consistency},
\rpath{alg:dh:sec:notation}, \rpath{alg:dh:sec:duplicates},
\rpath{alg:dh:sec:formal}, \rpath{alg:dh:tab:notation}); the
manuscript's seven research questions (\rpath{alg:dh:q:quadratic},
\rpath{alg:dh:q:remainder}, \rpath{alg:dh:q:uniform},
\rpath{alg:dh:q:coefficient}, \rpath{alg:dh:q:numberfield},
\rpath{alg:dh:q:exceptions}, \rpath{alg:dh:q:library}); its
Conclusion and appendices (\rpath{alg:dh:sec:conclusion},
\rpath{alg:dh:sec:ledger}, \rpath{alg:dh:sec:repro}); and the two
subsections of this appendix (\rpath{alg:dh:app:provenance},
\rpath{alg:dh:app:rerun}). In Part~I five labels were added to
existing subsections of Section~\ref{alg:sec:questions}
(\rpath{alg:sub:q-dyadic}, \rpath{alg:sub:q-split},
\rpath{alg:sub:q-height}, \rpath{alg:sub:q-density},
\rpath{alg:sub:q-library}); no existing label was renamed or removed,
and no Part~I number moved.
\item \emph{Numbering.} Manuscript Section~$k$ ($k=1,\ldots,9$) is printed
as Section~$k+14$, so its Theorem~$k.m$ and equation~($k.m$) are
Theorem~$(k+14).m$ and equation~($(k+14).m$). Its seven research
questions, numbered 1--7 by a counter of their own, are Research
questions~\ref{alg:dh:q:quadratic}--\ref{alg:dh:q:library} under the
report's shared counter; its Conclusion is
Section~\ref{alg:dh:sec:conclusion} and its Appendices~A and~B are
Sections~\ref{alg:dh:sec:ledger} and~\ref{alg:dh:sec:repro}. The titles
of Sections~\ref{alg:dh:sec:questions}, \ref{alg:dh:sec:conclusion}
and~\ref{alg:dh:sec:repro} gained ``of Part~II''.
\item \emph{Scope of the counting claim.} The sentence after
Theorem~\ref{alg:dh:thm:headline-count}, ``This settles the completely
split sector of the prior report's Research Question~12.3'', and the
matching sentence of the abstract (Section~\ref{alg:dh:sec:source}) are
printed with the restriction to square boxes $|A|,|B|\le T$ added and
marked \wtag. The theorem is unchanged; the sentence overstated it,
because Research question~\ref{alg:sub:q-height} asks for boxes
$|A|\le H_A$, $|B|\le H_B$. The delivered
\rpath{02-dense-hasse-STATUS.md} keeps the unrestricted wording.
\item \emph{Notation.} The manuscript's sum $S$ is printed $\Sigma$ (in
the proof of Lemma~\ref{alg:dh:lem:dyadic-lift} and before
\eqref{alg:dh:eq:thirdbranch}); nothing else was renamed
(Section~\ref{alg:dh:sec:notation}). Macros: the manuscript's \verb|\pin|
is \verb|\dhsha| here and its \verb|\pathtext| is the report's
\verb|\rpath|; its definitions of \verb|\Z|, \verb|\Q|, \verb|\R|,
\verb|\C|, \verb|\A|, \verb|\Zp|, \verb|\Qp|, \verb|\F|, \verb|\Ztwo|,
\verb|\im|, \verb|\vol|, \verb|\lcm| and \verb|\target| duplicate the
report's and were dropped; its unused \verb|\disc|, \verb|\ord|,
\verb|\rad| and \verb|\sgn| were not carried over; \verb|\Hsplit|,
\verb|\Isplit| and \verb|\calB| were added, and the package
\texttt{needspace} is loaded for its \verb|\Needspace|. The manuscript's
separate \texttt{question} counter was replaced by the report's.
\item \emph{Citations.} The manuscript's \texttt{proveitformal},
\texttt{priorreport} and \texttt{dlmf} are the entries
\cite{dh:proveitformal}, \cite{dh:priorreport} and \cite{dh:dlmf}; its
\texttt{priorreport} is Part~I itself. Its entry for Gao's paper was
merged into Part~I's entry \cite{gao}, with its note kept.
\item \emph{Notes marked \wtag} were added in
Section~\ref{alg:dh:sec:intro} (Part~I and the shipped name of
\texttt{SOURCES.md}), at the two square-box sentences, after the program
list of Section~\ref{alg:dh:sec:certificates}, after three research
questions (cross-references to Part~I), and at the end of
Section~\ref{alg:dh:sec:repro} (shipped names). The manuscript's title
page is printed as the abstract and status note of
Section~\ref{alg:dh:sec:source}.
\item \emph{Part~I} gained dated notes after Research
questions~\ref{alg:sub:q-split}, \ref{alg:sub:q-height}
and~\ref{alg:sub:q-density}, on the title page, in
Appendix~\ref{alg:app:replay} and in Appendix~\ref{alg:app:onesource};
the Part~II heading, its running head and two contents entries are new.
\end{enumerate}
No mathematical statement of the manuscript was changed.

Where the addition had to choose: (a) the manuscript became Part~II of
this report, not a new report, because it studies the same map through the
same chart and cubic and extends Part~I's $C=2$ family; (b) its
restatements of Part~I are kept in place and marked
(Section~\ref{alg:dh:sec:duplicates}), since all of its labels are kept;
(c) of the two colliding letters, the sum $S$ was renamed rather than the
height $T$, which occurs throughout and whose other use (the indeterminate
of $g$) is Part~I's own; (d) the counting claim was restricted to square
boxes in the printed text.

\begin{center}\small
\begin{tabular}{>{\raggedright}p{0.40\textwidth}>{\raggedright\arraybackslash}p{0.52\textwidth}}
\toprule
Delivered file & In this report\\
\midrule
\texttt{article.tex} & printed as Part~II
(Sections~\ref{alg:dh:sec:intro}--\ref{alg:dh:sec:repro})\\
\texttt{README.md}, \texttt{article.pdf} & not shipped (the report README
and \rpath{article.pdf} replace them)\\
\texttt{SOURCES.md}, \texttt{STATUS.md} &
\rpath{02-dense-hasse-SOURCES.md}, \rpath{02-dense-hasse-STATUS.md}\\
\texttt{Makefile} & \rpath{code/02-dense-hasse-Makefile}\\
\texttt{requirements.txt} & \rpath{data/02-dense-hasse-requirements.txt}\\
\texttt{code/split\_fibres.py}, \texttt{verify\_results.py},
\texttt{count\_heights.py} & \rpath{code/02-dense-hasse-split_fibres.py},
\rpath{code/02-dense-hasse-verify_results.py},
\rpath{code/02-dense-hasse-count_heights.py}\\
\texttt{data/build\_review.json}, \texttt{example.json},
\texttt{height\_counts.json}, \texttt{height\_counts.txt},
\texttt{verification.json}, \texttt{verification.txt} &
the same names under \rpath{data/} with the prefix
\texttt{02-dense-hasse-}\\
\bottomrule
\end{tabular}
\end{center}
Every shipped file is byte-identical to the delivery; no checksum ledger
was delivered. The shipped \rpath{02-dense-hasse-SOURCES.md},
\rpath{02-dense-hasse-STATUS.md} and \rpath{code/02-dense-hasse-Makefile}
still name the delivered paths. Two of the recorded outputs are not what
the delivered programs write: \rpath{data/02-dense-hasse-example.json}
is the standard output of \texttt{split\_fibres.py 1 1 3 -{}-modulus
18144000} (identical at the write), whereas \texttt{verify\_results.py}
writes an \texttt{example.json} without its \texttt{"modular\_witness"}
block; and \rpath{data/02-dense-hasse-verification.txt} is a captured
console log of a different run (0.524~s) from the one that wrote
\rpath{data/02-dense-hasse-verification.json} (0.511~s).

\subsection{Rerunning the checks of Part~II}\label{alg:dh:app:rerun}
\wtag{} [Added 30 September 2026, batch~58.] Rerun on a copy with the
delivered names, never in the report: \texttt{count\_heights.py} imports
\texttt{split\_fibres} by its delivered name, and
\texttt{verify\_results.py} and \texttt{count\_heights.py} write
\texttt{verification.json}, \texttt{example.json},
\texttt{height\_counts.json} and \texttt{height\_counts.txt} into
\rpath{../data/} relative to their own directory.
\begin{lstlisting}[language=bash]
W=<scratch>              # outside the repository
mkdir -p "$W/code" "$W/data"
for f in count_heights split_fibres verify_results; do
  cp "code/02-dense-hasse-$f.py" "$W/code/$f.py"; done
cd "$W"
uv run --no-project --with sympy==1.14.0 --with numpy==2.3.5 \
  --with mpmath==1.3.0 python code/verify_results.py
uv run --no-project --with numpy==2.3.5 --with mpmath==1.3.0 \
  python code/count_heights.py
py code/split_fibres.py 1 1 3 --modulus 18144000
\end{lstlisting}
At the write these commands were run so, with Python~3.13.5 under
\texttt{uv}, SymPy~1.14.0, NumPy~2.3.5 and mpmath~1.3.0: all twelve
groups passed (1.5~s); the regenerated \texttt{verification.json} and
\texttt{height\_counts.json} equal the shipped ones apart from their
\texttt{elapsed\_seconds} fields and Windows line endings,
\texttt{height\_counts.txt} is identical apart from line endings, and
the regenerated \texttt{example.json} is the shipped one without its
\texttt{"modular\_witness"} block, which the third command prints.
\texttt{split\_fibres.py} needs only the standard library.
\rpath{data/02-dense-hasse-build_review.json} is a static record of the
delivered 20-page PDF, not regenerated.

\subsection{Part~III: the third manuscript}\label{alg:ns:app:provenance}
\wtag{} [Added 30 September 2026, batch~64.] Part~III is built from
\emph{one} manuscript: manuscript~02 of batch~64 of ProveIt's incoming
reports, the archive \texttt{ProveIt\_Beyond\_Split\_Fibers} (inner
directory \texttt{ProveIt\_Nonsplit\_Hasse}, main file
\texttt{article.tex}, a 24-page US-Letter PDF dated 30~September 2026,
author line and PDF metadata ``AI-assisted research draft prepared for
Vladimir Reshetnikov''), pinned to ProveIt commit \nssha. The archive
arrived in commit \texttt{7747fcfdd} and was placed in
\texttt{3025c15df} as an addition to this report, with the file prefix
\texttt{03-nonsplit-hasse-}. The report is therefore built from three
manuscripts: Part~I (batch~36, manuscript~03;
Appendix~\ref{alg:app:onesource}), Part~II (batch~58, manuscript~01;
Appendix~\ref{alg:dh:app:provenance}) and Part~III. Nothing was merged
across the Parts and nothing was selected away: every result, proof,
example, table, question and limitation of the manuscript is printed in
Sections~\ref{alg:ns:sec:intro}--\ref{alg:ns:app:dyadic}, and it
contributed everything there except the text marked \wtag.
Section~\ref{alg:ns:sec:relation} is new.

Parts~I and~II are unchanged from Part~III's pin to the placement commit
(the only change to this directory in between is the placed files), so
the manuscript's references to Part~I's Research questions 12.2 and 12.3
and to Part~II's open questions and split constant are current. The
manuscript also read this directory's README and the collection's
\rpath{SetTheory/Cardinals/docs/reports/README.md} and
\rpath{SetTheory/Cardinals/README.md} at its pin, for scope and placement
only (\rpath{03-nonsplit-hasse-SOURCES.md}).

The editorial changes are these.
\begin{enumerate}[label=(\arabic*)]
\item \emph{Labels.} The manuscript's 78 labels carry the prefix
\texttt{alg:ns:}, unchanged after it. Twenty were added for Part~III: the
relation section, its six subsections and the notation table
(\rpath{alg:ns:sec:relation}, \rpath{alg:ns:sec:source},
\rpath{alg:ns:sec:answers}, \rpath{alg:ns:sec:consistency},
\rpath{alg:ns:sec:notation}, \rpath{alg:ns:sec:duplicates},
\rpath{alg:ns:sec:formalproject}, \rpath{alg:ns:tab:notation}); the
manuscript's nine research questions (\rpath{alg:ns:q:second},
\rpath{alg:ns:q:power}, \rpath{alg:ns:q:slices}, \rpath{alg:ns:q:unequal},
\rpath{alg:ns:q:quadfields}, \rpath{alg:ns:q:dyadicprec},
\rpath{alg:ns:q:beyond}, \rpath{alg:ns:q:numberfield},
\rpath{alg:ns:q:software}); its Conclusion
(\rpath{alg:ns:sec:conclusion}); and the two subsections of this
appendix (\rpath{alg:ns:app:provenance}, \rpath{alg:ns:app:rerun}). In
Part~I the label \rpath{alg:sub:q-extensions} was added to the existing
subsection 12.5. No existing label was renamed or removed, and no number
of Part~I or Part~II moved.
\item \emph{Numbering.} Manuscript Section~$k$ ($k=1,\ldots,14$) is
printed as Section~$k+28$, so its Theorem~$k.m$ and equation~($k.m$) are
Theorem~$(k+28).m$ and equation~($(k+28).m$); its Appendices~A--D are
Sections~\ref{alg:ns:app:package}--\ref{alg:ns:app:dyadic}, with the
manuscript's reference to ``Appendix~D'' printed as a section reference.
Its research questions already shared the theorem counter. The titles of
its research-question, Conclusion and appendix sections gained ``of
Part~III''.
\item \emph{Citations.} The manuscript's entry \texttt{prior} reads:
``\emph{Arithmetic Local--Global Dichotomies for ProveIt's Three-Variable
Keller Map}, with Part~II, \emph{Zariski-dense integral Hasse failures:
complete split-fiber criteria and sharp target-height asymptotics},
AI-assisted research reports dated September 24 and 29, 2026, merged in
ProveIt, snapshot [the pin]. Unrefereed and not formalized. See Part~I,
Sections~12.2--12.3, and the Part~II discussion of remaining nonsplit and
coefficient-certificate questions.'' It is this report, so its four
citations became cross-references to Parts~I and~II, marked \wtag, and
Part~I's section numbers 12.2 and 12.3 became cross-references (the
printed numbers are unchanged). Its entry \texttt{formal} names the same
Lean file at its own pin, identical to Part~II's entry, and is cited as
\cite{dh:proveitformal}. Its entries for Gao and Milne were merged into
Part~I's entries \cite{gao} and \cite{milne}, with their notes kept.
\item \emph{Notation.} The split count $S(T)$ is printed $\Sigma(T)$,
and the implied constant $C$ in the proof of
Lemma~\ref{alg:ns:lem:range} is printed $c_0$
(Section~\ref{alg:ns:sec:notation}); nothing else was renamed. Macros: the
manuscript's \verb|\pin| is \verb|\nssha| here; \verb|\NS| and
\verb|\spc| were added; its definitions of \verb|\Z|, \verb|\Q|,
\verb|\R|, \verb|\C|, \verb|\F|, \verb|\A|, \verb|\Zp|, \verb|\Qp|,
\verb|\Oh|, \verb|\Gal|, \verb|\Disc|, \verb|\im|, \verb|\lcm|,
\verb|\target| and \verb|\rpath| duplicate the report's and were dropped;
its unused \verb|\odd| was not carried over.
\item \emph{Notes marked \wtag} were added in
Section~\ref{alg:ns:sec:intro} (the report and the pin, the formal
declarations, Part~II's split constant), at the four former citations,
after the Remark following Proposition~\ref{alg:ns:prop:two} (Part~I's
residue counts), after Lemma~\ref{alg:ns:lem:split-nointeger} (Part~I's
Theorem~\ref{alg:thm:plane}, which the manuscript does not cite), after
Lemma~\ref{alg:ns:lem:mean} (Part~I's density), at the constant $c_0$,
after six research questions (cross-references to Parts~I and~II), at
the heads of Sections~\ref{alg:ns:app:package}--\ref{alg:ns:app:dyadic},
and at the end of Section~\ref{alg:ns:app:package} (shipped names). The
manuscript's title page is printed as the abstract and evidence note of
Section~\ref{alg:ns:sec:source}; its author line is recorded only here.
\item \emph{Parts~I and~II} gained dated notes on the title page, after
Research questions~\ref{alg:sub:q-dyadic}, \ref{alg:sub:q-split},
\ref{alg:sub:q-height} and~\ref{alg:sub:q-extensions}, in the
introductory paragraph of Part~II, in
Section~\ref{alg:dh:sec:answers}, after
Theorem~\ref{alg:dh:thm:headline-count}, after Part~II's Research
questions~\ref{alg:dh:q:quadratic}, \ref{alg:dh:q:coefficient}
and~\ref{alg:dh:q:numberfield}, in Section~\ref{alg:dh:sec:conclusion},
in Appendices~\ref{alg:app:replay} and~\ref{alg:app:onesource}, and in
the bibliography entries \cite{gao}, \cite{milne}
and~\cite{dh:proveitformal}. The title page's vertical spaces were
reduced so that its note fits on one page. The Part~III heading, its
running head and its contents entry are new.
\end{enumerate}
No mathematical statement of the manuscript was changed.

Where the addition had to choose: (a) the manuscript became Part~III of
this report, not a new report, because it studies the same map through
the same chart and cubic and answers questions of Parts~I and~II;
(b) its restatements of Parts~I and~II are kept in place and marked
(Section~\ref{alg:ns:sec:duplicates}), since all of its labels are kept;
in particular its proof that split fibers on $C=2$ have no integral point
is kept beside a pointer to Part~I's Theorem~\ref{alg:thm:plane};
(c) of the colliding letters, the count $S(T)$ was renamed rather than
the indeterminate $S$ or the set of places, following Part~II's rename of
its sum, and the constant $C$ was renamed rather than the coordinate;
(d) Part~II's split-sector theorems are kept as proved, and the sentences
that restrict the answer to Research question~\ref{alg:sub:q-height} to
the split sector are re-scoped by dated notes, not edited; (e) the
bibliography entry for this report itself became cross-references.

\begin{center}\small
\begin{tabular}{>{\raggedright}p{0.40\textwidth}>{\raggedright\arraybackslash}p{0.52\textwidth}}
\toprule
Delivered file & In this report\\
\midrule
\texttt{article.tex} & printed as Part~III
(Sections~\ref{alg:ns:sec:intro}--\ref{alg:ns:app:dyadic})\\
\texttt{README.md}, \texttt{article.pdf}, \texttt{SHA256SUMS} & not
shipped (the report README and \rpath{article.pdf} replace the first
two)\\
\texttt{SOURCES.md}, \texttt{STATUS.md} &
\rpath{03-nonsplit-hasse-SOURCES.md}, \rpath{03-nonsplit-hasse-STATUS.md}\\
\texttt{Makefile} & \rpath{code/03-nonsplit-hasse-Makefile}\\
\texttt{requirements.txt} & \rpath{data/03-nonsplit-hasse-requirements.txt}\\
\texttt{code/arithmetic.py}, \texttt{classify.py}, \texttt{count.py},
\texttt{symbolic.py}, \texttt{verify.py} & the same names under
\rpath{code/} with the prefix \texttt{03-nonsplit-hasse-}\\
\texttt{data/build\_review.json}, \texttt{classification\_example.json},
\texttt{counts.json}, \texttt{symbolic.json}, \texttt{verification.json} &
the same names under \rpath{data/} with the prefix
\texttt{03-nonsplit-hasse-}\\
\bottomrule
\end{tabular}
\end{center}
Every shipped file is byte-identical to the delivery, and matches the
delivered \texttt{SHA256SUMS}. The shipped
\rpath{03-nonsplit-hasse-SOURCES.md}, \rpath{03-nonsplit-hasse-STATUS.md}
and \rpath{code/03-nonsplit-hasse-Makefile}, and the manuscript's own text
(Sections~\ref{alg:ns:sec:computation} and~\ref{alg:ns:app:package}),
still name the delivered paths. \rpath{data/03-nonsplit-hasse-build_review.json}
records the delivered 24-page PDF, which is not shipped, and
\rpath{data/03-nonsplit-hasse-classification_example.json} is the
standard output of \texttt{classify.py 0 4 2}; no program writes that
file.

\subsection{Rerunning the checks of Part~III}\label{alg:ns:app:rerun}
\wtag{} [Added 30 September 2026, batch~64.] Rerun on a copy with the
delivered names, never in the report: \texttt{classify.py},
\texttt{count.py}, \texttt{symbolic.py} and \texttt{verify.py} import
\texttt{arithmetic} by its delivered name, and \texttt{verify.py},
\texttt{symbolic.py} and \texttt{count.py} (without \texttt{-{}-output})
write \texttt{verification.json}, \texttt{symbolic.json} and
\texttt{counts.json} into \rpath{../data/} relative to their own
directory.
\begin{lstlisting}[language=bash]
W=<scratch>              # outside the repository
mkdir -p "$W/code" "$W/data"
for f in arithmetic classify count symbolic verify; do
  cp "code/03-nonsplit-hasse-$f.py" "$W/code/$f.py"; done
cd "$W"
py code/verify.py
uv run --no-project --with sympy==1.14.0 python code/symbolic.py
py code/count.py
uv run --no-project --with sympy==1.14.0 python code/classify.py 0 4 2
\end{lstlisting}
At the write these commands were run so, with Python~3.14.4 for the
standard-library programs and Python~3.13.5 under \texttt{uv} with
SymPy~1.14.0: all exit with status zero, \texttt{verify.py} and
\texttt{symbolic.py} report \texttt{"status": "PASS"}, and the
regenerated \texttt{verification.json}, \texttt{symbolic.json} and
\texttt{counts.json} and the printed classification equal the shipped
files apart from Windows line endings. \texttt{verify.py},
\texttt{count.py} and \texttt{arithmetic.py} need only the standard
library.

\clearpage
\begin{thebibliography}{9}
\bibitem{proveit}
Vladimir Reshetnikov, \emph{ProveIt: Jacobian conjecture project},
\nolinkurl{Algebra/JacobianConjecture/README.md} and associated Lean/Rocq
sources. Repository snapshot \psha, inspected September 24, 2026.
\href{https://github.com/VladimirReshetnikov/ProveIt/blob/e21766d04c2b8a9b2cdba0cd43e563024b3bd1b9/Algebra/JacobianConjecture/README.md}{Pinned project documentation}.

\bibitem{proveitresearch}
Vladimir Reshetnikov, \emph{Searching for a simpler counterexample},
\nolinkurl{Algebra/JacobianConjecture/Research/README.md}, in ProveIt,
same pinned snapshot.
\href{https://github.com/VladimirReshetnikov/ProveIt/blob/e21766d04c2b8a9b2cdba0cd43e563024b3bd1b9/Algebra/JacobianConjecture/Research/README.md}{Pinned research documentation}.

\bibitem{gao}
Shuhong Gao, \emph{Counterexamples to the Jacobian conjecture in dimensions
greater than two}, arXiv:2608.00222v1 (2026), especially
Theorems 3.3--3.4 and Appendix A.1.
\url{https://arxiv.org/html/2608.00222v1}.
\wtag{} Part~II cites the abstract page
\url{https://arxiv.org/abs/2608.00222} (arXiv:2608.00222, 2026); its
entry reads: ``Used for attribution of the pre-existing geometric
setting; none of its generalizations is an assumption of the proofs
here.''
\wtag{} Part~III cites arXiv:2608.00222v1, July 31, 2026, with the note:
``Theorems~3.3--3.4 provide the relevant degree-three map and fiber
geometry; the map is attributed to Alp\"oge's announcement.''

\bibitem{finitefield}
\texttt{royvanrijn/jacobian-research}, \emph{Exact finite-field value
distribution}, archived research note,
\nolinkurl{research/archive/legacy-notes/FINITE_FIELD_VALUE_DISTRIBUTION.md}.
Repository snapshot \rsha, inspected September 24, 2026.
\href{https://github.com/royvanrijn/jacobian-research/blob/caf5685db1099cb90e608556d4b8d825734f29cb/research/archive/legacy-notes/FINITE_FIELD_VALUE_DISTRIBUTION.md}{Pinned finite-field note}.
The repository note is cited as prior public work, not as an independently
peer-reviewed source.

\bibitem{degreefive}
\texttt{royvanrijn/jacobian-research}, \emph{An integral Jacobian-one Hasse
failure}, archived research note,
\nolinkurl{research/archive/non-elliptic/verified/INTEGRAL_HASSE_JACOBIAN_ONE.md},
same pinned snapshot.
\href{https://github.com/royvanrijn/jacobian-research/blob/caf5685db1099cb90e608556d4b8d825734f29cb/research/archive/non-elliptic/verified/INTEGRAL_HASSE_JACOBIAN_ONE.md}{Pinned degree-five note}.
Its distinct map and rational obstruction are used only for comparison,
not as assumptions in any theorem of this article.

\bibitem{milne}
James S. Milne, \emph{Algebraic Number Theory}, version 3.08,
July 19, 2020, Theorem 8.31 and Example 8.36.
\url{https://www.jmilne.org/math/CourseNotes/ANT.pdf}.
\wtag{} Part~III cites the same version, adding ``printed pages
147--148''.

\bibitem{dh:proveitformal}
V. Reshetnikov, \emph{ProveIt: Jacobian-conjecture formal development},
repository at commit \dhsha, inspected September 29, 2026.
\url{https://github.com/VladimirReshetnikov/ProveIt/blob/1085b506d65e207a05b7e9c861bb1fe88432fe38/Algebra/JacobianConjecture/Lean/JacobianConjecture/Counterexample.lean}.
\wtag{} Cited by Part~II, and by Part~III at its pin \nssha, where the
file is identical.

\bibitem{dh:priorreport}
\emph{Arithmetic Local--Global Dichotomies for ProveIt's Three-Variable
Keller Map: Split integral Hasse failures, a dense parameter family,
and exact $p$-adic fiber laws}, research draft prepared for
Vladimir Reshetnikov, September 24, 2026; ProveIt research-report
collection, with later repository additions.
Sections 2, 4, and 12, especially Research Questions 12.2--12.4.
\url{https://github.com/VladimirReshetnikov/ProveIt/blob/1085b506d65e207a05b7e9c861bb1fe88432fe38/SetTheory/Cardinals/docs/reports/jacobian-conjecture/arithmetic-local-global-fibers/article.tex}.
\wtag{} Cited by Part~II; this is Part~I of the present report, at
Part~II's pin.

\bibitem{dh:dlmf}
NIST Digital Library of Mathematical Functions, Section 5.12,
\emph{Beta Function}, especially the beta--gamma identity and Euler's
integral. \url{https://dlmf.nist.gov/5.12}. \wtag{} Cited by Part~II.
\end{thebibliography}
\end{document}
